% Production et traduction de textes liés à la vie économique — Allemand Tle A — Cours signé OnBuch+
% Programme officiel MINESEC. Compiler avec : tectonic AL10-production-et-traduction-de-textes-lies-a-la-vie-economique.tex
\newcommand{\DOCMATIERE}{Allemand}
\newcommand{\DOCNIVEAU}{Tle A}
\newcommand{\DOCMODULE}{Chapitre 2 : Vie économique}
\newcommand{\DOCLECON}{AL10}
\newcommand{\DOCTITRE}{Production et traduction de textes liés à la vie économique}
% OnBuch+ — préambule « cours signés » ENRICHI (compilé avec tectonic / XeLaTeX)
% Système visuel « Pop » d'OnBuch+ : fond crème (jamais blanc), bordures encre
% épaisses, ombres « dures » décalées, titres Archivo Black en capitales.
% Version MATHÉMATIQUES : graphes (pgfplots/tikz), boîtes pédagogiques
% (définition, propriété, méthode, attention, le savais-tu, exercice résolu,
% exercice, TP, bilan), page de garde et signature OnBuch+ ancrée en bas.
%
% Macros attendues dans main.tex AVANT \input{preamble} :
%   \DOCMATIERE  \DOCNIVEAU  \DOCMODULE  \DOCLECON (numéro)  \DOCTITRE
\documentclass[11pt]{article}

\usepackage{fontspec}
\usepackage[ngerman,french]{babel} % français = langue principale ; l'allemand via \all{..} / textebox
\usepackage[a4paper,margin=2cm,top=3.4cm,bottom=2.8cm,headheight=1.4cm,footskip=1.3cm]{geometry}
\usepackage{amsmath,amssymb,amsfonts}
\usepackage{mathtools} % \xmapsto, \coloneqq... souvent utilisés par les modèles
\usepackage{cancel} % simplifications barrées (bilans de forces)
% \abs{z} (module, valeur absolue) et \norm{u} (norme), utilisés par les modèles
\providecommand{\abs}{}\let\abs\relax\DeclarePairedDelimiter{\abs}{\lvert}{\rvert}
\providecommand{\norm}{}\let\norm\relax\DeclarePairedDelimiter{\norm}{\lVert}{\rVert}
\usepackage[table]{xcolor} % [table] : fournit aussi \rowcolors (alternance de lignes)
\usepackage{pagecolor}
\usepackage{tikz}
\usetikzlibrary{arrows.meta,positioning,calc,shapes.geometric,shapes.misc,shapes.arrows,decorations.pathmorphing,decorations.pathreplacing,decorations.markings,patterns,shadows,mindmap,trees,fit,backgrounds,matrix,intersections,angles,quotes}
\usepackage{pgfplots}
\usepackage[version=4]{mhchem} % équations nucléaires \ce{^{235}_{92}U}
\usepackage[european,siunitx]{circuitikz} % schémas électriques
\pgfplotsset{compat=1.18}
\usepgfplotslibrary{fillbetween,groupplots} % aires entre courbes (calcul intégral)
\usepackage{enumitem}
\usepackage{fancyhdr}
% [most] charge toutes les bibliothèques tcolorbox (listings, minted, raster,
% documentation...) et gonfle inutilement la mémoire TeX sur un document long
% (~300 boîtes) : on ne charge que ce qui est réellement utilisé.
\usepackage[skins,breakable]{tcolorbox}
\usepackage{graphicx}
\usepackage{colortbl}
\usepackage{booktabs}
\usepackage{tabularx}
\usepackage{diagbox} % case d'en-tête diagonale des tableaux à double entrée
% Colonnes extensibles centrée / gauche / droite (C, L, R), souvent utilisées par les modèles
\newcolumntype{C}{>{\centering\arraybackslash}X}
\newcolumntype{L}{>{\raggedright\arraybackslash}X}
\newcolumntype{R}{>{\raggedleft\arraybackslash}X}
\usepackage{array}
\usepackage{multicol}
\usepackage{multirow}
\usepackage{siunitx}
% parse-numbers=false : accepte aussi \SI{2,5 \times 10^{-3}}{...}
\sisetup{locale=FR,output-decimal-marker={,},group-minimum-digits=4,parse-numbers=false}
\DeclareSIUnit\ohm{\ensuremath{\symup{\Omega}}} % U+2126 absent de la police
\DeclareSIUnit\division{div} % réglages d'oscilloscope (ms/div, V/div)
\DeclareSIUnit\tr{tr} % tours (vitesse de rotation, ex. \SI{1500}{\tr\per\minute})
\DeclareSIUnit\molar{M} % concentration molaire (ex. \SI{3}{\molar}, \SI{1}{\micro\molar})
\DeclareSIUnit\an{an} % année (le modèle confond parfois avec \year, un compteur TeX) — ex. \SI{5}{\kilo\meter\per\an}
\DeclareSIUnit\FCFA{FCFA} % franc CFA (ex. \SI{500}{\FCFA\per\kilo\gram})
\DeclareSIUnit\degreeBrix{\textdegree Brix} % teneur en sucre d'un sirop/jus (réfractométrie)
\DeclareSIUnit\month{mois} % le modèle utilise parfois \month, un compteur TeX, comme unité de durée
\usepackage{qrcode}
% \degree : commande fréquemment utilisée par les modèles pour un angle ou une
% température en dehors de \SI{}{} (non fournie par les paquets chargés).
\providecommand{\degree}{^{\circ}}
% \qed (fin de démonstration), utilisé par les modèles sans amsthm
\providecommand{\qed}{\hfill$\square$}
% \barre : double barre des valeurs interdites dans un tableau de variations (tkz-tab)
\providecommand{\barre}{\ensuremath{\|}}
% environnement proof (amsthm non chargé) utilisé par les modèles
\ifdefined\proof\else\newenvironment{proof}[1][Démonstration]{\par\noindent\textit{#1.}\ }{\qed\par}\fi
\usepackage{lastpage}
\usepackage{hyperref}
\hypersetup{hidelinks}

% ============================================================
% Palette « Pop » OnBuch+ — mêmes hex que la classe P (plus_ui.dart)
% ============================================================
\definecolor{poporange}{HTML}{F59321}
\definecolor{popdark}{HTML}{E07A0C}
\definecolor{popcream}{HTML}{FFF6E8}
\definecolor{popink}{HTML}{1C1714}
\definecolor{poppurple}{HTML}{7A5AE0}
\definecolor{popgreen}{HTML}{2E9E6A}
\definecolor{poppink}{HTML}{E0457B}
\definecolor{popblue}{HTML}{2F7DE1}
\definecolor{popgold}{HTML}{C98A12}
\definecolor{popmuted}{HTML}{8A7F76}
\definecolor{popline}{HTML}{EADFD2}
\definecolor{popgray}{HTML}{8A7F76}
\colorlet{popgrey}{popgray}
% Alias tolérants (le modèle nomme parfois les couleurs différemment)
\colorlet{popwhite}{white}
\colorlet{popblack}{popink}
\colorlet{popred}{poppink}
\colorlet{popyellow}{popgold}
% Teintes claires (fonds de tableaux/encadrés) — le modèle nomme parfois
% les couleurs avec un suffixe L (ex. popblueL) sans les avoir définies.
\colorlet{poporangeL}{poporange!20}
\colorlet{popdarkL}{popdark!20}
\colorlet{poppurpleL}{poppurple!20}
\colorlet{popgreenL}{popgreen!20}
\colorlet{poppinkL}{poppink!20}
\colorlet{popblueL}{popblue!20}
\colorlet{popgoldL}{popgold!20}
\colorlet{popredL}{popred!20}
\colorlet{popgrayL}{popgray!20}
% Teintes claires pour fonds de tableaux / zones
\colorlet{popblueL}{popblue!10!white}
\colorlet{poporangeL}{poporange!14!white}
\colorlet{popgreenL}{popgreen!12!white}
\colorlet{poppurpleL}{poppurple!12!white}
\colorlet{poppinkL}{poppink!12!white}
\colorlet{popgoldL}{popgold!14!white}

\pagecolor{popcream}

% ============================================================
% Typographie
% ============================================================
\defaultfontfeatures{Ligatures=TeX}
\setmainfont{PlusJakartaSans-Medium.ttf}[Path=../../../fonts/,BoldFont=PlusJakartaSans-Bold.ttf]
% Symboles, API et emojis absents de la police principale : repli sur DejaVu Sans (caractères actifs)
\newfontface\symfont{DejaVuSans.ttf}[Path=../../../fonts/]
\makeatletter
\newcommand\symactive[1]{\catcode#1=13 \begingroup\lccode`\~=#1 \lowercase{\endgroup\def~}{{\symfont\char#1}}}
\newcommand\symblank[1]{\catcode#1=13 \begingroup\lccode`\~=#1 \lowercase{\endgroup\def~}{}}
\newcommand\symhyph[1]{\catcode#1=13 \begingroup\lccode`\~=#1 \lowercase{\endgroup\def~}{\mbox{-}}}
\newcount\symc
\newcommand\symrange[2]{\symc=#1 \@whilenum\symc<#2 \do{\expandafter\symactive\expandafter{\the\symc}\advance\symc by 1 }}
\newcommand\symblankrange[2]{\symc=#1 \@whilenum\symc<#2 \do{\expandafter\symblank\expandafter{\the\symc}\advance\symc by 1 }}
\makeatother
\symrange{"0250}{"02FF}\symactive{"2713}\symactive{"2714}\symactive{"2717}\symactive{"2718}\symactive{"2610}\symactive{"2611}\symactive{"2612}
\symhyph{"2011}\symhyph{"2E1A}\symblankrange{"1F300}{"1FAFF}\symblank{"FE0F}
% Maths en sans-serif (Fira Math) avec chiffres et lettres droites de Plus
% Jakarta, pour que les formules se fondent dans le texte.
\usepackage[math-style=ISO,bold-style=ISO]{unicode-math}
\setmathfont{FiraMath-Regular.otf}[Path=../../../fonts/]
\setmathfont{PlusJakartaSans-Medium.ttf}[Path=../../../fonts/,range={up/num,up/{Latin,latin},\mathup}]
\usepackage{icomma} % virgule décimale sans espace en mode math
% Caractères mathématiques Unicode absents des polices de texte (Plus Jakarta,
% Archivo Black) : rendus par leur équivalent mathématique, en texte comme en
% formule — sinon ils s'affichent en carré vide (ex. « Arithmétique dans ℤ »).
\usepackage{newunicodechar}
\newunicodechar{ℝ}{\ensuremath{\mathbb{R}}}
\newunicodechar{ℤ}{\ensuremath{\mathbb{Z}}}
\newunicodechar{ℂ}{\ensuremath{\mathbb{C}}}
\newunicodechar{ℕ}{\ensuremath{\mathbb{N}}}
\newunicodechar{ℚ}{\ensuremath{\mathbb{Q}}}
\newunicodechar{𝔻}{\ensuremath{\mathbb{D}}}
\newunicodechar{α}{\ensuremath{\alpha}}
\newunicodechar{β}{\ensuremath{\beta}}
\newunicodechar{γ}{\ensuremath{\gamma}}
\newunicodechar{δ}{\ensuremath{\delta}}
\newunicodechar{ε}{\ensuremath{\varepsilon}}
\newunicodechar{θ}{\ensuremath{\theta}}
\newunicodechar{λ}{\ensuremath{\lambda}}
\newunicodechar{μ}{\ensuremath{\mu}}
\newunicodechar{σ}{\ensuremath{\sigma}}
\newunicodechar{τ}{\ensuremath{\tau}}
\newunicodechar{φ}{\ensuremath{\varphi}}
\newunicodechar{ψ}{\ensuremath{\psi}}
\newunicodechar{ω}{\ensuremath{\omega}}
\newunicodechar{Σ}{\ensuremath{\Sigma}}
% majuscules produites par \MakeUppercase dans les titres (α-aminés -> Α)
\newunicodechar{Α}{\ensuremath{\alpha}}
\newunicodechar{Β}{\ensuremath{\beta}}
\newunicodechar{Γ}{\ensuremath{\Gamma}}
\newunicodechar{Δ}{\ensuremath{\Delta}}
\newunicodechar{Ε}{\ensuremath{\varepsilon}}
\newunicodechar{Θ}{\ensuremath{\Theta}}
\newunicodechar{Λ}{\ensuremath{\Lambda}}
\newunicodechar{Μ}{\ensuremath{\mu}}
\newunicodechar{Τ}{\ensuremath{\tau}}
\newunicodechar{Φ}{\ensuremath{\Phi}}
\newunicodechar{Ψ}{\ensuremath{\Psi}}
\newunicodechar{Ω}{\ensuremath{\Omega}}
\newunicodechar{↦}{\ensuremath{\mapsto}}
\newunicodechar{∈}{\ensuremath{\in}}
\newunicodechar{∉}{\ensuremath{\notin}}
\newunicodechar{∧}{\ensuremath{\wedge}}
\newunicodechar{⇔}{\ensuremath{\Leftrightarrow}}
\newunicodechar{⟺}{\ensuremath{\Longleftrightarrow}}
\newunicodechar{⇒}{\ensuremath{\Rightarrow}}
\newunicodechar{⟹}{\ensuremath{\Longrightarrow}}
\newunicodechar{∘}{\ensuremath{\circ}}
\newunicodechar{∪}{\ensuremath{\cup}}
\newunicodechar{∩}{\ensuremath{\cap}}
\newunicodechar{≡}{\ensuremath{\equiv}}
\newunicodechar{⊂}{\ensuremath{\subset}}
\newunicodechar{⊥}{\ensuremath{\perp}}
\newunicodechar{∃}{\ensuremath{\exists}}
\newunicodechar{∀}{\ensuremath{\forall}}
\newunicodechar{∛}{\ensuremath{\sqrt[3]{\,}}}
\newunicodechar{∜}{\ensuremath{\sqrt[4]{\,}}}
\newunicodechar{⃗}{\ensuremath{{}^{\to}}}
\newunicodechar{ⁿ}{\ifmmode^{n}\else\textsuperscript{\ensuremath{n}}\fi}
\newunicodechar{ˣ}{\ifmmode^{x}\else\textsuperscript{\ensuremath{x}}\fi}
\newunicodechar{ᵇ}{\ifmmode^{b}\else\textsuperscript{\ensuremath{b}}\fi}
\newunicodechar{ʳ}{\ifmmode^{r}\else\textsuperscript{\ensuremath{r}}\fi}
\newunicodechar{ᵘ}{\ifmmode^{u}\else\textsuperscript{\ensuremath{u}}\fi}
\newunicodechar{ʸ}{\ifmmode^{y}\else\textsuperscript{\ensuremath{y}}\fi}
\newunicodechar{ᵃ}{\ifmmode^{a}\else\textsuperscript{\ensuremath{a}}\fi}
\newunicodechar{ᵉ}{\ifmmode^{e}\else\textsuperscript{\ensuremath{e}}\fi}
\newunicodechar{ᵏ}{\ifmmode^{k}\else\textsuperscript{\ensuremath{k}}\fi}
\newunicodechar{ᵖ}{\ifmmode^{p}\else\textsuperscript{\ensuremath{p}}\fi}
\newunicodechar{ᴺ}{\ifmmode^{N}\else\textsuperscript{\ensuremath{N}}\fi}
\newunicodechar{ᶜ}{\ifmmode^{c}\else\textsuperscript{\ensuremath{c}}\fi}
\newunicodechar{ʰ}{\ifmmode^{h}\else\textsuperscript{\ensuremath{h}}\fi}
\newunicodechar{ᵠ}{\ifmmode^{q}\else\textsuperscript{\ensuremath{q}}\fi}
\newunicodechar{ₙ}{\ifmmode_{n}\else\textsubscript{\ensuremath{n}}\fi}
\newunicodechar{ₐ}{\ifmmode_{a}\else\textsubscript{\ensuremath{a}}\fi}
\newunicodechar{ᵢ}{\ifmmode_{i}\else\textsubscript{\ensuremath{i}}\fi}
\newunicodechar{ᵣ}{\ifmmode_{r}\else\textsubscript{\ensuremath{r}}\fi}
\newunicodechar{ₖ}{\ifmmode_{k}\else\textsubscript{\ensuremath{k}}\fi}
\newunicodechar{ᵦ}{\ifmmode_{\beta}\else\textsubscript{\ensuremath{\beta}}\fi}
\newunicodechar{ₓ}{\ifmmode_{x}\else\textsubscript{\ensuremath{x}}\fi}
\newfontfamily\popdisplay{ArchivoBlack-Regular.ttf}[Path=../../../fonts/]
\newfontfamily\poplogo{SpaceGrotesk-Bold.ttf}[Path=../../../fonts/]
\newfontfamily\popsemi{PlusJakartaSans-SemiBold.ttf}[Path=../../../fonts/]
\newfontfamily\popxbold{PlusJakartaSans-ExtraBold.ttf}[Path=../../../fonts/]
\newfontfamily\popxboldls{PlusJakartaSans-ExtraBold.ttf}[Path=../../../fonts/,LetterSpace=3.0]

\setlist[enumerate]{leftmargin=1.7em,itemsep=5pt,topsep=4pt}
\setlist[itemize]{leftmargin=1.7em,itemsep=4pt,topsep=4pt,label={\color{poporange}\textbullet}}
\setlength{\parindent}{0pt}
\setlength{\parskip}{5pt}
\renewcommand{\arraystretch}{1.25}

% pgfplots : style Pop par défaut
\pgfplotsset{
  every axis/.append style={axis line style={popink,line width=1pt},tick style={popink},
    grid style={popline,line width=0.5pt},label style={font=\small},tick label style={font=\footnotesize}},
  popcurve/.style={poporange,line width=1.8pt,no markers,smooth},
}
% Même style disponible dans un \draw TikZ simple (hors pgfplots)
\tikzset{popcurve/.style={poporange,line width=1.8pt}}
\tikzset{popgrid/.style={popline,very thin}}
\colorlet{popcurve}{poporange} % le nom du style est parfois utilisé comme couleur

% ============================================================
% Pill / squiggle
% ============================================================
\newcommand{\poppill}[3]{%
  \tikz[baseline=-0.55ex]\node[draw=popink,line width=1.2pt,rounded corners=2.6pt,%
    fill=#2,inner xsep=6.5pt,inner ysep=3pt]{%
    \popxboldls\fontsize{7}{7}\selectfont\color{#3}\MakeUppercase{#1}};%
}
\newcommand{\popsquiggle}[1]{%
  \tikz[baseline=-0.4ex]{
    \draw[#1,line width=2.6pt,line cap=round]
      plot [smooth] coordinates {(0,0.09) (0.34,0.01) (0.68,0.21) (1.02,0.01) (1.36,0.10)};
  }%
}
% Mot-clé surligné (marqueur orange sous le texte)
\newcommand{\cle}[1]{{\popxbold\color{popdark}#1}}

% ============================================================
% En-tête / pied
% ============================================================
\pagestyle{fancy}
\fancyhf{}
\renewcommand{\headrulewidth}{0pt}
\renewcommand{\footrulewidth}{0pt}
\lhead{\raisebox{-0.15ex}{\poplogo\fontsize{13}{13}\selectfont\color{popink}OnBuch\color{poporange}+}}
\rhead{\poppill{\DOCMATIERE\ \textbullet\ \DOCNIVEAU\ \textbullet\ Leçon \DOCLECON}{white}{popink}}
\lfoot{\popxbold\fontsize{7.6}{7.6}\selectfont\color{popmuted}\MakeUppercase{Cours signé OnBuch+}\ \textbullet\ {\popsemi\DOCTITRE}}
\rfoot{\tikz[baseline=-0.6ex]\node[rounded corners=8.5pt,fill=popink,minimum height=17pt,minimum width=17pt,inner xsep=5pt,inner sep=0]{\popxbold\fontsize{8}{8}\selectfont\color{popcream}\thepage\,/\,\pageref*{LastPage}};}

% ============================================================
% Titres
% ============================================================
\newcommand{\chaptitre}[2]{%
  \begin{tcolorbox}[enhanced,colback=poporange,colframe=popink,boxrule=2.6pt,arc=11pt,
      drop shadow=popink,left=15pt,right=15pt,top=12pt,bottom=11pt,before skip=3pt,after skip=18pt]
    \poppill{#2}{popink}{popcream}\par\vspace{9pt}
    {\raggedright\hyphenpenalty=10000\exhyphenpenalty=10000\popdisplay\fontsize{22}{24}\selectfont\color{popcream}\MakeUppercase{#1}\par}
  \end{tcolorbox}
}
% Section numérotée : pastille orange avec le numéro + titre Archivo
\newcounter{popsec}
\newcommand{\coursec}[1]{%
  \stepcounter{popsec}\setcounter{popsub}{0}%
  \par\vspace{16pt}\noindent
  \begin{minipage}{\linewidth}
  \tikz[baseline=-0.7ex]\node[draw=popink,line width=1.6pt,fill=poporange,rounded corners=4pt,minimum size=22pt,inner sep=0]{\popdisplay\fontsize{12}{12}\selectfont\color{popcream}\thepopsec};%
  \hspace{8pt}{\popdisplay\fontsize{15}{17}\selectfont\color{popink}\MakeUppercase{#1}}\par
  \vspace{4pt}\noindent{\color{poporange}\rule{36pt}{3.6pt}}
  \end{minipage}\par\nopagebreak\vspace{8pt}
}
\newcounter{popsub}
\newcommand{\courssub}[1]{%
  \stepcounter{popsub}%
  \par\vspace{9pt}\noindent{\popxbold\fontsize{11.5}{13}\selectfont\color{popink}{\color{poporange}\thepopsec.\thepopsub}\hspace{6pt}#1}\par\nopagebreak\vspace{4pt}
}

% ============================================================
% Boîtes pédagogiques — même recette : fond blanc, bordure épaisse colorée,
% ombre dure encre, badge plein. Titre optionnel [#1].
% ============================================================
\tcbset{popbase/.style={breakable,enhanced,colback=white,boxrule=2.3pt,arc=9pt,
  shadow={3pt}{-3pt}{0pt}{popink},left=14pt,right=14pt,top=11pt,bottom=11pt,before skip=11pt,after skip=15pt,
  fontupper=\popsemi\fontsize{10.6}{14.5}\selectfont\color{popink},
  fontlower=\popsemi\fontsize{10.6}{14.5}\selectfont\color{popink}}}
% Coche dessinée (le glyphe ✓ n'existe pas dans les polices du document)
\newcommand{\popcheck}[1]{\tikz[baseline=-0.5ex]\draw[#1,line width=1.6pt,line cap=round,line join=round](-0.13,0.0)--(-0.04,-0.09)--(0.13,0.1);}
\providecommand{\coche}{\popcheck{popgreen}}
\providecommand{\resultat}[1]{\par\smallskip\noindent\fcolorbox{popgreen}{popgreenL}{\parbox{\dimexpr\linewidth-2\fboxsep-2\fboxrule\relax}{\textbf{Résultat.} #1}}\par\smallskip}
\newcommand{\popboxhead}[3]{\poppill{#1}{#2}{white}\ifx\relax#3\relax\else\hspace{6pt}{\popxbold\fontsize{10.6}{12}\selectfont\color{popink}#3}\fi\par\vspace{7pt}}

\newtcolorbox{aretenir}[1][]{popbase,colframe=popgreen,before upper={\popboxhead{À retenir}{popgreen}{#1}}}
% Texte en allemand (document de lecture, dialogue, extrait) : cadre
% d'étude. Un titre optionnel (source, auteur) s'affiche dans l'en-tête.
\newtcolorbox{textebox}[1][]{popbase,colframe=popblue,before upper={\popboxhead{Text}{popblue}{#1}\begin{otherlanguage}{ngerman}},after upper={\end{otherlanguage}}}
% Marques ✓ / ✗ (forme correcte / fautive) souvent utilisées par les modèles
\providecommand{\cross}{\textcolor{poppink}{\ensuremath{\times}}}
\providecommand{\xmark}{\cross}\providecommand{\crossmark}{\cross}\providecommand{\cmark}{\ensuremath{\checkmark}}
% Ligne de réponse / justification à compléter (inventée par les modèles)
\providecommand{\justification}[1]{\par\noindent\textit{Justification~:} \dotfill\par}
% \textipa (package tipa non chargé) : la police ne couvre pas l'API -> simple passage
\providecommand{\textipa}[1]{#1}
% Soulignement (ulem non chargé) : \uline, \ul
\providecommand{\uline}[1]{\underline{#1}}\providecommand{\ul}[1]{\underline{#1}}
% \glossaire{\item ...} inventée par les modèles : liste de glossaire
\providecommand{\glossaire}[1]{\begin{itemize}#1\end{itemize}}
% \alert (beamer) inventée par les modèles : texte mis en évidence
\providecommand{\alert}[1]{\textbf{\textcolor{poppink}{#1}}}
% variantes ASCII d'ordinaux écrites en macro
\providecommand{\iere}{\textsuperscript{re}}\providecommand{\ieme}{\textsuperscript{e}}\providecommand{\ere}{\textsuperscript{re}}\providecommand{\eme}{\textsuperscript{e}}\providecommand{\er}{\textsuperscript{er}}\providecommand{\ie}{\textsuperscript{e}}
% Ordinaux français écrits en macro par les modèles : 1\ière, 3\ième
\newcommand{\ière}{\textsuperscript{re}}\newcommand{\ième}{\textsuperscript{e}}
% \exemple (inventée par les modèles) : étiquette « Exemple : »
\providecommand{\exemple}{\textbf{Exemple~:}\ }
% \square utilisable aussi hors mode math (cases à cocher)
\AtBeginDocument{\let\squareMath\square\renewcommand{\square}{\ensuremath{\squareMath}}}
% Mot ou phrase en allemand à l'intérieur d'un paragraphe français
\newcommand{\all}[1]{\ifmmode\text{\textit{#1}}\else\begingroup\selectlanguage{ngerman}\itshape #1\endgroup\fi}
\let\esp\all % alias toléré
\newtcolorbox{exemplebox}[1][]{popbase,colframe=poporange,before upper={\popboxhead{Exemple}{poporange}{#1}}}
\newtcolorbox{definition}[1][]{popbase,colframe=popblue,before upper={\popboxhead{Définition}{popblue}{#1}}}
\newtcolorbox{propriete}[1][]{popbase,colframe=poppurple,before upper={\popboxhead{Propriété}{poppurple}{#1}}}
\newtcolorbox{methode}[1][]{popbase,colframe=popgold,before upper={\popboxhead{Méthode}{popgold}{#1}}}
\newtcolorbox{attention}[1][]{popbase,colframe=poppink,before upper={\popboxhead{Attention}{poppink}{#1}}}
\newtcolorbox{experience}[1][]{popbase,colframe=popblue,colback=popblueL,before upper={\popboxhead{Expérience}{popblue}{#1}}}
% « Le savais-tu ? » : carte violette pleine (contexte, culture, Cameroun)
\newtcolorbox{savaistu}[1][]{popbase,colback=poppurple,colframe=popink,
  fontupper=\popsemi\fontsize{10.4}{14.2}\selectfont\color{white},
  before upper={\poppill{Le savais-tu\,?}{white}{poppurple}\ifx\relax#1\relax\else\hspace{6pt}{\popxbold\color{white}#1}\fi\par\vspace{7pt}}}
% Exercice résolu : énoncé en haut, \tcblower puis la solution sur fond crème
\newcounter{popexo}\newcounter{popexr}
\newtcolorbox{exoresolu}[1][]{popbase,colframe=poporange,colbacklower=popcream,
  segmentation style={popink,line width=1.2pt,dashed},
  before upper={\stepcounter{popexr}\popboxhead{Exercice résolu \thepopexr}{poporange}{#1}},
  before lower={\poppill{Solution}{popink}{popcream}\par\vspace{6pt}}}
% Exercice (non résolu en place) — la correction est dans la partie Corrigés
\newtcolorbox{exercice}[1][]{popbase,colframe=popink,
  before upper={\stepcounter{popexo}\popboxhead{Exercice \thepopexo}{popink}{#1}}}
% Corrigé
\newtcolorbox{corrige}[1][]{popbase,colframe=popgreen,colback=popgreenL,
  before upper={\popboxhead{Corrigé}{popgreen}{#1}}}
% Objectifs / prérequis / bilan
\newtcolorbox{objectifs}{popbase,colframe=popink,colback=poporangeL,
  before upper={\popboxhead{Objectifs de la leçon}{popink}{}}}
\newtcolorbox{prerequis}{popbase,colframe=popink,colback=popblueL,
  before upper={\popboxhead{Prérequis}{popblue}{}}}
\newtcolorbox{bilan}{popbase,colframe=popink,colback=white,boxrule=2.8pt,
  before upper={\popboxhead{Fiche bilan}{poporange}{L'essentiel à connaître pour l'examen}}}
% Figure encadrée avec légende
\newtcolorbox{popfigure}[1][]{enhanced,breakable=false,colback=white,colframe=popline,boxrule=1.4pt,arc=8pt,
  left=10pt,right=10pt,top=10pt,bottom=8pt,before skip=10pt,after skip=12pt,halign=center,
  lower separated=false,halign lower=center,
  fontlower=\popsemi\fontsize{9}{11.5}\selectfont\color{popmuted},
  #1}
\newcommand{\legende}[1]{\par\vspace{6pt}{\popsemi\fontsize{9}{11.5}\selectfont\color{popmuted}#1}}

% ============================================================
% Signature OnBuch+
% ============================================================
\newcommand{\poplogoinline}[1]{{\poplogo\fontsize{#1}{#1}\selectfont\color{popink}OnBuch\color{poporange}+}}
\newcommand{\signatureonbuch}{%
  \begin{tcolorbox}[enhanced,colback=popink,colframe=popink,boxrule=2.2pt,arc=10pt,
      drop shadow=poporange,left=14pt,right=14pt,top=10pt,bottom=10pt,before skip=0pt,after skip=0pt]
    \tikz[baseline=-0.7ex]\node[circle,fill=popgreen,draw=popcream,line width=1.3pt,minimum size=22pt,inner sep=0]{\popcheck{white}};%
    \hspace{9pt}\begin{minipage}[c]{0.62\linewidth}
      {\popxbold\fontsize{12}{13}\selectfont\color{popcream}Signé\ \textbullet\ {\poplogo OnBuch\color{poporange}+}}\par\vspace{2pt}
      {\raggedright\popsemi\fontsize{8}{10.5}\selectfont\color{popline}\MakeUppercase{Équipe pédagogique OnBuch+}\\\MakeUppercase{Conforme au programme officiel MINESEC}\par}
    \end{minipage}\hfill
    {\poplogo\fontsize{22}{22}\selectfont\color{popcream}OnBuch\color{poporange}+}
  \end{tcolorbox}%
}

% ============================================================
% Page de garde
% #1 titre · #2 matière · #3 niveau · #4 module · #5 n° leçon · #6 durée ·
% #7 description / objectifs (texte ou itemize)
% ============================================================
\newcommand{\pagegarde}[7]{%
  \thispagestyle{empty}
  \newgeometry{a4paper,margin=1.7cm,top=1.6cm,bottom=1.4cm}%
  \begin{tikzpicture}[remember picture,overlay]
    \fill[poporange!18] ([xshift=-3.2cm,yshift=-3.6cm]current page.north east) circle (4.6cm);
    \fill[poppurple!14] ([xshift=2.4cm,yshift=7.6cm]current page.south west) circle (3.8cm);
  \end{tikzpicture}%
  {\poplogo\fontsize{30}{30}\selectfont\color{popink}OnBuch\color{poporange}+}\hfill
  \raisebox{0.6ex}{\poppill{Cours signé \textbullet\ Premium}{popink}{popcream}}\par
  \vspace{1pt}\popsquiggle{poppurple}\par
  \vspace{8pt}
  \begin{tcolorbox}[enhanced,colback=poporange,colframe=popink,boxrule=2.8pt,arc=13pt,
      drop shadow=popink,left=18pt,right=18pt,top=12pt,bottom=13pt,after skip=10pt]
    \poppill{#2 \textbullet\ #3}{popink}{popcream}\hspace{5pt}\poppill{#4}{white}{popink}\par\vspace{8pt}
    {\popdisplay\fontsize{14}{14}\selectfont\color{popink}LEÇON #5}\par\vspace{4pt}
    {\raggedright\hyphenpenalty=10000\exhyphenpenalty=10000\popdisplay\fontsize{25}{27}\selectfont\color{popcream}\MakeUppercase{#1}\par}
    \vspace{8pt}
    {\popxbold\fontsize{9.5}{9.5}\selectfont\color{popink}\MakeUppercase{Durée conseillée : #6}}
  \end{tcolorbox}
  \begin{tcolorbox}[enhanced,colback=white,colframe=popink,boxrule=2.2pt,arc=10pt,
      drop shadow=popink,left=16pt,right=16pt,top=10pt,bottom=10pt,after skip=8pt]
    \poppill{Dans cette leçon}{poppurple}{white}\par\vspace{5pt}
    \popsemi\fontsize{9.3}{12.6}\selectfont\color{popink}#7
  \end{tcolorbox}
  \noindent\begin{minipage}[t]{0.487\linewidth}
    \begin{tcolorbox}[enhanced,colback=popgreen,colframe=popink,boxrule=2.2pt,arc=10pt,
        drop shadow=popink,left=11pt,right=11pt,top=10pt,bottom=10pt,height=3.1cm,valign=center]
      \tcbox[colback=white,colframe=popink,boxrule=1.3pt,arc=5pt,boxsep=0pt,nobeforeafter,
        left=3pt,right=3pt,top=3pt,bottom=3pt,valign=center]{\qrcode[height=1.9cm]{https://play.google.com/store/apps/details?id=com.luvvix.onbuch&hl=fr}}%
      \hspace{8pt}\begin{minipage}[c]{0.5\linewidth}
        {\popdisplay\fontsize{11}{12.5}\selectfont\color{white}\MakeUppercase{Télécharge OnBuch}}\par\vspace{4pt}
        {\raggedright\popsemi\fontsize{8.4}{11}\selectfont\color{white}Gratuit \textbullet\ Google Play\\Tuteur IA, annales, concours, résultats\par}
      \end{minipage}
    \end{tcolorbox}
  \end{minipage}\hfill
  \begin{minipage}[t]{0.487\linewidth}
    \begin{tcolorbox}[enhanced,colback=poppurple,colframe=popink,boxrule=2.2pt,arc=10pt,
        drop shadow=popink,left=11pt,right=11pt,top=10pt,bottom=10pt,height=3.1cm,valign=center]
      \tikz[baseline=-0.7ex]\node[circle,fill=white,draw=popink,line width=1.3pt,minimum size=1.6cm,inner sep=0]{\popdisplay\fontsize{24}{24}\selectfont\color{poppurple}+};%
      \hspace{8pt}\begin{minipage}[c]{0.6\linewidth}
        {\popdisplay\fontsize{11}{12.5}\selectfont\color{white}\MakeUppercase{Passe à OnBuch+}}\par\vspace{4pt}
        {\raggedright\popsemi\fontsize{8.4}{11}\selectfont\color{white}Tous les cours signés, Tuteur IA illimité, fiches PDF — onglet OnBuch+ de l'app\par}
      \end{minipage}
    \end{tcolorbox}
  \end{minipage}
  \vspace{8pt}
  \popcontenu
  \vfill
  \signatureonbuch
  \restoregeometry
  \newpage
}

% Rangée « Ce cours contient » (page de garde)
\newcommand{\poptile}[3]{% #1 couleur #2 gros texte #3 légende
  \begin{tcolorbox}[enhanced,colback=#1,colframe=popink,boxrule=2pt,arc=9pt,shadow={2.5pt}{-2.5pt}{0pt}{popink},
      left=6pt,right=6pt,top=9pt,bottom=9pt,halign=center,valign=center,height=1.75cm,width=\linewidth,nobeforeafter]
    {\popdisplay\fontsize{12.5}{13}\selectfont\color{white}\MakeUppercase{#2}}\par\vspace{3pt}
    {\centering\popsemi\fontsize{7.8}{9.5}\selectfont\color{white}#3\par}
  \end{tcolorbox}}
\newcommand{\popcontenu}{%
  \par\noindent{\popxboldls\fontsize{7.5}{7.5}\selectfont\color{popmuted}CE COURS SIGNÉ CONTIENT}\par\vspace{7pt}
  \noindent\begin{minipage}[t]{0.235\linewidth}\poptile{popblue}{Cours}{complet, illustré, pas à pas}\end{minipage}\hfill
  \begin{minipage}[t]{0.235\linewidth}\poptile{poporange}{Méthodes}{et exercices résolus}\end{minipage}\hfill
  \begin{minipage}[t]{0.235\linewidth}\poptile{poppink}{Exercices}{type examen + corrigés}\end{minipage}\hfill
  \begin{minipage}[t]{0.235\linewidth}\poptile{popgreen}{Fiche bilan}{l'essentiel à retenir}\end{minipage}\par}

% Sommaire
\newtcolorbox{sommaire}{enhanced,colback=white,colframe=popink,boxrule=2.2pt,arc=10pt,
  drop shadow=popink,left=18pt,right=18pt,top=13pt,bottom=13pt,before skip=6pt,after skip=20pt,
  before upper={\poppill{Sommaire}{poppurple}{white}\par\vspace{9pt}\popsemi\fontsize{10.6}{14.5}\selectfont\color{popink}}}

% Signature de fin de cours — poussée tout en bas de la dernière page
\newcommand{\findecours}{%
  \par\vfill\nopagebreak
  \begin{center}\popsquiggle{poporange}\end{center}\vspace{4pt}
  \signatureonbuch
}
\AtBeginDocument{\def\llbracket{\mathopen{[\mkern-3.5mu[}}\def\rrbracket{\mathclose{]\mkern-3.5mu]}}} % intervalles d'entiers (glyphe ⟦ absent de la police)
% Glyphes absents de Fira Math : redéfinis à partir de glyphes disponibles
\AtBeginDocument{%
  \def\setminus{\mathbin{\backslash}}\let\smallsetminus\setminus
  \def\vdots{\mathord{\vbox{\baselineskip3.2pt\lineskiplimit0pt\kern4pt\hbox{.}\hbox{.}\hbox{.}}}}%
  \def\ddots{\mathinner{\mkern1mu\raise6.5pt\hbox{.}\mkern2mu\raise3.6pt\hbox{.}\mkern2mu\raise0.7pt\hbox{.}\mkern1mu}}%
  \def\triangle{\mathord{\scalebox{0.9}{$\symup{\Delta}$}}}%
  \def\nabla{\mathord{\raisebox{\depth}{\scalebox{1}[-1]{$\symup{\Delta}$}}}}%
  \def\checkmark{\popcheck{popdark}}%
  \def\star{\mathbin{\ast}}\def\bigstar{\mathord{\ast}}%
  \def\diamond{\mathbin{\scalebox{0.6}{$\Diamond$}}}%
  \def\bigcup{\mathop{\mathchoice{\raisebox{-0.3em}{\scalebox{1.7}{$\cup$}}}{\scalebox{1.25}{$\cup$}}{\cup}{\cup}}\displaylimits}%
  \def\bigcap{\mathop{\mathchoice{\raisebox{-0.3em}{\scalebox{1.7}{$\cap$}}}{\scalebox{1.25}{$\cap$}}{\cap}{\cap}}\displaylimits}%
  \def\lesseqqgtr{\mathrel{\lessgtr}}\def\bigoplus{\mathop{\oplus}\displaylimits}%
}
\providecommand{\ssi}{si et seulement si} % abréviation fréquente
\AtBeginDocument{\providecommand{\wideparen}{\overparen}} % arc de cercle

\begin{document}

\pagegarde{Production et traduction de textes liés à la vie économique}{Allemand}{Tle A}{Chapitre 2 : Vie économique}{AL10}{2 h}{Cette leçon mixte entraîne les élèves à lire, commenter, produire et traduire des textes liés à la vie économique (lettre commerciale, compte rendu, article de presse). Elle révise trois points de grammaire essentiels au Bac — comparatif, conditionnel, passif — à travers le thème et la version.
\vspace{4pt}
\begin{multicols}{2}\small
\begin{itemize}
\item À la fin de cette leçon, je sais identifier les caractéristiques d'un texte économique (lettre commerciale, compte rendu, article)
\item je sais utiliser le lexique de la vie économique (entreprise, commerce, emploi, monnaie) dans un texte
\item je sais rédiger une lettre commerciale allemande conforme aux conventions
\item je sais rédiger un compte rendu bref et structuré en allemand
\item je sais former et employer le comparatif et le superlatif des adjectifs
\item je sais employer le conditionnel (Konjunktiv II avec würde + Infinitiv, hätte, wäre, könnte)
\end{itemize}\end{multicols}}

\begin{sommaire}
\begin{multicols}{2}
\begin{enumerate}
\item Texte support : lecture et compréhension
\item Lexique thématique : la vie économique
\item Grammaire 1 : le comparatif et le superlatif
\item Grammaire 2 : le conditionnel (Konjunktiv II)
\item Grammaire 3 : le passif (Vorgangspassiv)
\item La lettre commerciale et le compte rendu
\item Traduction : thème et version
\item Activité d'intégration
\item Méthodes et exercices résolus
\item Exercices
\item Corrigés des exercices
\item Fiche bilan
\end{enumerate}
\end{multicols}
\end{sommaire}

\begin{objectifs}
\begin{itemize}
\item À la fin de cette leçon, je sais identifier les caractéristiques d'un texte économique (lettre commerciale, compte rendu, article)
\item je sais utiliser le lexique de la vie économique (entreprise, commerce, emploi, monnaie) dans un texte
\item je sais rédiger une lettre commerciale allemande conforme aux conventions
\item je sais rédiger un compte rendu bref et structuré en allemand
\item je sais former et employer le comparatif et le superlatif des adjectifs
\item je sais employer le conditionnel (Konjunktiv II avec würde + Infinitiv, hätte, wäre, könnte)
\item je sais transformer une phrase active en passive (Vorgangspassiv) et la traduire correctement
\item je sais traduire de courtes phrases et un texte économique du français vers l'allemand (thème) et de l'allemand vers le français (version)
\item je sais relire et corriger ma production (orthographe, majuscules, ordre des mots, cas)
\end{itemize}
\end{objectifs}

\begin{prerequis}
\begin{itemize}
\item Déclinaison des articles et des adjectifs (Première)
\item Ordre des mots : place du verbe en position 2, verbes à particule, subordonnées (Première)
\item Présent, Präteritum et Perfekt des verbes forts et faibles (Première)
\item Notions de base sur le Konjunktiv II vues en début d'année de Terminale
\item Lexique général du travail et des professions (Seconde/Première)
\end{itemize}
\end{prerequis}

\newpage

\coursec{Texte support : lecture et compréhension}

\courssub{Situation d'accroche et problématique}

Aminatou, élève de Terminale A à Yaoundé, rêve de travailler dans une entreprise d'import-export avec l'Allemagne. Un jour, son oncle lui montre un courriel d'un fournisseur de Munich : elle comprend quelques mots, mais la structure de la lettre commerciale et les formulations polies lui échappent. De plus, un article de presse germanophone sur l'économie camerounaise doit être traduit pour un exposé. Comment rédiger une lettre commerciale correcte ? Comment traduire fidèlement un texte économique ? Cette leçon donne les clés : le lexique de la vie économique, les structures de la lettre commerciale et du compte rendu, et la révision du comparatif, du conditionnel et du passif.

Dans cette première étape, nous allons lire ensemble un article de presse économique, le comprendre en profondeur et y repérer les structures grammaticales que nous réviserons ensuite. Avant de lire, rappelons la bonne démarche.

\begin{methode}[Lire un texte économique en allemand]
\begin{enumerate}
\item \textbf{Identifier le type et la source} : s'agit-il d'un article de presse, d'un rapport, d'une publicité ? Qui écrit, pour qui, quand ?
\item \textbf{Repérer les chiffres et les acteurs} : les nombres, les pourcentages, les noms de pays, d'entreprises et de personnes donnent le squelette du texte.
\item \textbf{Dégager la thèse} : quelle est l'idée principale ? L'auteur est-il optimiste, critique, neutre ?
\item \textbf{Relever le vocabulaire-clé} : noter les mots du champ lexical de l'économie (\all{der Handel}, \all{die Nachfrage}, \all{der Gewinn}...) avec leur article et leur pluriel.
\end{enumerate}
\end{methode}

\begin{attention}[Le piège de la traduction mot à mot]
Ne traduis jamais mot à mot lors de la première lecture ! Cherche d'abord le \cle{sens global} : lis le titre, le premier et le dernier paragraphe, puis devine le sens des mots inconnus grâce au contexte. Le verbe \all{bekommen} ne signifie pas « devenir » mais « recevoir » (faux ami avec l'anglais \emph{to become}) ; \all{die Fabrik} n'est pas « la fabrique de tissu » mais « l'usine » ; \all{der Kredit} est « le crédit, le prêt », et non une « carte de crédit » (\all{die Kreditkarte}). Le sens global d'abord, le détail ensuite.
\end{attention}

\courssub{Lecture du texte : Der kamerunische Kaffee erobert den deutschen Markt}

\begin{textebox}[Article de presse économique — rubrique « Wirtschaft »]
\textbf{Der kamerunische Kaffee erobert den deutschen Markt}

Kamerun exportiert jedes Jahr Tausende Tonnen Kaffee nach Deutschland. Die Nachfrage ist in den letzten Jahren größer geworden, weil die Qualität der Bohnen besser ist als früher. Besonders der Arabica-Kaffee aus der Westregion, der auf vulkanischen Böden wächst, wird von deutschen Röstereien sehr geschätzt.

In Hamburg, dem wichtigsten Kaffeehafen Europas, wird der kamerunische Kaffee entladen, geprüft und geröstet. Von dort aus wird er in ganz Deutschland verkauft. „Der kamerunische Kaffee hat einen stärkeren Geschmack als viele andere Sorten“, erklärt eine Händlerin aus Hamburg. „Die deutschen Kunden mögen ihn immer mehr.“

Doch es gibt auch Probleme. Die Bauern verdienen oft weniger Geld als die Zwischenhändler. Wenn die Bauern mehr Unterstützung bekämen, könnten sie ihre Produktion steigern und bessere Maschinen kaufen. „Wir bräuchten mehr Kredite und bessere Straßen“, sagt ein Bauer aus Dschang. „Dann würden wir mehr Kaffee exportieren und mehr Gewinn machen.“

Die Regierung in Yaoundé möchte den Handel mit Deutschland ausbauen. Neue Verträge wurden letztes Jahr unterschrieben, und die Ausfuhr von Kakao und Bananen soll ebenfalls erhöht werden. Experten sind optimistisch: Wenn der Wechselkurs stabil bleibt, wird der kamerunische Kaffee in Zukunft noch bekannter werden.

Der Erfolg des Kaffees zeigt: Die wirtschaftliche Zusammenarbeit zwischen Kamerun und Deutschland hat eine große Zukunft.
\end{textebox}

\textbf{Glossaire} (à noter avec article et pluriel) :

\begin{center}
\begin{tabularx}{\linewidth}{|X|X|}
\hline
\rowcolor{popblueL} \textbf{Deutsch} & \textbf{Français} \\
\hline
die Ausfuhr, -en & l'exportation \\
\hline
der Handel & le commerce \\
\hline
die Nachfrage (meist Singular) & la demande \\
\hline
die Bohne, -n & le grain (de café) \\
\hline
der Boden, die Böden & le sol, la terre \\
\hline
die Rösterei, -en & la torréfaction (atelier) \\
\hline
der Geschmack, die Geschmäcker & le goût \\
\hline
der Zwischenhändler, - & l'intermédiaire (commerçant) \\
\hline
der Kredit, -e & le crédit, le prêt \\
\hline
der Gewinn, -e & le bénéfice, le gain \\
\hline
der Wechselkurs, -e & le taux de change \\
\hline
die Zusammenarbeit (Singular) & la coopération \\
\hline
der Vertrag, die Verträge & le contrat \\
\hline
die Unterstützung, -en & le soutien, l'aide \\
\hline
\end{tabularx}
\end{center}

\courssub{Repérage du type de texte}

Avant toute question, identifions le document :

\begin{itemize}
\item \textbf{Type de texte} : un \cle{article de presse économique} (\all{ein Zeitungsartikel}), rubrique « Wirtschaft ».
\item \textbf{Source} : un journal germanophone (type quotidien national comme \emph{Die Zeit} ou \emph{Süddeutsche Zeitung}).
\item \textbf{Thème} : les exportations de café camerounais vers l'Allemagne et la coopération économique entre les deux pays.
\item \textbf{Ton} : globalement \cle{optimiste} et informatif, mais l'article signale aussi des problèmes (revenus faibles des paysans, manque de crédits).
\item \textbf{Structure typique} : titre accrocheur (\all{erobert} = « conquiert », image de conquête), faits et chiffres, citations de témoins, problèmes, perspectives d'avenir.
\end{itemize}

\begin{aretenir}[Le titre d'un article économique]
Le titre utilise souvent une \cle{métaphore} pour attirer le lecteur : \all{Der Kaffee erobert den Markt} (« le café conquiert le marché »). Pour le traduire, cherche une image équivalente en français : « Le café camerounais à la conquête du marché allemand ». Évite la traduction plate « conquiert » si une formulation française plus naturelle existe.
\end{aretenir}

\courssub{Compréhension globale : qui ? quoi ? où ? pourquoi ?}

\begin{exoresolu}[Questions de compréhension globale]
Réponds en français, en phrases complètes :
\begin{enumerate}
\item De quoi parle cet article ?
\item Quels sont les acteurs mentionnés ?
\item Où se déroule l'action principale ?
\item Pourquoi le café camerounais a-t-il du succès en Allemagne ?
\end{enumerate}
\tcblower
\textbf{Solution détaillée :}
\begin{enumerate}
\item L'article parle des exportations de café du Cameroun vers l'Allemagne, de leur succès croissant, des difficultés des producteurs et des perspectives de coopération économique.
\item Les acteurs sont : les paysans camerounais (\all{die Bauern}), les intermédiaires (\all{die Zwischenhändler}), les torréfacteurs et commerçants allemands (\all{die Röstereien}, \all{eine Händlerin aus Hamburg}), les clients allemands, le gouvernement camerounais et les experts.
\item L'action se déroule entre le Cameroun (la région de l'Ouest, Dschang, Yaoundé) et l'Allemagne (le port de Hambourg, puis tout le pays).
\item Parce que la qualité des grains s'est améliorée (\all{die Qualität ist besser als früher}) et que son goût est plus prononcé que celui d'autres variétés (\all{einen stärkeren Geschmack}).
\end{enumerate}
\end{exoresolu}

\courssub{Compréhension détaillée : chiffres, acteurs, causes et conséquences}

\begin{exercice}[À toi de chercher dans le texte]
\begin{enumerate}
\item Quelle quantité de café le Cameroun exporte-t-il chaque année ?
\item Quel type de café est particulièrement apprécié, et où pousse-t-il ?
\item Quelles sont les trois opérations faites à Hambourg sur le café ?
\item Quels sont les deux problèmes des paysans camerounais ?
\item Que demande le paysan de Dschang ?
\item Quels autres produits le Cameroun veut-il exporter davantage ?
\item À quelle condition le café camerounais deviendra-t-il encore plus connu ?
\end{enumerate}
\end{exercice}

\begin{corrige}[Exercice de compréhension détaillée]
\begin{enumerate}
\item Des milliers de tonnes (\all{Tausende Tonnen}) chaque année.
\item Le café Arabica, qui pousse sur les sols volcaniques de la région de l'Ouest (\all{aus der Westregion, auf vulkanischen Böden}).
\item Il est déchargé (\all{entladen}), contrôlé (\all{geprüft}) et torréfié (\all{geröstet}).
\item Ils gagnent souvent moins d'argent que les intermédiaires, et ils manquent de soutien (crédits, machines, routes).
\item Il demande plus de crédits et de meilleures routes : \all{Wir bräuchten mehr Kredite und bessere Straßen.}
\item Le cacao et les bananes (\all{die Ausfuhr von Kakao und Bananen}).
\item Si le taux de change reste stable : \all{Wenn der Wechselkurs stabil bleibt.}
\end{enumerate}
\end{corrige}

\begin{savaistu}[Le Cameroun et l'Allemagne, partenaires de longue date]
L'Allemagne est l'un des premiers partenaires commerciaux du Cameroun en Europe. Le café et le cacao camerounais sont très appréciés à Hambourg, premier port caféier d'Europe : c'est là qu'arrivent la plupart des cafés destinés au marché allemand. D'ailleurs, le Cameroun fut un protectorat allemand de 1884 à 1916 ; le mot allemand \all{Kamerun} désigne toujours le pays, et certains mots allemands ont laissé des traces dans la mémoire collective camerounaise.
\end{savaistu}

\courssub{Observation guidée : grammaire en contexte}

Le texte contient exactement les trois structures que nous réviserons dans cette leçon. Cherche-les d'abord toi-même, puis vérifie avec le tableau.

\begin{center}
\begin{tabularx}{\linewidth}{|X|X|X|}
\hline
\rowcolor{popblueL} \textbf{Structure} & \textbf{Exemple du texte} & \textbf{Traduction} \\
\hline
Comparatif & \all{Die Nachfrage ist größer geworden.} & La demande est devenue plus grande. \\
\hline
Comparatif & \all{besser ist als früher} & est meilleure qu'autrefois \\
\hline
Comparatif & \all{einen stärkeren Geschmack als viele andere Sorten} & un goût plus fort que beaucoup d'autres variétés \\
\hline
Comparatif & \all{weniger Geld als die Zwischenhändler} & moins d'argent que les intermédiaires \\
\hline
Conditionnel (Konj. II) & \all{Wenn die Bauern mehr Unterstützung bekämen, könnten sie ihre Produktion steigern.} & Si les paysans recevaient plus de soutien, ils pourraient augmenter leur production. \\
\hline
Conditionnel (Konj. II) & \all{Wir bräuchten mehr Kredite.} & Nous aurions besoin de plus de crédits. \\
\hline
Conditionnel (Konj. II) & \all{Dann würden wir mehr Kaffee exportieren.} & Alors nous exporterions plus de café. \\
\hline
Passif & \all{wird von deutschen Röstereien sehr geschätzt} & est très apprécié par les torréfacteurs allemands \\
\hline
Passif & \all{wird entladen, geprüft und geröstet} & est déchargé, contrôlé et torréfié \\
\hline
Passif (Präteritum) & \all{Neue Verträge wurden unterschrieben.} & De nouveaux contrats ont été signés. \\
\hline
Passif (avec modal) & \all{die Ausfuhr soll erhöht werden} & l'exportation doit être augmentée \\
\hline
\end{tabularx}
\end{center}

\begin{propriete}[Repérer les trois structures dans un texte]
\begin{itemize}
\item \textbf{Comparatif} : adjectif + \all{-er} + \all{als} → \all{größer als}, \all{besser als} (irrégulier : \all{gut – besser – am besten} ; \all{viel – mehr – am meisten}).
\item \textbf{Conditionnel (Konjunktiv II)} : \all{würde} + infinitif, ou formes propres \all{könnte}, \all{bräuchte}, \all{bekäme}, \all{hätte}, \all{wäre}.
\item \textbf{Passif (Vorgangspassiv)} : \all{werden} (conjugué) + participe II ; l'agent est introduit par \all{von} + datif : \all{wird von deutschen Röstereien geschätzt}.
\end{itemize}
\end{propriete}

\begin{attention}[Erreurs fréquentes des francophones]
\begin{itemize}
\item Ne confonds pas \all{als} (comparaison d'inégalité : \all{größer als} = « plus grand que ») et \all{wie} (égalité : \all{so groß wie} = « aussi grand que »). Après \all{als} et \all{wie}, le nom reprend le cas exigé par sa fonction dans la phrase : \all{Die Bauern verdienen weniger als die Zwischenhändler} (nominatif, car on compare deux sujets).
\item Au passif, le français « par » se traduit par \all{von} (+ datif) pour l'agent, mais par \all{durch} (+ accusatif) pour le moyen : \all{Der Kaffee wird per Schiff transportiert.} ou \all{Der Kaffee wird durch den Hafen Hamburg importiert.} (« Le café est transporté par bateau / via le port de Hambourg. »)
\item \all{bekommen} = « recevoir », jamais « devenir » : \all{Wenn die Bauern mehr Unterstützung bekämen} = « Si les paysans \emph{recevaient} plus de soutien ». « Devenir » se dit \all{werden}.
\item Attention à la place du verbe dans la subordonnée : \all{..., weil die Qualität besser ist als früher} — le verbe conjugué \all{ist} va à la fin, et \all{als früher} se place après lui.
\end{itemize}
\end{attention}

\courssub{Le champ lexical de l'économie : carte mentale}

Le texte mobilise le champ lexical de \all{die Wirtschaft} (l'économie). Organisons-le visuellement pour mieux le mémoriser.

\begin{popfigure}
\begin{tikzpicture}[mindmap, grow cyclic, every node/.style={concept, align=center, font=\small}, root concept/.append style={concept color=popblue, text=white, font=\bfseries}, level 1/.append style={level distance=3.2cm, sibling angle=72}]
\node[root concept]{Wirtschaft}
  child[concept color=poporange] { node[align=center]{Handel\\ \footnotesize der Export\\ \footnotesize die Ausfuhr\\ \footnotesize der Vertrag} }
  child[concept color=popgreen] { node[align=center]{Unternehmen\\ \footnotesize die Firma\\ \footnotesize der Gewinn\\ \footnotesize der Kredit} }
  child[concept color=poppurple] { node[align=center]{Arbeit\\ \footnotesize der Bauer\\ \footnotesize der Händler\\ \footnotesize der Kunde} }
  child[concept color=popgold] { node[align=center]{Geld\\ \footnotesize die Währung\\ \footnotesize der Wechselkurs\\ \footnotesize der Preis} }
  child[concept color=poppink] { node[align=center]{Märkte\\ \footnotesize der Markt\\ \footnotesize die Nachfrage\\ \footnotesize das Angebot} };
\end{tikzpicture}
\legende{Carte mentale du champ lexical de l'économie (die Wirtschaft)}
\end{popfigure}

\courssub{Le circuit de l'exportation : du producteur au consommateur}

Pour bien comprendre le parcours du café décrit dans l'article, voici le circuit de l'exportation entre le Cameroun et l'Allemagne.

\begin{popfigure}
\begin{tikzpicture}[node distance=0.9cm, every node/.style={align=center, font=\small}]
\node[draw, fill=popgreenL, rounded corners, minimum width=2.6cm, minimum height=1cm, align=center] (prod){\textbf{Produzent}\\ der Bauer in Dschang};
\node[draw, fill=poporangeL, rounded corners, minimum width=2.6cm, minimum height=1cm, right=of prod, align=center] (exp){\textbf{Exporteur}\\ Hafen Douala};
\node[draw, fill=popblueL, rounded corners, minimum width=2.6cm, minimum height=1cm, right=of exp, align=center] (imp){\textbf{Importeur}\\ Hafen Hamburg};
\node[draw, fill=poppinkL, rounded corners, minimum width=2.6cm, minimum height=1cm, right=of imp, align=center] (kunde){\textbf{Kunde}\\ in Deutschland};
\draw[-{Stealth[length=3mm]}, very thick, popdark] (prod) -- (exp);
\draw[-{Stealth[length=3mm]}, very thick, popdark] (exp) -- (imp) node[midway, above, font=\footnotesize]{Schiff};
\draw[-{Stealth[length=3mm]}, very thick, popdark] (imp) -- (kunde);
\end{tikzpicture}
\legende{Le circuit de l'exportation du café : Cameroun → Allemagne}
\end{popfigure}

Ce schéma permet de réutiliser le passif : \all{Der Kaffee wird in Douala verladen, nach Hamburg transportiert, dort geröstet und schließlich an die Kunden verkauft.} (« Le café est chargé à Douala, transporté à Hambourg, torréfié sur place et finalement vendu aux clients. ») Remarque la chaîne de participes II (\all{verladen, transportiert, geröstet, verkauft}) avec un seul auxiliaire \all{wird} : c'est une structure très fréquente dans les textes économiques.

\courssub{Reformulation des idées principales}

Savoir reformuler (\all{umformulieren}) avec ses propres mots est la compétence-clé du commentaire de texte et de la traduction. Voici un modèle.

\begin{exemplebox}[Reformulation : du texte à tes propres mots]
\begin{itemize}
\item \textbf{Idée 1} — Texte : \all{Die Nachfrage ist größer geworden, weil die Qualität besser ist als früher.}\par
Reformulation : \all{Immer mehr deutsche Kunden kaufen kamerunischen Kaffee, denn er schmeckt heute besser.} (« De plus en plus de clients allemands achètent du café camerounais, car il a meilleur goût aujourd'hui. »)
\item \textbf{Idée 2} — Texte : \all{Die Bauern verdienen oft weniger Geld als die Zwischenhändler.}\par
Reformulation : \all{Die Produzenten bekommen für ihre Arbeit nicht so viel Geld wie die Händler.} (« Les producteurs ne reçoivent pas autant d'argent pour leur travail que les commerçants. »)
\item \textbf{Idée 3} — Texte : \all{Die Zusammenarbeit zwischen Kamerun und Deutschland hat eine große Zukunft.}\par
Reformulation : \all{Kamerun und Deutschland werden in Zukunft noch mehr Handel treiben.} (« Le Cameroun et l'Allemagne commerceront encore davantage à l'avenir. »)
\end{itemize}
\end{exemplebox}

\begin{methode}[Reformuler une phrase allemande en quatre étapes]
\begin{enumerate}
\item \textbf{Comprendre} la phrase source (sujet, verbe, compléments).
\item \textbf{Remplacer} les mots par des synonymes : \all{die Nachfrage wächst} → \all{die Kunden kaufen mehr}.
\item \textbf{Changer la structure} : transformer une subordonnée en deux phrases, ou une phrase passive en active : \all{Der Kaffee wird gekauft} → \all{Die Firmen kaufen den Kaffee}.
\item \textbf{Vérifier} : le sens doit rester identique, la grammaire correcte (place du verbe, cas, majuscules des noms).
\end{enumerate}
\end{methode}

\begin{exoresolu}[Reformule avec tes propres mots]
Reformule en allemand la phrase suivante, extraite du texte :\par
\all{Wenn die Bauern mehr Unterstützung bekämen, könnten sie ihre Produktion steigern.}
\tcblower
\textbf{Solution détaillée :}\par
Étape 1 : sens = condition irréelle (les paysans ne reçoivent pas assez de soutien).\par
Étapes 2-3 : synonymes et autre structure, par exemple :\par
\all{Mit mehr Hilfe könnten die Bauern mehr Kaffee produzieren.} (« Avec plus d'aide, les paysans pourraient produire plus de café. »)\par
Autre possibilité : \all{Die Bauern brauchen mehr Unterstützung, um ihre Produktion zu steigern.} (« Les paysans ont besoin de plus de soutien pour augmenter leur production. »)\par
Étape 4 : vérification — le sens est conservé, le verbe conjugué est à la deuxième place (\all{könnten}, \all{brauchen}), et les noms \all{Hilfe}, \all{Unterstützung}, \all{Produktion} prennent la majuscule.
\end{exoresolu}

\begin{experience}[Activité de classe : le journaliste économique]
Par groupes de trois : l'un joue le paysan de Dschang, l'autre la commerçante de Hambourg, le troisième le journaliste. Le journaliste pose trois questions en allemand (avec \all{Wie}, \all{Warum}, \all{Was}) ; les deux témoins répondent en réutilisant au moins un comparatif, une phrase au conditionnel et une phrase au passif. Ensuite, chaque groupe rédige un résumé de cinq phrases de l'article avec ses propres mots.
\end{experience}

\begin{aretenir}[Bilan de l'étape 1]
\begin{itemize}
\item Un article économique se lit en quatre étapes : type et source, chiffres et acteurs, thèse, vocabulaire-clé.
\item Le texte étudié présente les exportations de café camerounais vers l'Allemagne, leurs causes (meilleure qualité) et leurs obstacles (revenus des paysans, crédits, routes).
\item Il illustre les trois révisions grammaticales de la leçon : le \cle{comparatif} (\all{größer als}), le \cle{conditionnel} (\all{würden}, \all{könnten}, \all{bekämen}) et le \cle{passif} (\all{wird geschätzt}, \all{wurden unterschrieben}).
\item Reformuler, c'est dire la même chose autrement, sans jamais trahir le sens.
\end{itemize}
\end{aretenir}

\coursec{Lexique thématique : la vie économique}

Dans la section précédente, nous avons lu et commenté un texte de presse sur la vie économique. Vous avez remarqué que ce type de texte repose sur un vocabulaire précis et récurrent : entreprise, commerce, banque, emploi, production. C'est ce vocabulaire que nous allons maintenant organiser, approfondir et mémoriser, car il constitue la boîte à outils indispensable pour produire (lettre commerciale, compte rendu) et pour traduire (thème et version) des textes économiques. Pour chaque champ lexical, nous observerons d'abord les mots en contexte, puis nous les fixerons dans un tableau avec article et pluriel — car en allemand, on n'apprend jamais un nom sans son genre.

\begin{propriete}[La majuscule obligatoire des noms]
En allemand, \cle{tous les noms communs} prennent une majuscule, quelle que soit leur place dans la phrase : \all{der Markt}, \all{die Wirtschaft}, \all{das Produkt}, \all{die Arbeitslosigkeit}. C'est une différence majeure avec le français, où seuls les noms propres sont capitalisés. En production écrite comme en thème, l'oubli de la majuscule est une faute comptée. Astuce : tout mot précédé de \all{der}, \all{die} ou \all{das} (ou qui pourrait l'être) s'écrit avec une majuscule.
\end{propriete}

\courssub{Champ 1 : L'entreprise et sa structure}

Observons d'abord deux phrases qui anticipent sur la grammaire du passif (voir Grammaire 3) et sur le comparatif (voir Grammaire 1) :

\all{Die Firma wurde 1998 in Douala gegründet.} « L'entreprise a été fondée en 1998 à Douala. » (Passif \all{Präteritum})

\all{Die Chefin leitet den Betrieb mit 50 Mitarbeitern.} « La directrice dirige l'entreprise avec 50 employés. »

Trois mots se recouvrent partiellement : \all{das Unternehmen} (l'entreprise, terme général et abstrait), \all{die Firma} (la firme, la société, terme commercial courant) et \all{der Betrieb} (l'exploitation, l'établissement, l'usine vue comme lieu de travail). Attention au genre : on dit \all{der Betrieb} mais \all{die Firma}.

\begin{center}
\begin{tabularx}{\linewidth}{|X|X|X|}
\hline
\rowcolor{popblueL} Deutsch (Sing., Pl.) & Français & Remarque \\
\hline
das Unternehmen, - & l'entreprise & terme général, neutre, pas de pluriel usuel \\
\hline
die Firma, die Firmen & la firme, la société & terme commercial \\
\hline
der Betrieb, die Betriebe & l'exploitation, l'établissement & lieu de production \\
\hline
der Chef, die Chefs / die Chefinnen & le chef, la cheffe & féminin en \all{-in}, plur. \all{-nen} \\
\hline
der Geschäftsführer, die Geschäftsführer & le gérant, le directeur général & composé : \all{Geschäft} + \all{Führer} \\
\hline
der Mitarbeiter, die Mitarbeiter & l'employé, le collaborateur & pluriel identique (n-décl.) \\
\hline
die Belegschaft, die Belegschaften & le personnel & nom collectif, féminin \\
\hline
gründen & fonder & verbe faible : \all{gründete, hat gegründet} \\
\hline
leiten & diriger & \all{leitete, hat geleitet} \\
\hline
\end{tabularx}
\end{center}

\courssub{Champ 2 : Le commerce international}

\all{Kamerun exportiert Kakao und Kaffee nach Europa.} « Le Cameroun exporte du cacao et du café vers l'Europe. »

\all{Die Konkurrenz auf dem Weltmarkt ist stark.} « La concurrence sur le marché mondial est forte. »

Notez la double appellation : \all{der Export} et \all{die Ausfuhr} désignent l'exportation ; \all{der Import} et \all{die Einfuhr} désignent l'importation. Les verbes correspondants sont \all{exportieren} et \all{importieren}, réguliers, avec le préfixe \all{-ieren} typique des emprunts internationaux (toujours faibles, sans \all{ge-} au participe : \all{hat exportiert}).

\begin{center}
\begin{tabularx}{\linewidth}{|X|X|X|}
\hline
\rowcolor{popblueL} Deutsch (Sing., Pl.) & Français & Remarque \\
\hline
der Handel & le commerce & invariable au singulier courant \\
\hline
der Export / die Ausfuhr & l'exportation & deux synonymes \\
\hline
der Import / die Einfuhr & l'importation & deux synonymes \\
\hline
exportieren / importieren & exporter / importer & faibles, sans \all{ge-} au Partizip II \\
\hline
der Kunde, die Kunden & le client & n-déclinaison : \all{den Kunden} (Akk./Dat.) \\
\hline
der Lieferant, die Lieferanten & le fournisseur & n-déclinaison \\
\hline
der Markt, die Märkte & le marché & Umlaut au pluriel \\
\hline
die Konkurrenz & la concurrence & aussi « les concurrents » (coll.) \\
\hline
die Bestellung, die Bestellungen & la commande & \all{eine Bestellung aufgeben} \\
\hline
der Katalog, die Kataloge & le catalogue & \\
\hline
\end{tabularx}
\end{center}

\begin{exemplebox}[Exemples en contexte]
\all{Die Firma exportiert Kakao nach Europa.} « L'entreprise exporte du cacao vers l'Europe. »

\all{Unser Lieferant schickt die Waren jede Woche.} « Notre fournisseur envoie les marchandises chaque semaine. »

\all{Der Kunde ist mit der Qualität zufrieden.} « Le client est satisfait de la qualité. » (\all{zufrieden mit} + Dat.)
\end{exemplebox}

\courssub{Champ 3 : L'argent, la banque et la comptabilité}

\all{Das Unternehmen hat einen Kredit bei der Bank aufgenommen.} « L'entreprise a contracté un crédit auprès de la banque. »

\all{Im letzten Jahr machte die Firma einen großen Gewinn.} « L'année dernière, la société a réalisé un gros bénéfice. »

Distinguez bien \all{der Gewinn} (le bénéfice, le gain) et \all{der Verlust} (la perte) : ce sont des contraires. \all{Die Währung} désigne la monnaie d'un pays (l'euro, le franc CFA) ; \all{das Geld} est l'argent en général. Attention à \all{die Bank} (la banque) vs \all{die Bank, die Bänke} (le banc) !

\begin{center}
\begin{tabularx}{\linewidth}{|X|X|X|}
\hline
\rowcolor{popblueL} Deutsch (Sing., Pl.) & Français & Remarque / Collocation \\
\hline
das Geld, die Gelder & l'argent & pluriel rare (\all{Gelder} = fonds) \\
\hline
die Währung, die Währungen & la monnaie, la devise & \all{Wechselkurs} = taux de change \\
\hline
der Preis, die Preise & le prix & aussi « le prix » (récompense) \\
\hline
der Gewinn, die Gewinne & le bénéfice, le gain & contraire : \all{der Verlust} \\
\hline
der Verlust, die Verluste & la perte & \\
\hline
die Bank, die Banken & la banque & \all{die Bank, die Bänke} = le banc ! \\
\hline
das Konto, die Konten & le compte & \all{ein Konto eröffnen / führen} \\
\hline
der Kredit, die Kredite & le crédit, le prêt & \all{einen Kredit aufnehmen / gewähren} \\
\hline
die Rechnung, die Rechnungen & la facture & \all{eine Rechnung stellen / bezahlen} \\
\hline
der Zins, die Zinsen & l'intérêt (taux) & \all{Zinsen zahlen} \\
\hline
investieren & investir & faible, sans \all{ge-} (\all{hat investiert}) \\
\hline
\end{tabularx}
\end{center}

\courssub{Champ 4 : L'emploi et les ressources humaines}

\all{Die Arbeitslosigkeit ist gesunken.} « Le chômage a baissé. » (Perfekt avec \all{sein}, verbe fort \all{sinken})

\all{Die Firma stellt neue Arbeiter ein.} « L'entreprise embauche de nouveaux ouvriers. » (Verbe à particule \all{einstellen})

Notez la paire de contraires : \all{einstellen} (embaucher, particule séparable) et \all{entlassen} (licencier, verbe fort : \all{entlässt, entließ, hat entlassen}). Distinguez \all{der Arbeiter} (l'ouvrier, travail manuel) et \all{der Angestellte} (l'employé, salarié de bureau) — ce dernier se décline comme un adjectif : \all{ein Angestellter, der Angestellte, die Angestellten}.

\begin{center}
\begin{tabularx}{\linewidth}{|X|X|X|}
\hline
\rowcolor{popblueL} Deutsch (Sing., Pl.) & Français & Remarque \\
\hline
die Arbeit, die Arbeiten & le travail & \\
\hline
der Arbeitsplatz, die Arbeitsplätze & l'emploi, le poste de travail & \all{Arbeitsplätze schaffen} \\
\hline
die Arbeitslosigkeit & le chômage & nom abstrait, sans pluriel \\
\hline
der Arbeiter, die Arbeiter & l'ouvrier & travail manuel \\
\hline
der Angestellte, die Angestellten & l'employé (bureau) & déclinaison adjectivale (n-décl.) \\
\hline
das Gehalt, die Gehälter & le salaire (employé) & mensuel, fixe, Umlaut plur. \\
\hline
der Lohn, die Löhne & le salaire (ouvrier) & souvent horaire, Umlaut plur. \\
\hline
die Bewerbung, die Bewerbungen & la candidature & \all{sich bewerben um} + Akk. \\
\hline
der Arbeitsvertrag, die Arbeitsverträge & le contrat de travail & \\
\hline
einstellen & embaucher & particule séparable : \all{stellt ... ein} \\
\hline
entlassen & licencier & fort : \all{entlässt, entließ, hat entlassen} \\
\hline
kündigen & démissionner / résilier & \all{kündigen} (faible) + Dat. (\all{dem Chef}) \\
\hline
\end{tabularx}
\end{center}

\begin{attention}[Gehalt ou Lohn ? bekommen ou werden ?]
Deux pièges classiques pour les francophones :
\begin{itemize}
\item \all{das Gehalt} est le salaire mensuel d'un employé ou d'un cadre ; \all{der Lohn} est le salaire (souvent horaire ou hebdomadaire) d'un ouvrier. On dit donc \all{das Gehalt des Angestellten} mais \all{der Lohn des Arbeiters}.
\item \all{bekommen} signifie \cle{recevoir} et non « devenir » : \all{Ich bekomme mein Gehalt.} « Je reçois mon salaire. » « Devenir » se dit \all{werden} : \all{Er wird Chef.} « Il devient chef. » Ne traduisez jamais « devenir » par \all{bekommen} !
\end{itemize}
\end{attention}

\courssub{Champ 5 : La production, l'offre et la demande}

\all{Die Nachfrage nach Kaffee steigt, aber das Angebot ist klein.} « La demande de café augmente, mais l'offre est faible. » (Comparatif d'infériorité \all{kleiner als} — voir Grammaire 1)

\all{Die Fabrik produziert Waren von guter Qualität.} « L'usine produit des marchandises de bonne qualité. » (Déclinaison adjectivale après préposition — \all{von guter Qualität})

\all{Die Nachfrage} (la demande) et \all{das Angebot} (l'offre) forment le couple central de l'économie de marché. \all{Die Ware} désigne la marchandise concrète ; \all{das Produkt} est le produit (résultat de la production) ; \all{verbrauchen} signifie « consommer » (le consommateur : \all{der Verbraucher}).

\begin{center}
\begin{tabularx}{\linewidth}{|X|X|X|}
\hline
\rowcolor{popblueL} Deutsch (Sing., Pl.) & Français & Remarque / Structure \\
\hline
die Produktion & la production & \all{in Produktion gehen} \\
\hline
produzieren & produire & faible, sans \all{ge-} \\
\hline
das Produkt, die Produkte & le produit & résultat de la production \\
\hline
die Ware, die Waren & la marchandise & concret, stockable \\
\hline
die Qualität, die Qualitäten & la qualité & \all{von hoher Qualität} \\
\hline
die Nachfrage & la demande & \all{die Nachfrage nach} + Dat. \\
\hline
das Angebot, die Angebote & l'offre & aussi « offre commerciale écrite » \\
\hline
verbrauchen & consommer & \all{der Verbraucher} = le consommateur \\
\hline
der Verbraucher, die Verbraucher & le consommateur & n-déclinaison \\
\hline
\end{tabularx}
\end{center}

\begin{attention}[Faux amis à retenir absolument]
\begin{itemize}
\item \all{die Fabrik} = l'usine (et non « fabrique » au sens religieux) : \all{Er arbeitet in einer Fabrik.} « Il travaille dans une usine. »
\item \all{das Angebot} = l'offre (commerciale), pas seulement « l'offre » abstraite : \all{Ihr Angebot vom 5. Mai} « votre offre du 5 mai ».
\item \all{aktuell} = actuel (adjectif), qui a lieu maintenant ; « actuellement » (adverbe) se traduit par \all{zurzeit} ou \all{gegenwärtig}. \all{aktuell} ne s'emploie pas comme adverbe : \all{die aktuellen Preise} « les prix actuels ».
\item \all{konkurrent} = concurrent (adj./nom) ; \all{die Konkurrenz} = la concurrence (collectif) ou les concurrents.
\end{itemize}
\end{attention}

\courssub{Expressions utiles pour la lettre commerciale}

Pour rédiger ou traduire une lettre commerciale (voir section « La lettre commerciale »), mémorisez ces formules figées. Notez l'emploi du \cle{Konjunktiv II} (\all{möchten}, \all{wären}) pour la politesse (voir Grammaire 2).

\begin{exemplebox}[Formules types de la correspondance commerciale]
\all{Sehr geehrte Damen und Herren,} « Madame, Monsieur, » (formule d'appel sans nom)

\all{Sehr geehrter Herr Müller,} « Monsieur Müller, » (avec nom)

\all{Wir bedanken uns für Ihr Angebot vom 10. Mai.} « Nous vous remercions de votre offre du 10 mai. »

\all{Wir bedanken uns für Ihre Bestellung.} « Nous vous remercions de votre commande. »

\all{Wir möchten 100 Säcke Kakao bestellen.} « Nous souhaitons commander 100 sacs de cacao. » (\all{möchten} = Konj. II de \all{mögen})

\all{Bitte senden Sie uns Ihren Katalog.} « Veuillez nous envoyer votre catalogue. » (Impératif formel)

\all{Wir bitten um ein Angebot für ...} « Nous vous prions de bien vouloir nous faire une offre pour ... »

\all{Mit freundlichen Grüßen} « Cordialement / Nous vous prions d'agréer nos salutations distinguées. » (Sans point final !)
\end{exemplebox}

Notez la construction réflexive \all{sich bedanken für} + accusatif (se remercier de quelque chose) et la formule de politesse finale \all{Mit freundlichen Grüßen}, sans point final, suivie du nom de l'entreprise et de la signature.

\courssub{Organigramme d'une entreprise (Visualisation)}

\begin{popfigure}
\begin{center}
\begin{tikzpicture}[
  chef/.style={rectangle, rounded corners, draw=popblue, fill=popblueL, align=center, font=\bfseries, minimum width=3.5cm, minimum height=1cm},
  abt/.style={rectangle, rounded corners, draw=poporange, fill=poporangeL, align=center, minimum width=3.2cm, minimum height=0.9cm},
  mit/.style={rectangle, rounded corners, draw=popgreen, fill=popgreenL, align=center, minimum width=2.8cm, minimum height=0.8cm},
  lin/.style={-{Stealth}, thick, popmuted}
]
\node[chef, align=center] (chef){Geschäftsführung\\ (Chef / Chefin)};
\node[abt, below left=1.5cm and 0.5cm of chef, align=center] (verkauf){Vertrieb / Verkauf\\ (Abteilung Verkauf)};
\node[abt, below=1.5cm of chef, align=center] (prod){Produktion\\ (Abteilung Produktion)};
\node[abt, below right=1.5cm and 0.5cm of chef, align=center] (buch){Finanzwesen / Buchhaltung\\ (Abteilung Buchhaltung)};
\node[mit, below=1cm of verkauf] (m1) {Mitarbeiter / Angestellte};
\node[mit, below=1cm of prod] (m2) {Arbeiter / Fachkräfte};
\node[mit, below=1cm of buch] (m3) {Buchhalter / Controller};
\draw[lin] (chef) -- (verkauf);
\draw[lin] (chef) -- (prod);
\draw[lin] (chef) -- (buch);
\draw[lin] (verkauf) -- (m1);
\draw[lin] (prod) -- (m2);
\draw[lin] (buch) -- (m3);
\end{tikzpicture}
\end{center}
\legende{Organigramme simplifié : la direction (\all{Geschäftsführung}), les départements (\all{Abteilungen}) et le personnel (\all{Mitarbeiter}). Vocabulaire clé : \all{der Vertrieb} (vente), \all{die Buchhaltung} (comptabilité), \all{die Fachkraft} (spécialiste qualifié).}
\end{popfigure}

\courssub{Texte d'application : une entreprise camerounaise}

Le texte suivant réinvestit l'essentiel du vocabulaire vu ci-dessus et anticipe sur les structures grammaticales du chapitre (Passif, Conditionnel, Comparatif). Lisez-le à voix haute, puis répondez aux questions.

\begin{textebox}[Ein Erfolg aus Douala]
Die Firma „Kamerun Kakao GmbH“ wurde 2005 in Douala gegründet. Am Anfang hatte der Betrieb nur zehn Mitarbeiter, heute beschäftigt die Belegschaft über zweihundert Arbeiter und Angestellte. Die Chefin, Frau Mbarga, leitet das Unternehmen mit viel Energie.

Die Firma produziert Kakao und Schokolade von hoher Qualität. Sie exportiert ihre Waren nach Deutschland, in die Schweiz und nach Österreich. Die Nachfrage nach Schokolade aus Kamerun steigt jedes Jahr, und das Angebot der Firma wird immer größer. Auf dem Weltmarkt ist die Konkurrenz stark, aber die Kunden sind mit dem Produkt zufrieden.

Letztes Jahr hat die Firma bei der Bank einen Kredit aufgenommen, um in neue Maschinen zu investieren. Der Gewinn ist gestiegen, und das Unternehmen konnte fünfzig neue Arbeiter einstellen. Die Gehälter der Angestellten und die Löhne der Arbeiter wurden erhöht. Die Arbeitslosigkeit in der Region ist dadurch gesunken.

„Wir bedanken uns bei unseren Kunden und Lieferanten“, sagt Frau Mbarga. „Ohne sie wäre dieser Erfolg nicht möglich.“
\end{textebox}

\textbf{Glossaire.} \all{gründen} (faible) : fonder ; \all{beschäftigen} (faible) : employer ; \all{die Belegschaft} : le personnel ; \all{der Kredit} : le crédit ; \all{aufnehmen} (fort : \all{nimmt auf, nahm auf, hat aufgenommen}) : contracter (un crédit) ; \all{investieren} (faible, sans \all{ge-}) : investir ; \all{einstellen} (séparable) : embaucher ; \all{erhöhen} (faible) : augmenter ; \all{sinken, sank, ist gesunken} (fort, aux. \all{sein}) : baisser ; \all{der Erfolg} : le succès ; \all{möglich} : possible.

\textbf{Questions de compréhension (à répondre en allemand).}
\begin{enumerate}
\item Wann und wo wurde die Firma gegründet?
\item Wohin exportiert die Firma ihre Waren?
\item Warum hat die Firma einen Kredit aufgenommen?
\item Was ist mit den Gehältern und Löhnen passiert? (Passif !)
\item Was sagt Frau Mbarga am Ende? (Conditionnel !)
\end{enumerate}

\begin{exoresolu}[Thème guidé : phrases économiques]
Traduisez en allemand :
1. L'entreprise exporte du cacao vers l'Europe.
2. Le chômage a baissé.
3. Nous vous remercions de votre commande.
4. La banque a accordé un crédit à la firme.
\tcblower
1. \all{Die Firma exportiert Kakao nach Europa.} — Verbe en 2\textsuperscript{e} position ; direction \all{nach} + nom de pays sans article (sauf \all{die Schweiz}, \all{der Iran}...).

2. \all{Die Arbeitslosigkeit ist gesunken.} — Perfekt avec auxiliaire \all{sein} (verbe de changement d'état) ; participe \all{gesunken} (verbe fort \all{sinken, sank, gesunken}).

3. \all{Wir bedanken uns für Ihre Bestellung.} — Verbe réflexif \all{sich bedanken für} + accusatif ; \all{Ihre} avec majuscule (vouvoiement formel).

4. \all{Die Bank hat der Firma einen Kredit gegeben.} — \all{geben} construit avec un datif (\all{der Firma}) pour le bénéficiaire et un accusatif (\all{einen Kredit}) pour l'objet. Variante plus formelle : \all{Die Bank hat der Firma einen Kredit gewährt.} (\all{gewähren}, faible).
\end{exoresolu}

\begin{savaistu}[Wirtschaft : économie ou auberge ?]
Le mot \all{die Wirtschaft} signifie à la fois « l'économie » (d'un pays) et, dans le langage courant en Allemagne du Sud (\all{Bayern}) et en Autriche, « l'auberge », « le petit restaurant » ! Les deux sens ont la même racine : \all{wirtschaften} signifiait « gérer, tenir une maison, faire ménage ». Le terme \all{das Wirtschaftswunder} (« le miracle économique ») désigne la reconstruction spectaculaire de l'Allemagne de l'Ouest après 1945 : en une quinzaine d'années, le pays est devenu l'une des premières puissances industrielles du monde. Le Cameroun entretient d'ailleurs d'anciennes relations économiques avec l'Allemagne, héritées de la période du « Kamerun » allemand (1884--1916) : le port de Douala et les plantations de cacao en sont des témoins historiques.
\end{savaistu}

\begin{exercice}[À vous de jouer — Entraînement Baccalauréat]
\textbf{1. Articles et genres.} Complétez avec l'article défini correct (\all{der}, \all{die}, \all{das}) : \all{\ldots{} Markt, \ldots{} Währung, \ldots{} Konto, \ldots{} Konkurrenz, \ldots{} Gehalt, \ldots{} Ware, \ldots{} Verlust, \ldots{} Zins}.

\textbf{2. Version (Allemand $\to$ Français).} Traduisez :
\begin{enumerate}
\item \all{Der Lieferant hat die Waren zu spät geschickt.}
\item \all{Die Firma stellt zwanzig neue Mitarbeiter ein.}
\item \all{Wir möchten Ihre Produkte bestellen.}
\item \all{Ohne den Kredit wäre die Investition nicht möglich.}
\end{enumerate}

\textbf{3. Thème (Français $\to$ Allemand).} Traduisez :
\begin{enumerate}
\item L'offre est plus grande que la demande. (Comparatif)
\item Nous souhaitons commander vos produits. (Conditionnel de politesse)
\item L'usine a été fondée en 2010. (Passif \all{Präteritum})
\item Le directeur a licencié l'ouvrier. (Verbe fort \all{entlassen})
\end{enumerate}
\end{exercice}

\begin{corrige}[Exercice « À vous de jouer »]
\textbf{1.} \all{der Markt, die Währung, das Konto, die Konkurrenz, das Gehalt, die Ware, der Verlust, der Zins}.

\textbf{2.}
\begin{enumerate}
\item Le fournisseur a envoyé les marchandises trop tard.
\item L'entreprise embauche vingt nouveaux employés.
\item Nous souhaitons commander vos produits.
\item Sans le crédit, l'investissement ne serait pas possible. (Konjunktiv II \all{wäre})
\end{enumerate}

\textbf{3.}
\begin{enumerate}
\item \all{Das Angebot ist größer als die Nachfrage.} (Comparatif \all{größer als})
\item \all{Wir möchten Ihre Produkte bestellen.} (\all{möchten} = Konj. II)
\item \all{Die Fabrik wurde 2010 gegründet.} (Passif \all{Präteritum} : \all{wurde} + Partizip II)
\item \all{Der Chef hat den Arbeiter entlassen.} (Perfekt verbe fort : \all{hat entlassen})
\end{enumerate}
\end{corrige}

\begin{aretenir}[L'essentiel du lexique économique — Fiche de révision]
\begin{itemize}
\item \textbf{Cinq champs lexicaux} à maîtriser pour l'écrit et l'oral :
  \begin{itemize}
  \item Entreprise : \all{die Firma, das Unternehmen, der Betrieb, gründen, leiten, die Belegschaft}.
  \item Commerce : \all{der Handel, exportieren, importieren, der Markt, die Konkurrenz, der Kunde, der Lieferant}.
  \item Argent/Banque : \all{das Konto, der Kredit, der Gewinn, der Verlust, die Bank, investieren, die Rechnung}.
  \item Emploi : \all{die Arbeit, der Arbeitsplatz, die Arbeitslosigkeit, der Arbeiter, der Angestellte, das Gehalt, der Lohn, einstellen, entlassen}.
  \item Production : \all{die Ware, das Produkt, die Qualität, die Nachfrage, das Angebot, produzieren, verbrauchen}.
  \end{itemize}
\item \textbf{Règles d'or} : Tout nom $\to$ Majuscule + Article + Pluriel. Verbes \all{-ieren} $\to$ Pas de \all{ge-} au Partizip II. \all{sich bedanken für} + Akk. \all{Mit freundlichen Grüßen} sans point.
\item \textbf{Faux amis} : \all{Fabrik} = usine ; \all{bekommen} = recevoir ; \all{aktuell} = actuel (adj.) ; \all{konkurrent} = concurrent.
\item \textbf{Ponts vers la grammaire} : \all{größer als} (Comparatif), \all{wäre/möchten} (Conditionnel/Konj. II), \all{wurde gegründet/wurden erhöht} (Passif).
\end{itemize}
\end{aretenir}

\coursec{Grammaire 1 : le comparatif et le superlatif}

Dans la section précédente, nous avons constitué un lexique riche de la vie économique : \all{der Markt}, \all{die Produktion}, \all{der Preis}, \all{die Qualität}, \all{der Gewinn}... Or, dès qu'on parle d'économie, on compare : un prix est \emph{plus haut} qu'un autre, une qualité est \emph{meilleure}, une entreprise est \emph{la plus grande} du pays. Comparer, c'est précisément le rôle grammatical du \cle{comparatif} et du \cle{superlatif}. Cette section les reprend de manière systématique : formation, emploi, irrégularités, déclinaison et pièges de traduction.

\courssub{Observation : partir des exemples}

Lisons d'abord ces phrases, extraites d'un article économique imaginaire sur le commerce entre l'Allemagne et le Cameroun :

\begin{exemplebox}[Observer avant de comprendre]
\begin{enumerate}
\item \all{Der deutsche Markt ist größer als der kamerunische Markt.} \\
« Le marché allemand est plus grand que le marché camerounais. »
\item \all{Kaffee ist teurer als Kakao.} \\
« Le café est plus cher que le cacao. »
\item \all{Die beste Qualität kommt aus dem Westen.} \\
« La meilleure qualité vient de l'Ouest. »
\item \all{Die Produktion ist so hoch wie letztes Jahr.} \\
« La production est aussi élevée que l'année dernière. »
\end{enumerate}
\end{exemplebox}

Observons : dans les phrases 1 et 2, l'adjectif (\all{groß}, \all{teuer}) prend la terminaison \all{-er} et le second terme de la comparaison est introduit par \all{als} (« que »). Dans la phrase 3, l'adjectif prend la terminaison \all{-ste} et l'article défini : c'est le superlatif. Dans la phrase 4, on retrouve la structure \all{so ... wie} (« aussi ... que »). Trois constructions, trois degrés de comparaison : c'est toute la logique du système.

\courssub{Le comparatif de supériorité : adjectif + -er + als}

\begin{propriete}[Formation du comparatif]
Le \cle{comparatif} se forme en ajoutant la terminaison \all{-er} à l'adjectif. Le second terme de la comparaison est introduit par \all{als} (= « que ») :
\begin{center}
\all{adjectif + -er + als}
\end{center}
\all{klein} (petit) $\rightarrow$ \all{kleiner als} (plus petit que) ; \all{billig} (bon marché) $\rightarrow$ \all{billiger als} (moins cher que).
\end{propriete}

\begin{exemplebox}[Comparatifs en contexte économique]
\begin{itemize}
\item \all{Die Preise auf dem Markt in Douala sind höher als im Dorf.} \\
« Les prix au marché de Douala sont plus élevés qu'au village. »
\item \all{Ein Bendskin ist schneller als ein Bus im Stau.} \\
« Un bendskin est plus rapide qu'un bus dans les embouteillages. »
\item \all{Die neue Fabrik ist moderner als die alte.} \\
« La nouvelle usine est plus moderne que l'ancienne. »
\item \all{Importierte Waren sind oft teurer als lokale Produkte.} \\
« Les marchandises importées sont souvent plus chères que les produits locaux. »
\end{itemize}
\end{exemplebox}

\begin{attention}[Deux petites modifications de forme]
\begin{itemize}
\item Les adjectifs en \all{-el} perdent le \all{-e-} : \all{dunkel} $\rightarrow$ \all{dunkler} (et non \emph{dunkeler}).
\item Les adjectifs en \all{-er} comme \all{teuer} perdent aussi le \all{-e-} : \all{teuer} $\rightarrow$ \all{teurer}.
\item Les adjectifs en \all{-d}, \all{-t}, \all{-s}, \all{-ß}, \all{-z}, \all{-sch} prennent souvent \all{-er} sans difficulté, mais au superlatif ils prendront \all{-esten} (voir plus bas) : \all{heiß} $\rightarrow$ \all{heißer}.
\end{itemize}
\end{attention}

\courssub{Le superlatif : am ... -sten et der/die/das ... -ste}

Le superlatif exprime le degré le plus haut (« le plus »). L'allemand possède \textbf{deux formes} selon la place de l'adjectif dans la phrase.

\begin{propriete}[Formation du superlatif]
\begin{enumerate}
\item \textbf{Forme prédicative (attribut du sujet) / adverbiale} (après \all{sein}, \all{werden}, \all{bleiben} ou en fonction d'adverbe) : \all{am + adjectif + -sten}. \\
\all{Der Mount Cameroon ist am höchsten.} « Le mont Cameroun est le plus haut. »
\item \textbf{Forme attributive} (adjectif placé devant le nom) : \all{der/die/das + adjectif + -ste}, avec déclinaison adjectivale normale. \\
\all{die beste Qualität} « la meilleure qualité » ; \all{der größte Markt} « le plus grand marché ».
\end{enumerate}
\textbf{Exception de forme} : les adjectifs terminés par \all{-d}, \all{-t}, \all{-s}, \all{-ß}, \all{-z}, \all{-sch} prennent \all{-esten} et non \all{-sten} : \all{heiß} $\rightarrow$ \all{am heißesten} ; \all{alt} $\rightarrow$ \all{am ältesten} ; \all{frisch} $\rightarrow$ \all{am frischesten}.
\end{propriete}

\begin{exemplebox}[Les deux formes du superlatif]
\begin{itemize}
\item \all{Im Juli sind die Preise am niedrigsten.} \\
« En juillet, les prix sont les plus bas. » (forme avec \all{am})
\item \all{Die größte Bank des Landes hat ihren Sitz in Yaoundé.} \\
« La plus grande banque du pays a son siège à Yaoundé. » (forme attributive)
\item \all{Welches Produkt ist am billigsten?} \\
« Quel produit est le moins cher ? »
\item \all{Das teuerste Geschäft liegt im Stadtzentrum.} \\
« Le magasin le plus cher se trouve au centre-ville. »
\end{itemize}
\end{exemplebox}

\courssub{L'umlaut au comparatif et au superlatif}

\begin{propriete}[Umlaut sur a, o, u]
La plupart des adjectifs \textbf{monosyllabiques} dont la voyelle est \textbf{a}, \textbf{o} ou \textbf{u} prennent l'\cle{umlaut} au comparatif et au superlatif :
\begin{center}
\all{alt} $\rightarrow$ \all{älter} $\rightarrow$ \all{am ältesten} \\
\all{stark} $\rightarrow$ \all{stärker} $\rightarrow$ \all{am stärksten} \\
\all{hoch} $\rightarrow$ \all{höher} $\rightarrow$ \all{am höchsten} \\
\all{groß} $\rightarrow$ \all{größer} $\rightarrow$ \all{am größten} \\
\all{jung} $\rightarrow$ \all{jünger} $\rightarrow$ \all{am jüngsten} \\
\all{kurz} $\rightarrow$ \all{kürzer} $\rightarrow$ \all{am kürzesten}
\end{center}
Les adjectifs de plusieurs syllabes (\all{modern}, \all{interessant}, \all{billig}) ne prennent \textbf{jamais} l'umlaut : \all{moderner}, \all{interessanter}, \all{billiger}. Quelques monosyllabes échappent aussi à l'umlaut : \all{klar} $\rightarrow$ \all{klarer} ; \all{voll} $\rightarrow$ \all{voller}.
\end{propriete}

\begin{attention}[Attention à \all{hoch}]
L'adjectif \all{hoch} perd son \all{-c-} dès qu'il reçoit une terminaison : \all{höher}, \all{am höchsten}, \all{das höchste Gebäude}. On n'écrit jamais \emph{hocher}.
\end{attention}

\courssub{Les comparatifs irréguliers : à apprendre par cœur}

Quatre adjectifs/adverbes ont des formes totalement irrégulières, exactement comme en français (« bon / meilleur / le meilleur ») :

\begin{center}
\begin{tabularx}{\linewidth}{|X|X|X|X|}
\hline
\rowcolor{popblueL} Positif & Comparatif & Superlatif & Français \\
\hline
\all{gut} & \all{besser} & \all{am besten} & bon / meilleur / le meilleur \\
\hline
\all{viel} & \all{mehr} & \all{am meisten} & beaucoup / plus / le plus \\
\hline
\all{gern} & \all{lieber} & \all{am liebsten} & volontiers / plutôt / de préférence \\
\hline
\all{nah} & \all{näher} & \all{am nächsten} & proche / plus proche / le plus proche \\
\hline
\end{tabularx}
\end{center}

\begin{exemplebox}[Irréguliers en phrases]
\begin{itemize}
\item \all{Die Qualität dieses Kakaos ist besser.} \\
« La qualité de ce cacao est meilleure. »
\item \all{Diese Firma verdient mehr Geld als ihre Konkurrenten.} \\
« Cette entreprise gagne plus d'argent que ses concurrents. »
\item \all{Ich kaufe am liebsten auf dem Markt ein.} \\
« J'aime faire mes courses de préférence au marché. »
\item \all{Der nächste Geldautomat ist zwei Kilometer entfernt.} \\
« Le distributeur le plus proche est à deux kilomètres. »
\end{itemize}
\end{exemplebox}

\courssub{Le comparatif d'égalité : so ... wie}

\begin{propriete}[Comparaison d'égalité]
Pour dire « aussi ... que », l'allemand utilise la structure encadrante \all{so + adjectif (au positif) + wie} :
\begin{center}
\all{so ... wie} \hspace{1cm} négation : \all{nicht so ... wie}
\end{center}
L'adjectif reste au \textbf{positif} : pas de terminaison \all{-er} à l'intérieur de cette structure.
\end{propriete}

\begin{exemplebox}[Égalité et inégalité]
\begin{itemize}
\item \all{Die Produktion ist so hoch wie letztes Jahr.} \\
« La production est aussi élevée que l'année dernière. »
\item \all{Er verdient nicht so viel wie sein Kollege.} \\
« Il ne gagne pas autant que son collègue. »
\item \all{Der Markt in Bafoussam ist so groß wie der in Bamenda.} \\
« Le marché de Bafoussam est aussi grand que celui de Bamenda. »
\item \all{Unsere Preise sind nicht so niedrig wie früher.} \\
« Nos prix ne sont pas aussi bas qu'avant. »
\end{itemize}
\end{exemplebox}

\courssub{L'adjectif comparé en fonction d'épithète : la déclinaison}

Quand le comparatif ou le superlatif est placé \textbf{devant un nom}, il se décline comme n'importe quel adjectif épithète, selon l'article, le genre, le nombre et le cas.

\begin{exemplebox}[Déclinaison des adjectifs comparés]
\begin{itemize}
\item \all{ein größerer Markt} « un marché plus grand » (nominatif masculin, article indéfini)
\item \all{die bessere Qualität} « la meilleure qualité » (nominatif féminin, article défini)
\item \all{mit einem billigeren Produkt} « avec un produit moins cher » (datif)
\item \all{die jüngsten Arbeiter} « les ouvriers les plus jeunes » (nominatif pluriel)
\item \all{Wir suchen einen erfahreneren Mitarbeiter.} \\
« Nous cherchons un employé plus expérimenté. » (accusatif masculin)
\end{itemize}
\end{exemplebox}

\begin{aretenir}[La logique à retenir]
La terminaison de comparaison (\all{-er} / \all{-ste}) vient \textbf{d'abord}, la terminaison de déclinaison vient \textbf{ensuite} : \all{größ-} + \all{-er} (comparatif) + \all{-er} (déclinaison) = \all{größerer}. Ne jamais oublier la deuxième terminaison devant un nom.
\end{aretenir}

\courssub{Tableau récapitulatif des trois degrés}

\begin{center}
\begin{tabular}{|l|l|l|l|}
\hline
\rowcolor{popblueL} Positif & Comparatif & Superlatif (\all{am}) & Particularité \\
\hline
\all{klein} & \all{kleiner} & \all{am kleinsten} & régulier \\
\hline
\all{alt} & \all{älter} & \all{am ältesten} & umlaut \\
\hline
\all{hoch} & \all{höher} & \all{am höchsten} & umlaut, perte du -c- \\
\hline
\all{groß} & \all{größer} & \all{am größten} & umlaut \\
\hline
\all{gut} & \all{besser} & \all{am besten} & irrégulier \\
\hline
\all{viel} & \all{mehr} & \all{am meisten} & irrégulier \\
\hline
\all{gern} & \all{lieber} & \all{am liebsten} & irrégulier \\
\hline
\all{teuer} & \all{teurer} & \all{am teuersten} & perte du -e- \\
\hline
\end{tabular}
\end{center}

\begin{popfigure}
\begin{center}
\begin{tikzpicture}[scale=1, transform shape]
\draw[-{Stealth[length=3mm]}, very thick, popmuted] (0,0) -- (10.5,0) node[below left=0.1cm and 3.2cm, popmuted] {\footnotesize degré croissant};
\node[draw=popblue, fill=popblueL, rounded corners, align=center, minimum width=2.6cm, minimum height=1.3cm] (pos) at (1.6,1.2) {\textbf{Positif}\\ \all{klein}\\ \footnotesize « petit »};
\node[draw=poporange, fill=poporangeL, rounded corners, align=center, minimum width=2.6cm, minimum height=1.3cm] (comp) at (5.2,1.2) {\textbf{Comparatif}\\ \all{kleiner als}\\ \footnotesize « plus petit que »};
\node[draw=popgreen, fill=popgreenL, rounded corners, align=center, minimum width=2.6cm, minimum height=1.3cm] (sup) at (8.8,1.2) {\textbf{Superlatif}\\ \all{am kleinsten}\\ \footnotesize « le plus petit »};
\draw[-{Stealth[length=2.5mm]}, thick, poporange] (pos) -- (comp) node[midway, above, popdark] {\footnotesize + \all{-er}};
\draw[-{Stealth[length=2.5mm]}, thick, popgreen] (comp) -- (sup) node[midway, above, popdark] {\footnotesize \all{am} + \all{-sten}};
\node[align=center, popdark] at (5.2,-0.9) {\footnotesize Égalité : \all{so klein wie} « aussi petit que »};
\end{tikzpicture}
\end{center}
\legende{L'échelle des trois degrés de comparaison : positif, comparatif, superlatif.}
\end{popfigure}

\courssub{Comparaison des deux langues : français et allemand}

\begin{center}
\begin{tabularx}{\linewidth}{|X|X|X|}
\hline
\rowcolor{popblueL} Français & Deutsch & Remarque \\
\hline
plus grand que & \all{größer als} & \all{als} = « que », jamais \all{wie} \\
\hline
aussi grand que & \all{so groß wie} & ici c'est \all{wie}, jamais \all{als} \\
\hline
le plus grand (épithète) & \all{der größte} & article + \all{-ste} + déclinaison \\
\hline
le plus grand (attribué) & \all{am größten} & \all{am} + \all{-sten} \\
\hline
meilleur & \all{besser} & irrégulier, jamais \all{mehr gut} \\
\hline
de plus en plus cher & \all{immer teurer} & \all{immer} + comparatif \\
\hline
\end{tabularx}
\end{center}

\begin{exoresolu}[Mettez les adjectifs au degré demandé]
Énoncé : complétez avec la forme correcte de l'adjectif entre parenthèses.
\begin{enumerate}
\item \all{Der Hafen in Douala ist \dots\ als der Hafen in Kribi.} (\all{groß})
\item \all{Der Mount Cameroon ist der \dots\ Berg des Landes.} (\all{hoch})
\item \all{Dieses Produkt ist \dots\ als das importierte.} (\all{billig})
\item \all{Unsere Firma hat die \dots\ Qualität.} (\all{gut})
\item \all{Der Zug ist nicht \dots\ schnell \dots\ das Flugzeug.} (\all{so ... wie})
\end{enumerate}
\tcblower
Solution détaillée :
\begin{enumerate}
\item \all{größer} : comparatif de \all{groß} avec umlaut (\all{o} $\rightarrow$ \all{ö}) + \all{-er} ; le second terme est introduit par \all{als}. « Le port de Douala est plus grand que le port de Kribi. »
\item \all{höchste} : superlatif attributif devant le nom masculin \all{Berg} avec article défini : \all{der höchste Berg}. Perte du \all{-c-} de \all{hoch}. « Le mont Cameroun est la plus haute montagne du pays. »
\item \all{billiger} : comparatif régulier, sans umlaut (adjectif de deux syllabes). « Ce produit est moins cher que le produit importé. »
\item \all{beste} : superlatif irrégulier de \all{gut} (\all{gut} $\rightarrow$ \all{besser} $\rightarrow$ \all{best-}), avec déclinaison féminine : \all{die beste Qualität}. « Notre entreprise a la meilleure qualité. »
\item \all{so} ... \all{wie} : comparaison d'égalité niée, l'adjectif reste au positif. « Le train n'est pas aussi rapide que l'avion. »
\end{enumerate}
\end{exoresolu}

\courssub{Pièges fréquents pour les francophones}

\begin{attention}[Les quatre erreurs à bannir]
\begin{enumerate}
\item \textbf{Jamais} \all{mehr besser} ni \all{mehr größer} ! En français on dit « plus grand », mais l'allemand n'ajoute jamais \all{mehr} devant un comparatif : on dit simplement \all{größer}, \all{besser}, \all{billiger}. \all{Mehr} est lui-même le comparatif de \all{viel} et ne s'emploie que seul : \all{mehr Geld} « plus d'argent ».
\item \all{als} ≠ \all{wenn} : \all{als} sert à la comparaison d'inégalité (\all{größer als}) ; \all{wenn} signifie « quand / si ». Ne les confondez pas.
\item Comparaison d'inégalité = \all{als} ; comparaison d'égalité = \all{wie}. Dire \all{größer wie} est une erreur grave et fréquente.
\item Devant un nom, n'oubliez pas la déclinaison : \all{ein größerer Markt}, pas \all{ein größer Markt}.
\end{enumerate}
\end{attention}

\begin{exercice}[À vous de comparer]
Traduisez en allemand :
\begin{enumerate}
\item Le cacao camerounais est meilleur que le cacao importé.
\item Yaoundé n'est pas aussi grande que Douala.
\item Quelle est la banque la plus moderne du pays ?
\item Les prix sont de plus en plus hauts.
\item Ce magasin est le plus proche, mais ses produits sont plus chers.
\end{enumerate}
\end{exercice}

\begin{corrige}[Exercice « À vous de comparer »]
\begin{enumerate}
\item \all{Der kamerunische Kakao ist besser als der importierte Kakao.} — comparatif irrégulier de \all{gut} : \all{besser}.
\item \all{Yaoundé ist nicht so groß wie Douala.} — égalité niée : \all{nicht so ... wie}, adjectif au positif.
\item \all{Welche ist die modernste Bank des Landes?} — superlatif attributif : \all{die modernste Bank} (pas d'umlaut, adjectif de deux syllabes).
\item \all{Die Preise sind immer höher.} — \all{immer} + comparatif avec umlaut et perte du \all{-c-} de \all{hoch}.
\item \all{Dieses Geschäft ist am nächsten, aber seine Produkte sind teurer.} — superlatif irrégulier de \all{nah} : \all{am nächsten} ; comparatif régulier \all{teurer}.
\end{enumerate}
\end{corrige}

\begin{aretenir}[L'essentiel du comparatif et du superlatif]
\begin{itemize}
\item Comparatif : \all{adjectif + -er + als} (\all{größer als} « plus grand que »).
\item Superlatif : \all{am + adjectif + -sten} ou \all{der/die/das + adjectif + -ste}.
\item Umlaut sur \all{a}, \all{o}, \all{u} pour la plupart des monosyllabes : \all{alt} $\rightarrow$ \all{älter}, \all{hoch} $\rightarrow$ \all{höher}.
\item Irréguliers : \all{gut} $\rightarrow$ \all{besser} $\rightarrow$ \all{am besten} ; \all{viel} $\rightarrow$ \all{mehr} $\rightarrow$ \all{am meisten} ; \all{gern} $\rightarrow$ \all{lieber} $\rightarrow$ \all{am liebsten}.
\item Égalité : \all{so ... wie} ; négation : \all{nicht so ... wie}.
\item Devant un nom, l'adjectif comparé se décline : \all{die bessere Qualität}.
\end{itemize}
\end{aretenir}

Ces outils de comparaison seront immédiatement réinvestis dans la section suivante, consacrée au conditionnel (\all{Konjunktiv II}), qui permet de formuler des souhaits et des hypothèses économiques : « si les prix étaient plus bas, nous achèterions davantage... ».

\coursec{Grammaire 2 : le conditionnel (Konjunktiv II)}

Dans la section précédente, nous avons appris à comparer : \all{billiger als}, \all{am billigsten}. Comparer, c'est décrire la réalité économique telle qu'elle est. Mais la vie économique, c'est aussi le monde des projets, des négociations et des rêves : « Si j'avais plus d'argent, j'ouvrirais un magasin », « Pourriez-vous baisser le prix ? », « À votre place, j'accepterais l'offre ». Toutes ces phrases ont un point commun : elles expriment l'\emph{irréel}, le souhait, l'hypothèse ou la politesse. En allemand, c'est le domaine du \cle{Konjunktiv II}, que nous appelons en français « le conditionnel ».

\courssub{Observation : le conditionnel dans un dialogue commercial}

Lisons d'abord ce court dialogue entre une cliente camerounaise et un vendeur à Douala. Soulignez mentalement toutes les formes qui expriment un souhait, une hypothèse ou une formule de politesse.

\begin{textebox}[Am Markt in Douala — Au marché de Douala]
\textbf{Kundin:} Guten Tag! Ich \textbf{hätte} gern zehn Kilo Reis. Wie viel \textbf{kostet} das?\par
\textbf{Verkäufer:} Das macht 8000 CFA. Wenn Sie zwanzig Kilo \textbf{nähmen}, \textbf{könnte} ich Ihnen einen besseren Preis machen.\par
\textbf{Kundin:} Hm, das \textbf{wäre} zu viel für mich. Wenn ich mehr Geld \textbf{hätte}, \textbf{würde} ich sofort zwanzig Kilo kaufen. \textbf{Könnten} Sie mir vielleicht trotzdem einen Rabatt geben?\par
\textbf{Verkäufer:} An Ihrer Stelle \textbf{würde} ich das Angebot annehmen. Aber gut: Ich \textbf{würde} Ihnen 7500 CFA vorschlagen.\par
\textbf{Kundin:} Einverstanden! Wenn mein Geschäft besser \textbf{laufen würde}, \textbf{käme} ich jede Woche zu Ihnen.\par
\textbf{Verkäufer:} Das \textbf{wäre} schön! Hier ist Ihr Reis. Auf Wiedersehen!
\end{textebox}

\textbf{Glossaire :} der Reis (le riz, pas de pluriel usuel / \emph{die Reissorten} pour les variétés) ; der Preis, -e (le prix) ; der Rabatt, -e (la remise) ; das Angebot, -e (l'offre) ; das Geschäft, -e (le commerce, l'affaire) ; einverstanden (d'accord).

\textbf{Questions de compréhension :}
\begin{enumerate}
\item Que veut acheter la cliente, et en quelle quantité ?
\item Quelle proposition le vendeur fait-il pour obtenir un meilleur prix ?
\item Pourquoi la cliente n'achète-elle pas vingt kilos ?
\item Quel prix final le vendeur propose-t-il ?
\end{enumerate}

\textbf{Observation grammaticale :} dans ce dialogue, on trouve \all{ich hätte gern} (je voudrais), \all{könnten Sie...?} (pourriez-vous...?), \all{das wäre zu viel} (ce serait trop), \all{wenn ich mehr Geld hätte, würde ich kaufen} (si j'avais plus d'argent, j'achèterais). Toutes ces formes appartiennent au \cle{Konjunktiv II}. Remarquez deux choses : d'une part la construction \all{würde + infinitif} ; d'autre part des formes spéciales pour \all{haben} (hätte), \all{sein} (wäre) et \all{können} (könnte).

\courssub{Les emplois du Konjunktiv II}

Le Konjunktiv II sert à sortir de la réalité. Voici ses cinq emplois principaux, tous illustrés par le dialogue :

\begin{enumerate}
\item \textbf{Le souhait irréel} : \all{Wenn ich doch mehr Geld hätte!} « Si seulement j'avais plus d'argent ! » (mais je n'en ai pas).
\item \textbf{L'hypothèse / la condition irréelle} : \all{Wenn ich mehr Geld hätte, würde ich zwanzig Kilo kaufen.} « Si j'avais plus d'argent, j'achèterais vingt kilos. »
\item \textbf{Le conseil} : \all{An Ihrer Stelle würde ich das Angebot annehmen.} « À votre place, j'accepterais l'offre. »
\item \textbf{La politesse (demande atténuée)} : \all{Könnten Sie mir bitte einen Rabatt geben?} « Pourriez-vous me faire une remise, s'il vous plaît ? »
\item \textbf{Le rêve, l'imagination} : \all{Ich wäre gern reich.} « J'aimerais être riche. »
\end{enumerate}

\begin{popfigure}
\begin{center}
\begin{tikzpicture}[
  central/.style={rectangle, rounded corners, draw=popblue, fill=popblueL, thick, align=center, font=\bfseries, minimum width=3.6cm, minimum height=1cm},
  branch/.style={rectangle, rounded corners, draw=popdark, fill=popcream, align=center, minimum width=4.6cm, minimum height=0.9cm, font=\small},
  arr/.style={-{Stealth[length=3mm]}, thick, popdark}
]
\node[central, align=center] (c){Konjunktiv II\\ le conditionnel};
\node[branch, above left=1.1cm and 0.4cm of c, align=center] (s){Souhait irréel\\ \all{Wenn ich doch reich wäre!}};
\node[branch, below left=1.1cm and 0.4cm of c, align=center] (h){Hypothèse / condition\\ \all{Wenn ich Zeit hätte, ...}};
\node[branch, above right=1.1cm and 0.4cm of c, align=center] (co){Conseil\\ \all{An deiner Stelle würde ich...}};
\node[branch, below right=1.1cm and 0.4cm of c, align=center] (p){Politesse\\ \all{Ich hätte gern ...}};
\draw[arr] (c) -- (s);
\draw[arr] (c) -- (h);
\draw[arr] (c) -- (co);
\draw[arr] (c) -- (p);
\end{tikzpicture}
\end{center}
\legende{Carte mentale : les quatre grands emplois du Konjunktiv II.}
\end{popfigure}

\courssub{La formation : würde + infinitif}

\begin{propriete}[Formation du Konjunktiv II présent]
Pour la grande majorité des verbes, le Konjunktiv II présent se forme avec l'auxiliaire \cle{würde} (conjugué) + l'\cle{infinitif} placé \textbf{en fin de phrase} :
\begin{center}
sujet + \textbf{würde} (2\up{e} position) + ... + \textbf{infinitif} (dernière position)
\end{center}
\all{Ich \textbf{würde} ein Geschäft in Yaoundé \textbf{eröffnen}.} « J'ouvrirais un magasin à Yaoundé. »
\end{propriete}

Le verbe \all{würden} se conjugue ainsi (c'est le Präteritum de \all{werden} avec un Umlaut) :

\begin{center}
\begin{tabular}{|l|l|l|}
\hline
\rowcolor{popblueL} Personne & würden & Français \\
\hline
ich & würde & je ... -rais \\
\hline
du & würdest & tu ... -rais \\
\hline
er/sie/es & würde & il/elle ... -rait \\
\hline
wir & würden & nous ... -rions \\
\hline
ihr & würdet & vous ... -riez \\
\hline
sie/Sie & würden & ils/vous ... -raient \\
\hline
\end{tabular}
\end{center}

\begin{exemplebox}[würde + infinitif en contexte économique]
\all{Wir \textbf{würden} mehr Kakao \textbf{exportieren}, wenn der Preis besser wäre.} « Nous exporterions davantage de cacao si le prix était meilleur. »\par
\all{\textbf{Würdest} du in Deutschland \textbf{arbeiten}?} « Travaillerais-tu en Allemagne ? »\par
\all{Die Firma \textbf{würde} neue Arbeiter \textbf{einstellen}.} « L'entreprise embaucherait de nouveaux ouvriers. »
\end{exemplebox}

\courssub{Les formes fortes à connaître par cœur}

Pour un petit groupe de verbes très fréquents, on n'utilise \textbf{pas} \all{würde}, mais une forme dite « forte », dérivée du Präteritum avec un Umlaut (et parfois une modification de la voyelle). Ce sont les verbes les plus utilisés de la langue : \all{sein}, \all{haben}, les verbes modaux, \all{wissen}, \all{geben}, \all{werden}, \all{kommen}.

\begin{center}
\begin{tabular}{|l|l|l|l|}
\hline
\rowcolor{popblueL} Infinitif & Präteritum & Konjunktiv II & Français \\
\hline
sein & war & \textbf{wäre} & je serais \\
\hline
haben & hatte & \textbf{hätte} & j'aurais \\
\hline
können & konnte & \textbf{könnte} & je pourrais \\
\hline
müssen & musste & \textbf{müsste} & je devrais \\
\hline
dürfen & durfte & \textbf{dürfte} & il me serait permis \\
\hline
sollen & sollte & \textbf{sollte} & je devrais (conseil) \\
\hline
wollen & wollte & \textbf{wollte} & je voudrais (rare, voir \all{möchte}) \\
\hline
wissen & wusste & \textbf{wüsste} & je saurais \\
\hline
geben & gab & \textbf{es gäbe} & il y aurait \\
\hline
werden & wurde & \textbf{würde} & je deviendrais / auxiliaire \\
\hline
kommen & kam & \textbf{käme} & je viendrais \\
\hline
\end{tabular}
\end{center}

\begin{propriete}[Règle d'or]
Konjunktiv II présent = \all{würde} + infinitif ; \textbf{mais} pour \all{sein}, \all{haben} et les verbes modaux, on utilise toujours les formes fortes : \all{wäre}, \all{hätte}, \all{könnte}, \all{müsste}, \all{sollte}, \all{dürfte}, \all{wollte}...
\end{propriete}

\begin{exemplebox}[Les formes fortes en phrases]
\all{Wenn die Bauern mehr Geld \textbf{hätten}, \textbf{würden} sie Maschinen \textbf{kaufen}.} « Si les paysans avaient plus d'argent, ils achèteraient des machines. »\par
\all{Das \textbf{wäre} eine gute Investition.} « Ce serait un bon investissement. »\par
\all{\textbf{Könnten} Sie uns den Preis senken?} « Pourriez-vous baisser le prix pour nous ? »\par
\all{Du \textbf{solltest} mehr sparen.} « Tu devrais économiser davantage. » (conseil)\par
\all{Es \textbf{gäbe} weniger Arbeitslosigkeit, wenn es mehr Fabriken gäbe.} « Il y aurait moins de chômage s'il y avait plus d'usines. »
\end{exemplebox}

\courssub{La phrase conditionnelle avec wenn}

\begin{propriete}[Structure de la phrase conditionnelle]
La subordonnée introduite par \all{wenn} (si) rejette le \textbf{verbe conjugué en dernière position}. Si la subordonnée est placée en tête, le verbe de la principale vient immédiatement après la virgule (position 1) :
\begin{center}
\all{Wenn ich Zeit \textbf{hätte}, \textbf{würde} ich mehr arbeiten.}
\end{center}
« Si j'avais le temps, je travaillerais davantage. »
\end{propriete}

\begin{exemplebox}[Phrases conditionnelles]
\all{Wenn ich reich \textbf{wäre}, \textbf{würde} ich eine Schule \textbf{bauen}.} « Si j'étais riche, je construirais une école. »\par
\all{Wenn der Wechselkurs besser \textbf{wäre}, \textbf{könnten} wir mehr \textbf{importieren}.} « Si le taux de change était meilleur, nous pourrions importer davantage. »\par
\all{Ich \textbf{würde} das Produkt \textbf{kaufen}, wenn es billiger \textbf{wäre}.} « J'achèterais le produit s'il était moins cher. »
\end{exemplebox}

\courssub{Le conditionnel passé}

Pour parler d'une hypothèse \emph{dans le passé} (ce qui ne s'est pas produit), on utilise :

\begin{propriete}[Konjunktiv II passé]
\cle{hätte} ou \cle{wäre} (conjugué) + \textbf{participe II} en fin de phrase. On choisit \all{hätte} ou \all{wäre} selon la même règle qu'au Perfekt : \all{wäre} pour les verbes de mouvement/changement d'état, \all{hätte} pour tous les autres.
\begin{center}
\all{Ich \textbf{hätte} das Angebot \textbf{angenommen}.} « J'aurais accepté l'offre. »\par
\all{Wir \textbf{wären} nach Deutschland \textbf{gereist}.} « Nous serions partis en Allemagne. »
\end{center}
\end{propriete}

\begin{exemplebox}[Conditionnel passé]
\all{Wenn ich Geld \textbf{gehabt hätte}, \textbf{hätte} ich das Haus \textbf{gekauft}.} « Si j'avais eu de l'argent, j'aurais acheté la maison. » (mais je n'ai pas eu d'argent, donc je n'ai pas acheté)\par
\all{Ohne den Regen \textbf{wäre} die Ernte besser \textbf{gewesen}.} « Sans la pluie, la récolte aurait été meilleure. »
\end{exemplebox}

\courssub{La politesse : les formules à retenir}

Dans la vie économique (magasin, banque, lettre commerciale), le Konjunktiv II adoucit les demandes. Ces formules sont indispensables :

\begin{center}
\begin{tabularx}{\linewidth}{|X|X|}
\hline
\rowcolor{popblueL} Deutsch & Français \\
\hline
\all{Ich hätte gern einen Termin.} & Je voudrais un rendez-vous. \\
\hline
\all{Könnten Sie mir bitte helfen?} & Pourriez-vous m'aider, s'il vous plaît ? \\
\hline
\all{Ich würde vorschlagen, dass wir uns morgen treffen.} & Je proposerais que nous nous rencontrions demain. \\
\hline
\all{Dürfte ich Sie etwas fragen?} & Puis-je vous demander quelque chose ? \\
\hline
\all{Würden Sie bitte das Fenster öffnen?} & Voudriez-vous ouvrir la fenêtre, s'il vous plaît ? \\
\hline
\all{Sie sollten den Vertrag lesen.} & Vous devriez lire le contrat. \\
\hline
\end{tabularx}
\end{center}

\courssub{Comparaison français / allemand}

\begin{center}
\begin{tabularx}{\linewidth}{|X|X|X|}
\hline
\rowcolor{popblueL} Français & Deutsch & Remarque \\
\hline
j'achèterais & \all{ich würde kaufen} & würde + infinitif en fin \\
\hline
je serais & \all{ich wäre} & forme forte, jamais « würde sein » \\
\hline
j'aurais & \all{ich hätte} & forme forte, jamais « würde haben » \\
\hline
je pourrais & \all{ich könnte} & forme forte du modal \\
\hline
je voudrais & \all{ich hätte gern / ich möchte} & politesse au magasin / souhait \\
\hline
j'aurais acheté & \all{ich hätte gekauft} & conditionnel passé \\
\hline
\end{tabularx}
\end{center}

\begin{attention}[Erreurs fréquentes des francophones]
\begin{enumerate}
\item Ne dites jamais \all{ich würde sein} ni \all{ich würde haben} : on dit \all{ich \textbf{wäre}} et \all{ich \textbf{hätte}}. Même chose pour les modaux : \all{ich \textbf{könnte}}, pas « ich würde können ».
\item L'infinitif avec \all{würde} va \textbf{en toute fin de phrase} : \all{Ich würde morgen nach Douala \textbf{fahren}} (et non « Ich würde fahren morgen nach Douala »).
\item Dans la subordonnée avec \all{wenn}, le verbe conjugué va \textbf{en fin} : \all{Wenn ich reich \textbf{wäre}, ...} (et non « Wenn ich wäre reich »).
\item Attention au faux ami : \all{bekommen} signifie « recevoir », pas « devenir » ; \all{ich würde bekommen} = « je recevrais ».
\item Le français « je voudrais » ne se traduit pas par un verbe \all{wollen} au conditionnel, mais par \all{ich hätte gern} (objet) ou \all{ich möchte} (verbe d'action).
\item Verbes à particule : \all{annehmen} (accepter) est \textbf{séparable} (\all{er nimmt an}). À l'infinitif, il s'écrit en un mot : \all{annehmen}. Participe II : \all{angenommen}.
\end{enumerate}
\end{attention}

\begin{exoresolu}[Thème : traduisez en allemand]
Énoncé : traduisez les phrases suivantes en allemand. 1) Si j'avais le temps, je visiterais l'usine. 2) À votre place, j'accepterais le contrat. 3) Je voudrais un entretien avec le directeur. 4) Pourriez-vous m'envoyer le catalogue ? 5) Nous aurions vendu plus de cacao.
\tcblower
Solution détaillée :\par
1) \all{Wenn ich Zeit \textbf{hätte}, \textbf{würde} ich die Fabrik \textbf{besuchen}.} — Condition irréelle : forme forte \all{hätte} en fin de subordonnée ; dans la principale, \all{würde} en position 1 et infinitif \all{besuchen} en fin.\par
2) \all{An Ihrer Stelle \textbf{würde} ich den Vertrag \textbf{annehmen}.} — Formule de conseil ; verbe à particule séparable : à l'infinitif, \all{annehmen} s'écrit en un mot (participe \all{angenommen}).\par
3) \all{Ich \textbf{hätte} gern ein Gespräch mit dem Direktor.} — Politesse : \all{hätte gern} ; \all{das Gespräch} = l'entretien, la conversation.\par
4) \all{\textbf{Könnten} Sie mir bitte den Katalog \textbf{schicken}?} — Demande polie : \all{könnten Sie} + infinitif en fin.\par
5) \all{Wir \textbf{hätten} mehr Kakao \textbf{verkauft}.} — Conditionnel passé : \all{hätten} + participe II \all{verkauft} en fin de phrase.
\end{exoresolu}

\begin{exercice}[À vous de jouer]
Traduisez en allemand : 1) Si les jeunes avaient du travail, ils resteraient au village. 2) Je devrais économiser plus d'argent. 3) Voudriez-vous un café ? 4) Si nous étions riches, nous construirions un hôtel à Kribi. 5) Elle aurait acheté la voiture, mais elle était trop chère.
\end{exercice}

\begin{corrige}[Exercice : le conditionnel]
1) \all{Wenn die Jugendlichen Arbeit hätten, würden sie im Dorf bleiben.} 2) \all{Ich sollte mehr Geld sparen.} 3) \all{Hätten Sie gern einen Kaffee?} (ou \all{Würden Sie gern einen Kaffee trinken?}) 4) \all{Wenn wir reich wären, würden wir ein Hotel in Kribi bauen.} 5) \all{Sie hätte das Auto gekauft, aber es war zu teuer.}
\end{corrige}

\begin{aretenir}[L'essentiel du Konjunktiv II]
\begin{itemize}
\item Formation générale : \all{würde} + infinitif en fin de phrase.
\item Formes fortes obligatoires : \all{wäre} (serais), \all{hätte} (aurais), \all{könnte} (pourrais), \all{müsste} / \all{sollte} (devrais), \all{dürfte}, \all{wüsste}, \all{es gäbe} (il y aurait).
\item Phrase conditionnelle : \all{Wenn ich Zeit hätte, würde ich mehr arbeiten.} — verbe conjugué en fin de subordonnée.
\item Conditionnel passé : \all{hätte/wäre} + participe II : \all{ich hätte gekauft} = j'aurais acheté.
\item Politesse : \all{Ich hätte gern...}, \all{Könnten Sie bitte...?}, \all{Ich würde vorschlagen...}
\end{itemize}
\end{aretenir}

Ce conditionnel nous sera immédiatement utile : dans la section suivante, nous verrons comment le passif (\all{Vorgangspassiv}) permet de décrire les processus économiques — « le cacao est exporté », « les produits sont fabriqués » — avant d'aborder la rédaction de la lettre commerciale, où \all{würden} et \all{könnten} sont des outils de politesse indispensables.

\coursec{Grammaire 3 : le passif (Vorgangspassiv)}

Dans la section précédente, nous avons vu comment le \cle{Konjunktiv II} permet d'exprimer l'hypothèse, le souhait et le conseil : \all{Wenn ich mehr Geld hätte, würde ich ein Geschäft eröffnen.} (« Si j'avais plus d'argent, j'ouvrirais un magasin. »). Nous restons dans le domaine de la vie économique, mais nous changeons de perspective : dans les textes commerciaux, administratifs et journalistiques allemands, ce n'est souvent pas \emph{qui} fait l'action qui compte, mais \emph{ce qui est fait}. Pour cela, l'allemand utilise massivement le \cle{passif}. C'est l'objet de cette troisième révision grammaticale.

\courssub{Observation : le passif dans un texte économique}

Lisons d'abord un court article de presse économique. Soulignez mentalement toutes les formes avec \all{werden} ou \all{worden} suivies d'un participe.

\begin{textebox}[Ein Bericht aus der Wirtschaftsredaktion]
In der Kakaoregion um Yaoundé wird seit einigen Jahren mehr investiert. Neue Maschinen wurden von einer deutschen Firma geliefert, und die Arbeiter wurden in modernen Produktionsmethoden geschult. Der Kakao wird zuerst getrocknet, dann wird er in Säcke verpackt und schließlich nach Europa exportiert. Im Hafen von Douala werden die Container auf Schiffe verladen. „Früher wurde viel Zeit verloren“, erklärt ein Händler, „heute wird alles schneller erledigt.“ Auch die Qualität wird streng kontrolliert: Jede Lieferung muss geprüft werden, bevor sie bezahlt wird. Nächstes Jahr wird eine neue Fabrik gebaut werden. Dort werden dann Schokolade und Kakaobutter hergestellt. Die Gewerkschaft fordert, dass bessere Löhne gezahlt werden. In der Region wird viel diskutiert, aber es wird auch viel gearbeitet.
\end{textebox}

\textbf{Glossaire :} \all{die Maschine, -n} (la machine) ; \all{liefern} (livrer) ; \all{schulen} (former) ; \all{trocknen} (sécher) ; \all{der Sack, „-e} (le sac) ; \all{verpacken} (emballer) ; \all{der Hafen, „-} (le port) ; \all{verladen (verlädt, verlud, verladen)} (charger) ; \all{erledigen} (régler, traiter) ; \all{die Lieferung, -en} (la livraison) ; \all{prüfen} (contrôler) ; \all{die Kakaobutter} (le beurre de cacao) ; \all{die Gewerkschaft, -en} (le syndicat) ; \all{der Lohn, „-e} (le salaire).

\textbf{Questions de compréhension :}
\begin{enumerate}
\item Qui a livré les nouvelles machines ?
\item Que se passe-t-il dans le port de Douala ?
\item Qu'est-ce qui sera construit l'année prochaine ?
\end{enumerate}

Comptez les formes passives : \all{wird investiert, wurden geliefert, wurden geschult, wird getrocknet, wird verpackt, wird exportiert, werden verladen, wurde verloren, wird erledigt, wird kontrolliert, muss geprüft werden, wird bezahlt, wird gebaut werden, werden hergestellt, gezahlt werden, wird diskutiert, wird gearbeitet}. Dix-sept occurrences en onze lignes ! Le passif est donc une structure \emph{centrale} du texte économique allemand.

\courssub{Formation du Vorgangspassiv}

\begin{definition}[Le Vorgangspassiv]
Le \cle{Vorgangspassiv} (passif de processus) présente une action ou un processus du point de vue de celui qui la \emph{subit}. Il se forme avec l'auxiliaire \cle{werden}, conjugué au temps voulu, suivi du \cle{Partizip II} (participe passé) du verbe principal, placé en fin de phrase.
\end{definition}

\begin{propriete}[Règle de formation]
\begin{itemize}
\item \textbf{Présent :} \all{werden} (présent) + Partizip II → \all{Das Produkt wird verkauft.} (« Le produit est vendu. »)
\item \textbf{Präteritum :} \all{werden} (prétérit) + Partizip II → \all{Das Produkt wurde verkauft.} (« Le produit a été vendu. »)
\item \textbf{Perfekt :} \all{sein} (présent) + Partizip II + \all{worden} → \all{Das Produkt ist verkauft worden.} (« Le produit a été vendu. »)
\item \textbf{Futur :} \all{werden} (présent) + Partizip II + \all{werden} (infinitif) → \all{Das Produkt wird verkauft werden.} (« Le produit sera vendu. »)
\item \textbf{Avec modal :} modal conjugué + Partizip II + \all{werden} → \all{Das Produkt muss verkauft werden.} (« Le produit doit être vendu. »)
\item L'\textbf{agent} (celui qui agit) est introduit par \cle{von} (+ datif) ; le \textbf{moyen} ou l'instrument par \cle{durch} (+ accusatif).
\end{itemize}
\end{propriete}

Conjugaison complète de l'auxiliaire \all{werden} au passif :

\begin{center}
\begin{tabular}{|l|l|l|}
\hline
\rowcolor{popblueL} Personne & Präsens & Präteritum \\ \hline
ich & werde & wurde \\ \hline
du & wirst & wurdest \\ \hline
er/sie/es & wird & wurde \\ \hline
wir & werden & wurden \\ \hline
ihr & werdet & wurdet \\ \hline
sie/Sie & werden & wurden \\ \hline
\end{tabular}
\end{center}

\courssub{Le passif à tous les temps : tableau récapitulatif}

\begin{center}
\begin{tabularx}{\linewidth}{|l|X|X|}
\hline
\rowcolor{popblueL} Temps & Deutsch & Français \\ \hline
Präsens & Der Kaffee wird verkauft. & Le café est vendu. \\ \hline
Präteritum & Der Kaffee wurde verkauft. & Le café a été vendu / fut vendu. \\ \hline
Perfekt & Der Kaffee ist verkauft worden. & Le café a été vendu. \\ \hline
Futur I & Der Kaffee wird verkauft werden. & Le café sera vendu. \\ \hline
Mit Modalverb & Der Kaffee muss verkauft werden. & Le café doit être vendu. \\ \hline
Modal im Präteritum & Der Kaffee musste verkauft werden. & Le café devait être vendu. \\ \hline
\end{tabularx}
\end{center}

\begin{exemplebox}[Exemples traduits]
\all{Der Kaffee wird nach Deutschland exportiert.} (« Le café est exporté vers l'Allemagne. »)

\all{Die Waren wurden gestern geliefert.} (« Les marchandises ont été livrées hier. »)

\all{Der Brief ist von der Sekretärin geschrieben worden.} (« La lettre a été écrite par la secrétaire. »)

\all{In dieser Fabrik werden Autos hergestellt.} (« Dans cette usine, on fabrique des voitures. »)

\all{Die Rechnung muss sofort bezahlt werden.} (« La facture doit être payée immédiatement. »)

\all{Die Rechnungen werden per E-Mail geschickt.} (« Les factures sont envoyées par courriel. »)

\all{Der Vertrag wurde durch einen Fehler verzögert.} (« Le contrat a été retardé par une erreur. » — ici \all{durch} exprime la cause/le moyen.)
\end{exemplebox}

\courssub{Transformation actif → passif}

La transformation suit une mécanique précise :

\begin{enumerate}
\item Le \textbf{complément d'objet direct} (accusatif) de la phrase active devient le \textbf{sujet} (nominatif) de la phrase passive.
\item Le verbe actif se transforme : \all{werden} conjugué (au même temps que le verbe actif) + Partizip II en fin de phrase.
\item L'ancien \textbf{sujet} devient complément d'agent introduit par \all{von} (+ datif) ; il est souvent omis s'il est inconnu ou sans importance.
\end{enumerate}

\begin{popfigure}
\begin{center}
\begin{tikzpicture}[
  box/.style={draw=popblue, thick, rounded corners, fill=popblueL, align=center, inner sep=6pt},
  pbox/.style={draw=poporange, thick, rounded corners, fill=poporangeL, align=center, inner sep=6pt},
  lab/.style={font=\small\itshape, text=popdark}
]
\node[box, align=center] (aktiv) at (0,0){\textbf{Aktiv}\\ Der Kunde kauft den Kaffee.};
\node[pbox, align=center] (passiv) at (0,-3.2){\textbf{Passiv}\\ Der Kaffee wird vom Kunden gekauft.};
\draw[-{Stealth[length=3mm]}, very thick, popdark] (aktiv) -- (passiv);
\node[lab] at (3.6,-0.6) {Akkusativobjekt: \all{den Kaffee}};
\node[lab] at (3.6,-2.6) {→ Subjekt: \all{der Kaffee}};
\node[lab] at (-3.8,-0.6) {Subjekt: \all{der Kunde}};
\node[lab] at (-3.8,-2.6) {→ Agent: \all{vom Kunden}};
\node[lab, text=popgreen] at (0,-1.6) {kauft → wird gekauft};
\end{tikzpicture}
\end{center}
\legende{Schéma de transformation de la voix active à la voix passive}
\end{popfigure}

Autres exemples de transformation :

\begin{center}
\begin{tabularx}{\linewidth}{|X|X|}
\hline
\rowcolor{popgreenL} Aktiv & Passiv \\ \hline
Die Firma liefert die Waren. & Die Waren werden (von der Firma) geliefert. \\ \hline
Der Chef unterschrieb den Vertrag. & Der Vertrag wurde (vom Chef) unterschrieben. \\ \hline
Man hat die Preise gesenkt. & Die Preise sind gesenkt worden. \\ \hline
Die Bank wird den Kredit prüfen. & Der Kredit wird (von der Bank) geprüft werden. \\ \hline
\end{tabularx}
\end{center}

\courssub{Le passif impersonnel}

Quand le verbe n'a pas de complément d'objet direct (verbe intransitif), l'allemand construit un \cle{passif impersonnel} avec \all{es} (souvent omis en début de phrase si un autre élément occupe la première place) :

\begin{exemplebox}[Passiv ohne Subjekt]
\all{Es wird viel produziert.} (« On produit beaucoup. »)

\all{Hier wird gearbeitet.} (« Ici, on travaille. »)

\all{In der Region wird viel diskutiert.} (« Dans la région, on discute beaucoup. »)

\all{Am Montag wird nicht geliefert.} (« Le lundi, on ne livre pas. »)
\end{exemplebox}

C'est l'équivalent du français « on + verbe ». Notez que le verbe reste à la \textbf{3\up{e} personne du singulier} : \all{es wird}, jamais \all{es werden}.

\courssub{Vorgangspassiv ou Zustandspassiv ?}

Le français « la porte est fermée » est ambigu : s'agit-il d'une action (quelqu'un la ferme) ou d'un état (elle est fermée) ? L'allemand distingue clairement les deux :

\begin{center}
\begin{tabularx}{\linewidth}{|X|X|X|}
\hline
\rowcolor{popblueL} Français & Deutsch & Sens \\ \hline
La porte est fermée (action) & Die Tür wird geschlossen. & Vorgangspassiv : on la ferme \\ \hline
La porte est fermée (état) & Die Tür ist geschlossen. & Zustandspassiv : elle est fermée \\ \hline
Le magasin est ouvert à 8 h. & Das Geschäft wird um 8 Uhr geöffnet. & action d'ouvrir \\ \hline
Le magasin est ouvert. & Das Geschäft ist geöffnet. & état : il est ouvert \\ \hline
\end{tabularx}
\end{center}

\begin{propriete}[Zustandspassiv]
Le \cle{Zustandspassiv} (passif d'état) se forme avec \all{sein} + Partizip II et décrit le \textbf{résultat} ou l'\textbf{état} après l'action : \all{Die Rechnung ist bezahlt.} (« La facture est payée » — c'est réglé). Le Vorgangspassiv avec \all{werden} décrit l'\textbf{action en cours} : \all{Die Rechnung wird bezahlt.} (« La facture est en train d'être payée / on paie la facture »).
\end{propriete}

\begin{attention}[Pièges fréquents des francophones]
\begin{itemize}
\item \textbf{Ne confondez pas} \all{werden} auxiliaire du passif et \all{werden} verbe « devenir » : \all{Er wird Arzt.} (« Il devient médecin. ») mais \all{Er wird operiert.} (« Il est opéré. »). Au passif, \all{werden} est toujours suivi d'un \textbf{Partizip II}.
\item \textbf{Au Perfekt du passif}, le participe de \all{werden} est \all{worden}, jamais \all{geworden} : \all{Der Brief ist geschrieben worden.} ✓ — \all{Der Brief ist geschrieben geworden.} ✗. (On utilise \all{geworden} seulement pour « devenir » : \all{Er ist Arzt geworden.} « Il est devenu médecin. »)
\item \textbf{Place du participe} : le Partizip II se place en fin de phrase ; au Perfekt, \all{worden} vient après lui, tout à la fin : \all{Die Waren sind gestern von der Firma geliefert worden.}
\item \textbf{Avec un modal}, l'ordre est : modal conjugué en 2\up{e} position … Partizip II + \all{werden} à la toute fin : \all{Die Rechnung muss sofort bezahlt werden.} (et non \all{muss werden bezahlt}).
\item Le français « on » se rend souvent par le passif allemand : « Ici on parle allemand. » → \all{Hier wird Deutsch gesprochen.}
\end{itemize}
\end{attention}

\begin{savaistu}[Pourquoi tant de passifs dans les textes économiques ?]
Le passif est extrêmement fréquent dans les textes économiques, juridiques et administratifs allemands (\all{Amtssprache}), car il met l'accent sur l'\emph{action} et le \emph{résultat} plutôt que sur l'auteur : \all{Die Bestellung wurde bearbeitet.} (« La commande a été traitée. »). C'est aussi une stratégie de politesse ou de prudence : dire \all{Es wurde ein Fehler gemacht} (« Une erreur a été commise ») évite de désigner un coupable ! Dans vos lettres commerciales (section 6), le passif vous donnera un style professionnel et objectif, très apprécié dans la correspondance d'affaires germanophone.
\end{savaistu}

\courssub{Exercice résolu et entraînement}

\begin{exoresolu}[Transformation actif → passif]
Énoncé : Mettez les phrases suivantes au passif, au même temps. 1) \all{Der Bauer erntet den Kakao.} 2) \all{Die Firma hat die Maschinen geliefert.} 3) \all{Man muss die Rechnung bezahlen.} 4) \all{Der Händler verkaufte die Waren.}
\tcblower
Solution détaillée :
\begin{enumerate}
\item Présent : l'accusatif \all{den Kakao} devient sujet \all{der Kakao} ; \all{erntet} → \all{wird geerntet}. → \all{Der Kakao wird (vom Bauern) geerntet.} (« Le cacao est récolté par le paysan. »)
\item Perfekt : auxiliaire \all{sein} + Partizip II + \all{worden}. → \all{Die Maschinen sind (von der Firma) geliefert worden.} (« Les machines ont été livrées par l'entreprise. »)
\item Modal : modal conjugué + Partizip II + \all{werden}. → \all{Die Rechnung muss bezahlt werden.} (« La facture doit être payée. »)
\item Präteritum : \all{wurden} + Partizip II. → \all{Die Waren wurden (vom Händler) verkauft.} (« Les marchandises ont été vendues par le commerçant. »)
\end{enumerate}
\end{exoresolu}

\begin{exercice}[À vous de jouer]
1) Traduisez en allemand : a) Le café est exporté vers l'Allemagne. b) La lettre a été écrite par la secrétaire. c) Les salaires doivent être payés à temps. d) Ici, on travaille beaucoup.

2) Mettez au passif : a) \all{Die Bank prüft den Kredit.} b) \all{Man hat eine neue Fabrik gebaut.}
\end{exercice}

\begin{corrige}[Exercice « À vous de jouer »]
1) a) \all{Der Kaffee wird nach Deutschland exportiert.} b) \all{Der Brief ist von der Sekretärin geschrieben worden.} c) \all{Die Löhne müssen pünktlich gezahlt werden.} d) \all{Hier wird viel gearbeitet.}

2) a) \all{Der Kredit wird (von der Bank) geprüft.} b) \all{Eine neue Fabrik ist gebaut worden.}
\end{corrige}

\begin{aretenir}[L'essentiel du Vorgangspassiv]
Passif = \all{werden} conjugué + Partizip II en fin de phrase. Au Perfekt : \all{sein} + Partizip II + \all{worden}. Avec un modal : modal + Partizip II + \all{werden}. L'agent est introduit par \all{von} (+ Dativ), le moyen par \all{durch} (+ Akkusativ). Le passif impersonnel (\all{Es wird gearbeitet.}) traduit le « on » français. Distinguez \all{Die Tür wird geschlossen} (action) et \all{Die Tür ist geschlossen} (état).
\end{aretenir}

\coursec{Production et traduction de textes liés à la vie économique}

\courssub{Texte support : lecture et compréhension}

\begin{textebox}[Article de presse : Le cacao camerounais sur le marché allemand]
\textbf{Kameruns Kakao: Qualität trifft auf deutsche Technologie}\\
\\
Berlin/Douala. Die Nachfrage nach nachhaltigem Kakao aus Kamerun steigt auf dem deutschen Markt stetig. Im vergangenen Jahr importierte Deutschland rund 45.000 Tonnen Kakao aus Kamerun, was einem Anstieg von 12 \% gegen\up{u}ber dem Vorjahr entspricht. Deutsche Schokoladenhersteller wie Ritter Sport oder Lindt legen zunehmend Wert auf rückverfolgbare Lieferketten und faire Preise für die Kleinbauern in der Region Centre und Süd.\\
\\
Ein neues Kooperationsabkommen zwischen dem kamerunischen Ministerium für Handel (\all{MINCOMMERCE}) und dem Bundesministerium für wirtschaftliche Zusammenarbeit und Entwicklung (\all{BMZ}) soll die Verarbeitungskapazitäten vor Ort ausbauen. „Wir wollen nicht nur Rohstoffe exportieren, sondern Mehrwert im Land schaffen“, erklärte Minister Luc Magloire Mbarga Atanga bei der Unterzeichnung in Yaoundé. Geplant sind der Bau neuer Fermentierungszentren und die Schulung von 5.000 Kooperativen-Mitgliedern in nachhaltigen Anbaumethoden.\\
\\
Dennoch bleiben Herausforderungen: Die Logistik auf dem Weg zum Hafen Douala ist oft mangelhaft, und die Bürokratie verzögert manchmal den Export. Doch der Trend ist klar: Kamerun positioniert sich als verlässlicher Partner für die deutsche Süßwarenindustrie.
\end{textebox}

\textbf{Glossaire :} \all{die Nachfrage, -n} (la demande) ; \all{nachhaltig} (durable) ; \all{der Anstieg, -e} (l'augmentation) ; \all{der Vorjahr, -e} (l'année précédente) ; \all{der Schokoladenhersteller, -} (le fabricant de chocolat) ; \all{rückverfolgbar} (traçable) ; \all{die Lieferkette, -n} (la chaîne d'approvisionnement) ; \all{der Kleinbauer, -n} (le petit paysan) ; \all{das Kooperationsabkommen, -} (l'accord de coopération) ; \all{die Verarbeitungskapazität, -en} (la capacité de transformation) ; \all{vor Ort} (sur place) ; \all{der Mehrwert, -e} (la valeur ajoutée) ; \all{das Fermentierungszentrum, -en} (le centre de fermentation) ; \all{die Kooperative, -n} (la coopérative) ; \all{die Logistik} (la logistique) ; \all{mangelhaft} (déficient) ; \all{verlässlich} (fiable) ; \all{die Süßwarenindustrie} (l'industrie de la confiserie).

\textbf{Questions de compréhension :}
\begin{enumerate}
\item Um wie viel Prozent stiegen die Kakaouimporte aus Kamerun nach Deutschland? (12 \%)
\item Nennen Sie zwei Ziele des neuen Kooperationsabkommens. (Ausbau der Verarbeitungskapazitäten, Schulung von Kooperativen-Mitgliedern)
\item Welche zwei Probleme werden am Ende des Textes genannt? (Logistik, Bürokratie)
\item Finden Sie im Text drei verbes au passif (Vorgangspassiv) und notez leur infinitif. (\all{wird ... exportiert} $\rightarrow$ \all{exportieren} ; \all{wird ... geschaffen} $\rightarrow$ \all{schaffen} ; \all{werden ... verzögert} $\rightarrow$ \all{verzögern})
\end{enumerate}

\courssub{Lexique thématique : la vie économique}

\begin{propriete}[Vocabulaire de base : acteurs et actions]
\begin{center}
\begin{tabularx}{\linewidth}{|l|X|X|}
\hline
\rowcolor{popblueL} \textbf{Deutsch (Artikel, Plural)} & \textbf{Français} & \textbf{Contexte / Exemple} \\
\hline
\all{der Handel, -} ; \all{der Außenhandel, -} & le commerce ; le commerce extérieur & \all{Der Außenhandel wächst.} \\
\hline
\all{der Importeur, -e} / \all{der Exporteur, -e} & l'importateur / l'exportateur & \all{Der Importeur bestellt Ware.} \\
\hline
\all{der Großhändler, -} / \all{der Einzelhändler, -} & le grossiste / le détaillant & \all{Der Großhändler liefert an Einzelhändler.} \\
\hline
\all{die Ware, -n} ; \all{das Produkt, -e} & la marchandise ; le produit & \all{Die Ware ist angekommen.} \\
\hline
\all{der Preis, -e} ; \all{der Weltmarktpreis, -e} & le prix ; le cours mondial & \all{Der Preis ist verhandelbar.} \\
\hline
\all{die Bestellung, -en} & la commande & \all{Die Bestellung ist raus.} \\
\hline
\all{die Lieferung, -en} ; \all{der Liefertermin, -e} & la livraison ; le délai de livraison & \all{Die Lieferung erfolgt pünktlich.} \\
\hline
\all{die Rechnung, -en} ; \all{bezahlen} ; \all{die Zahlung, -en} & la facture ; payer ; le paiement & \all{Die Rechnung wurde bezahlt.} \\
\hline
\all{der Vertrag, -\"{e}r} ; \all{unterschreiben} & le contrat ; signer & \all{Wir unterschreiben den Vertrag.} \\
\hline
\all{die Bank, -en} ; \all{das Konto, -en} ; \all{der Kredit, -e} & la banque ; le compte ; le crédit & \all{Die Bank gewährt einen Kredit.} \\
\hline
\all{die Investition, -en} ; \all{investieren} & l'investissement ; investir & \all{Die Investition lohnt sich.} \\
\hline
\all{der Gewinn, -e} ; \all{der Verlust, -e} ; \all{rentabel} & le bénéfice ; la perte ; rentable & \all{Die Firma macht Gewinn.} \\
\hline
\all{die Wirtschaft, -en} ; \all{wirtschaftlich} & l'économie ; économique & \all{Die wirtschaftliche Lage ist gut.} \\
\hline
\all{der Arbeitsmarkt, -\"{a}rkte} ; \all{die Arbeitslosigkeit} & le marché du travail ; le chômage & \all{Der Arbeitsmarkt erholt sich.} \\
\hline
\end{tabularx}
\end{center}
\end{propriete}

\begin{savaistu}[Le cacao, moteur économique du Cameroun]
Le Cameroun est le 4\up{e}me producteur mondial de cacao (après la Côte d'Ivoire, le Ghana, l'Indonésie). La filière emploie environ 600.000 planteurs. L'Allemagne est le premier client européen du cacao camerounais (fèves brutes et beurre de cacao). Les termes \all{Kakao}, \all{Schokolade}, \all{fairer Handel} (commerce équitable) et \all{Nachhaltigkeit} (durabilité) sont au cœur des relations bilatérales actuelles.
\end{savaistu}

\courssub{Grammaire 1 : Le comparatif et le superlatif}

\begin{propriete}[Formation du comparatif et du superlatif]
\begin{enumerate}
\item \textbf{Comparatif d'inégalité} : \cle{Positiv} + \cle{-er} + \cle{als} (que).
      \all{schnell} $\rightarrow$ \all{schneller als} (plus rapide que).
      \all{teuer} $\rightarrow$ \all{teurer als} (plus cher que).
\item \textbf{Comparatif d'égalité} : \cle{so} + \cle{Positiv} + \cle{wie} (aussi ... que).
      \all{so schnell wie} (aussi rapide que).
\item \textbf{Superlatif} : \cle{am} + \cle{Positiv} + \cle{-(e)sten} (prédicatif) \textbf{ou} article défini + \cle{Positiv} + \cle{-(e)sten} + \cle{adjectival ending} (attributif).
      \all{am schnellsten} (le plus rapide -- prédicatif) ; \all{der schnellste Zug} (le train le plus rapide -- attributif).
\end{enumerate}
\end{propriete}

\begin{attention}[Irrégularités majeures (à connaître par cœur)]
\begin{center}
\begin{tabular}{|l|l|l|l|}
\hline
\rowcolor{popblueL} \textbf{Positif} & \textbf{Comparatif} & \textbf{Superlatif (prédicatif)} & \textbf{Superlatif (attributif)} \\
\hline
\all{gut} & \all{besser} & \all{am besten} & \all{der beste} \\
\hline
\all{viel} & \all{mehr} & \all{am meisten} & \all{die meisten} \\
\hline
\all{gerne} & \all{lieber} & \all{am liebsten} & --- \\
\hline
\all{hoch} & \all{höher} & \all{am höchsten} & \all{der höchste} \\
\hline
\all{nah} & \all{näher} & \all{am nächsten} & \all{der nächste} \\
\hline
\end{tabular}
\end{center}
\emph{Note :} \all{hoch, nah} prennent un Umlaut au comparatif/superlatif. \all{teuer, sauer} $\rightarrow$ \all{teurer, saurer} (pas d'Umlaut, \emph{e} muet tombe).
\end{attention}

\begin{exemplebox}[Comparatif et superlatif en contexte économique]
\begin{itemize}
\item \all{Der Kaffee aus Kamerun ist \textbf{günstiger als} der aus Brasilien.} (Le café du Cameroun est moins cher que celui du Brésil.)
\item \all{Diese Maschine ist \textbf{die effizienteste} in der Fabrik.} (Cette machine est la plus efficace de l'usine.)
\item \all{Wir exportieren \textbf{mehr Kakao als} im Vorjahr.} (Nous exportons plus de cacao que l'an dernier.)
\item \all{Unser Angebot ist \textbf{am attraktivsten}.} (Notre offre est la plus attractive.)
\end{itemize}
\end{exemplebox}

\courssub{Grammaire 2 : Le conditionnel (Konjunktiv II)}

\begin{propriete}[Formation du Konjunktiv II (Présent)]
Deux formes coexistent :
\begin{enumerate}
\item \textbf{Forme \all{würde} + Infinitif} (valable pour \textbf{tous} les verbes, très fréquent à l'oral et à l'écrit standard).
      \all{ich würde machen, du würdest machen, er würde machen, wir würden machen, ihr würdet machen, sie würden machen}.
\item \textbf{Forme simple (une seule mot)} : basée sur le radical du \cle{Präteritum} + Umlaut (si possible) + terminaisons \all{-e, -est, -e, -en, -et, -en}. Réservée aux verbes très fréquents (\all{sein, haben, werden}, modaux, \all{gehen, kommen, stehen, lassen}).
\end{enumerate}
\end{propriete}

\begin{center}
\begin{tabular}{|l|l|l|l|}
\hline
\rowcolor{popblueL} \textbf{Infinitif} & \textbf{Präteritum} & \textbf{Konjunktiv II (Forme simple)} & \textbf{Traduction (je)} \\
\hline
\all{sein} & \all{war} & \all{wäre} & je serais \\
\hline
\all{haben} & \all{hatte} & \all{hätte} & j'aurais \\
\hline
\all{werden} & \all{wurde} & \all{würde} & je deviendrais / j'aurais (auxiliaire) \\
\hline
\all{können} & \all{konnte} & \all{könnte} & je pourrais \\
\hline
\all{müssen} & \all{musste} & \all{müsste} & je devrais \\
\hline
\all{dürfen} & \all{durfte} & \all{dürfte} & je pourrais (autorisation) \\
\hline
\all{wollen} & \all{wollte} & \all{wollte} & je voudrais \\
\hline
\all{sollen} & \all{sollte} & \all{sollte} & je devrais (conseil) \\
\hline
\all{mögen} & \all{mochte} & \all{möchte} & je voudrais \\
\hline
\all{gehen} & \all{ging} & \all{ginge} & j'irais \\
\hline
\all{kommen} & \all{kam} & \all{käme} & je viendrais \\
\hline
\all{stehen} & \all{stand} & \all{stünde} & je serais (debout/situé) \\
\hline
\all{lassen} & \all{ließ} & \all{ließe} & je laisserais / je ferais faire \\
\hline
\end{tabular}
\end{center}

\begin{propriete}[Emplois principaux du Konjunktiv II]
\begin{enumerate}
\item \textbf{Politesse (Höflichkeit)} : atténue une demande, une offre, un refus. \all{\textbf{Könnten} Sie mir helfen?} (Pourriez-vous m'aider ?) ; \all{Ich \textbf{möchte} einen Termin vereinbaren.} (Je voudrais prendre rendez-vous.)
\item \textbf{Irréel du présent/futur (avec \all{wenn})} : \all{Wenn ich Zeit \textbf{hätte}, \textbf{würde} ich kommen.} (Si j'avais le temps, je viendrais.) $\rightarrow$ \all{wenn} + KII, principale + KII.
\item \textbf{Irréel du passé (avec \all{wenn})} : \all{Wenn ich es \textbf{gewusst hätte}, \textbf{wäre} ich gekommen.} (Si je l'avais su, je serais venu.) $\rightarrow$ \all{wenn} + KII (Perfekt), principale + KII (Perfekt).
\item \textbf{Style indirect (discours rapporté)} : \all{Er sagte, er \textbf{käme} morgen.} (Il a dit qu'il viendrait demain.)
\end{enumerate}
\end{propriete}

\begin{attention}[Pièges francophones : \all{wäre} vs \all{würde}]
\all{wäre} = conditionnel de \all{sein} (être). \all{würde} = conditionnel de \all{werden} (devenir / auxiliaire futur/passif).
\begin{itemize}
\item \all{Ich \textbf{wäre} gerne Arzt.} (Je serais gerne médecin.) -- \textbf{Être}
\item \all{Ich \textbf{würde} gerne Arzt \textbf{werden}.} (Je voudrais devenir médecin.) -- \textbf{Devenir}
\item \all{Das \textbf{wäre} gut.} (Ce serait bien.) -- \textbf{Être}
\item \all{Das \textbf{würde} mich freuen.} (Cela me ferait plaisir.) -- \textbf{Auxiliaire} + Infinitif
\end{itemize}
\end{attention}

\courssub{Grammaire 3 : Le passif (Vorgangspassiv)}

\begin{definition}[Vorgangspassiv vs Zustandspassiv]
\begin{itemize}
\item \textbf{Vorgangspassiv} (passif d'action / processus) : \cle{werden} + \cle{Partizip II}. Insiste sur l'action. \all{Die Fabrik \textbf{wird gebaut}.} (L'usine est en train d'être construite / est construite.)
\item \textbf{Zustandspassiv} (passif d'état / résultat) : \cle{sein} + \cle{Partizip II}. Insiste sur l'état résultant. \all{Die Fabrik \textbf{ist gebaut}.} (L'usine est construite / achevée.)
\end{itemize}
Seuls les verbes transitifs (accusatif) peuvent former un passif personnel. Verbes intransitifs $\rightarrow$ \cle{Passif impersonnel} : \all{Es wird gearbeitet.} (On travaille / Il est travaillé.)
\end{definition}

\begin{propriete}[Conjugaison du Vorgangspassiv (ex. \all{machen})]
\begin{center}
\begin{tabular}{|l|l|l|}
\hline
\rowcolor{popblueL} \textbf{Temps} & \textbf{Allemand} & \textbf{Français} \\
\hline
Präsens & \all{werden gemacht} & est fait / est en train d'être fait \\
\hline
Präteritum & \all{wurden gemacht} & a été fait / était fait \\
\hline
Perfekt & \all{sind gemacht worden} & a été fait (passé composé) \\
\hline
Plusquamperfekt & \all{waren gemacht worden} & avait été fait \\
\hline
Futur I & \all{werden gemacht werden} & sera fait \\
\hline
Modaux (Präsens) & \all{müssen gemacht werden} & doivent être faits \\
\hline
Modaux (Präteritum/KII) & \all{müssten gemacht werden} / \all{konnten gemacht werden} & devraient / pourraient être faits \\
\hline
\end{tabular}
\end{center}
\emph{Rappel :} Au Perfekt/Plusquamperfekt, pas de \all{ge-} devant \all{worden} (\all{geworden} n'existe pas au passif).
\end{propriete}

\begin{exemplebox}[Passif dans les textes économiques]
\begin{itemize}
\item \all{Die Preise \textbf{werden monatlich festgelegt}.} (Présens -- Les prix sont fixés mensuellement.)
\item \all{Der Vertrag \textbf{wurde gestern unterschrieben}.} (Präteritum -- Le contrat a été signé hier.)
\item \all{Die Maschinen \textbf{sind aus Deutschland importiert worden}.} (Perfekt -- Les machines ont été importées d'Allemagne.)
\item \all{Die Lieferung \textbf{kann per Schiff erfolgen} / \textbf{kann per Schiff geliefert werden}.} (Modal -- La livraison peut se faire par bateau.)
\item \all{Es \textbf{wurde beschlossen}, die Produktion zu steigern.} (Impersonnel -- Il a été décidé d'augmenter la production.)
\end{itemize}
\end{exemplebox}

\begin{attention}[Agent et instrument : \all{von} vs \all{durch}]
\begin{itemize}
\item \all{von} + Dativ : agent vivant, auteur de l'action. \all{Der Brief wurde \textbf{von dem Direktor} unterschrieben.}
\item \all{durch} + Accusatif : instrument, moyen, cause inanimée (souvent avec une nuance de "par le biais de"). \all{Die Daten wurden \textbf{durch einen Virus} gelöscht.}
\item Si l'agent n'est pas important (style objectif), on l'omet : \all{Es wird Deutsch gesprochen.}
\end{itemize}
\end{attention}

\courssub{La lettre commerciale et le compte rendu}

Dans les sections précédentes, nous avons révisé le comparatif, le conditionnel (Konjunktiv II) et le passif (Vorgangspassiv). Nous allons maintenant réinvestir cette grammaire dans les deux types de textes que l'examen et la vie professionnelle exigent : la \cle{lettre commerciale} (\all{der Geschäftsbrief}) et le \cle{compte rendu} (\all{der Bericht}).

\courssub{Observer d'abord : une lettre commerciale authentique}

Lisons d'abord un modèle complet. Repérez la mise en page, les formules figées et les emplois du conditionnel de politesse.

\begin{textebox}[Modèle de lettre commerciale : une offre]
Kamerun-Kaffee GmbH\\
Postfach 4521, Douala\\
\\
Müller Import AG\\
Hafenstraße 12\\
20457 Hamburg\\
\\
Douala, den 15. März 2025\\
\\
\textbf{Betreff: Ihr Angebot für Arabica-Kaffee}\\
\\
Sehr geehrte Damen und Herren,\\
\\
wir beziehen uns auf Ihre Anfrage vom 10. März. Wir \textbf{könnten} Ihnen 500 kg Arabica-Kaffee zum Preis von 8 Euro pro Kilo anbieten. Die Lieferung \textbf{würde} innerhalb von zwei Wochen erfolgen. Die Bezahlung kann per Banküberweisung erfolgen.\\
\\
Vielen Dank im Voraus. Wir freuen uns auf Ihre Bestellung.\\
\\
Mit freundlichen Grüßen\\
\\
A. Mbarga
\end{textebox}

\textbf{Glossaire :} \all{die Anfrage, -n} (la demande de renseignements) ; \all{das Angebot, -e} (l'offre) ; \all{die Bestellung, -en} (la commande) ; \all{die Lieferung, -en} (la livraison) ; \all{die Banküberweisung, -en} (le virement bancaire) ; \all{erfolgen} (avoir lieu, être effectué) ; \all{beziehen sich auf} + Acc (se référer à).

\textbf{Questions de compréhension :}
\begin{enumerate}
\item Wer schreibt diesen Brief, und an wen? (\all{Kamerun-Kaffee GmbH} écrit à \all{Müller Import AG}.)
\item Worauf bezieht sich der Absender? (À la \all{Anfrage vom 10. März}.)
\item Was bietet die Firma an, und zu welchem Preis? (500 kg Arabica à 8 €/kg.)
\item Welche Verbformen drücken die Höflichkeit aus? (\all{könnten anbieten}, \all{würde erfolgen} -- Konjunktiv II.)
\end{enumerate}

\textbf{Observation grammaticale :} Remarquez \all{wir \textbf{könnten} ... anbieten} et \all{die Lieferung \textbf{würde} erfolgen} : c'est le \cle{Konjunktiv II}, le conditionnel de politesse, obligatoire dans la correspondance commerciale. Notez aussi le passif possible : \all{Die Ware \textbf{wird geliefert}.} (Présens) ou \all{Die Ware \textbf{wurde geliefert}.} (Präteritum).

\courssub{La structure de la lettre commerciale (der Geschäftsbrief)}

La lettre commerciale allemande suit une mise en page rigide (norme DIN 5008). Toute erreur de structure est sanctionnée à l'examen.

\begin{propriete}[Les huit zones de la lettre commerciale]
\begin{enumerate}
\item \all{der Absender} : l'expéditeur (nom, adresse), en haut à gauche.
\item \all{der Empfänger} : le destinataire, sous l'expéditeur.
\item \all{Ort und Datum} : lieu et date, à droite : \all{Douala, den 15. März 2025} (accusatif de temps avec \all{den}).
\item \all{der Betreff} : l'objet, en gras, sans verbe conjugué, sans "Betreff:" dans la norme la plus récente mais toléré : \all{Betreff: Ihre Anfrage vom 10. März}.
\item \all{die Anrede} : la formule d'appel, suivie d'une \textbf{virgule}.
\item \all{der Brieftext} : le corps, qui commence par une \textbf{minuscule}.
\item \all{die Grußformel} : la formule de politesse finale, \textbf{sans virgule}.
\item \all{die Unterschrift} : la signature (manuscrite), nom tapé en dessous.
\end{enumerate}
\end{propriete}

\begin{popfigure}
\begin{tikzpicture}[font=\small, scale=0.85]
\draw[thick, popblue, rounded corners] (0,0) rectangle (12,14);
% Absender
\node[anchor=north west, draw=poporange, fill=poporangeL, rounded corners, align=left, text width=5cm] at (0.4,13.7) {Kamerun-Kaffee GmbH\\Postfach 4521\\Douala (Kamerun)};
\node[anchor=west, popdark, font=\bfseries\small] at (12.3,13.3) {Absender};
% Empfänger
\node[anchor=north west, draw=popgreen, fill=popgreenL, rounded corners, align=left, text width=5cm] at (0.4,11.8) {Müller Import AG\\Hafenstraße 12\\20457 Hamburg (Deutschland)};
\node[anchor=west, popdark, font=\bfseries\small] at (12.3,11.3) {Empfänger};
% Datum
\node[anchor=north east, draw=poppurple, fill=poppurpleL, rounded corners] at (11.6,10.5) {Douala, den 15. März 2025};
\node[anchor=west, popdark, font=\bfseries\small] at (12.3,10.5) {Datum};
% Betreff
\node[anchor=north west, draw=popgold, fill=popgoldL, rounded corners] at (0.4,9.8) {\textbf{Betreff: Ihr Angebot für Arabica-Kaffee}};
\node[anchor=west, popdark, font=\bfseries\small] at (12.3,9.8) {Betreff};
% Anrede
\node[anchor=north west, draw=popblue, fill=popblueL, rounded corners] at (0.4,9.0) {Sehr geehrte Damen und Herren,};
\node[anchor=west, popdark, font=\bfseries\small] at (12.3,9.0) {Anrede};
% Brieftext
\node[anchor=north west, draw=popink, fill=popcream, rounded corners, align=left, text width=10.5cm] at (0.4,8.2) {wir beziehen uns auf Ihre Anfrage vom 10. März.\\Wir \textbf{könnten} Ihnen 500 kg Kaffee anbieten.\\Die Lieferung \textbf{würde} in zwei Wochen erfolgen.\\[3pt]Vielen Dank im Voraus. Wir freuen uns auf Ihre Bestellung.};
\node[anchor=west, popdark, font=\bfseries\small] at (12.3,7.0) {Brieftext};
% Grußformel
\node[anchor=north west, draw=poppink, fill=poppinkL, rounded corners] at (0.4,4.8) {Mit freundlichen Grüßen};
\node[anchor=west, popdark, font=\bfseries\small] at (12.3,4.8) {Grußformel};
% Unterschrift
\node[anchor=north west, draw=popmuted, fill=popblueL, rounded corners, align=center] at (0.4,4.0){A. Mbarga\\Geschäftsführer};
\node[anchor=west, popdark, font=\bfseries\small] at (12.3,4.0) {Unterschrift};
\end{tikzpicture}
\legende{Maquette d'une lettre commerciale allemande : les huit zones obligatoires (norme DIN 5008).}
\end{popfigure}

\begin{attention}[Virgules et majuscules : les deux pièges classiques]
\begin{itemize}
\item Après \all{Sehr geehrte Damen und Herren,} il y a une \textbf{virgule}, et le corps de la lettre commence par une \textbf{minuscule} : \all{Sehr geehrte Damen und Herren,\\\textbf{w}ir beziehen uns auf \dots} — contrairement au français qui met une majuscule après la virgule.
\item \all{Mit freundlichen Grüßen} ne prend \textbf{jamais de virgule} à la fin. Écrire \all{Mit freundlichen Grüßen,} est une faute grave.
\item Abréviation \all{MfG} (\all{Mit freundlichen Grüßen}) : uniquement dans les courriels informels, \textbf{interdite} dans une lettre d'examen formelle.
\item Majuscules : \textbf{Tous les noms} prennent une majuscule (\all{Anfrage, Angebot, Lieferung, Banküberweisung, Grüßen}).
\end{itemize}
\end{attention}

\courssub{Les quatre types de lettres commerciales}

\begin{definition}[Les types de lettres commerciales]
\begin{itemize}
\item \all{die Anfrage (-n)} : la demande de renseignements (catalogue, prix, conditions).
\item \all{das Angebot (-e)} : l'offre (proposition de marchandises, prix, délai).
\item \all{die Bestellung (-en)} : la commande (quantité précise, référence).
\item \all{die Beschwerde / die Reklamation (-en)} : la réclamation (livraison défectueuse, retard, facture erronée).
\end{itemize}
\end{definition}

\begin{center}
\begin{tabularx}{\linewidth}{|l|X|X|}
\hline
\rowcolor{popblueL} \textbf{Type} & \textbf{Formules figées (Deutsch)} & \textbf{Français} \\
\hline
Anfrage & \all{Wir bitten Sie um Zusendung Ihres Katalogs.} & Nous vous prions de nous envoyer votre catalogue. \\
\hline
Anfrage & \all{Könnten Sie uns bitte eine Preisliste schicken?} & Pourriez-vous nous envoyer une liste de prix ? \\
\hline
Angebot & \all{Wir beziehen uns auf Ihre Anfrage vom \dots} & Nous nous référons à votre demande du \dots \\
\hline
Angebot & \all{Wir können Ihnen \dots zum Preis von \dots anbieten.} & Nous pouvons vous offrir \dots au prix de \dots \\
\hline
Angebot & \all{Die Lieferung erfolgt innerhalb von 14 Tagen.} & La livraison a lieu sous 14 jours. \\
\hline
Bestellung & \all{Wir möchten 100 Säcke Kakao bestellen.} & Nous aimerions commander 100 sacs de cacao. \\
\hline
Bestellung & \all{Bitte liefern Sie die Ware bis zum 30. April.} & Veuillez livrer la marchandise avant le 30 avril. \\
\hline
Reklamation & \all{Leider müssen wir uns über die Lieferung beschweren.} & Malheureusement, nous devons nous plaindre de la livraison. \\
\hline
Reklamation & \all{Die Ware war beschädigt. Wir fordern Ersatzlieferung.} & La marchandise était endommagée. Nous exigeons un remplacement. \\
\hline
Reklamation & \all{Wir bitten um Gutschrift / Nachbesserung.} & Nous demandons un avoir / une réparation. \\
\hline
Clôture & \all{Vielen Dank im Voraus.} & Merci d'avance. \\
\hline
Clôture & \all{Wir freuen uns auf Ihre Antwort / Bestellung.} & Nous attendons votre réponse/commande avec plaisir. \\
\hline
\end{tabularx}
\end{center}

\begin{propriete}[Le conditionnel de politesse dans la lettre]
Le \cle{Konjunktiv II} adoucit les demandes et les offres ; il est la marque de la courtoisie professionnelle. On privilégie \all{könnten, würden, möchten, hätten}.
\begin{itemize}
\item \all{\textbf{Könnten} Sie uns bitte ein Angebot \textbf{schicken}?} (Pourriez-vous nous envoyer une offre ?)
\item \all{Wir \textbf{würden} uns freuen, bald von Ihnen zu \textbf{hören}.} (Nous serions heureux d'avoir bientôt de vos nouvelles.)
\item \all{\textbf{Hätten} Sie 200 Säcke Kakao auf Lager?} (Auriez-vous 200 sacs de cacao en stock ?)
\item \all{Wir \textbf{könnten} Ihnen einen \textbf{günstigeren} Preis anbieten.} (Nous pourrions vous offrir un prix plus avantageux -- réinvestissement du comparatif \all{günstiger}.)
\item \all{Wir \textbf{möchten} Sie bitten, die Rechnung zu \textbf{prüfen}.} (Nous aimerions vous prier de vérifier la facture.)
\end{itemize}
\end{propriete}

\begin{exemplebox}[Exemples traduits : lettre commerciale]
\begin{itemize}
\item \all{Wir möchten 100 Säcke Kakao bestellen.} « Nous aimerions commander 100 sacs de cacao. »
\item \all{Leider müssen wir uns über die verspätete Lieferung beschweren.} « Malheureusement, nous devons nous plaindre de la livraison en retard. »
\item \all{Wir möchten Sie bitten, uns die Rechnung zu schicken.} « Nous aimerions vous prier de nous envoyer la facture. »
\item \all{Die Lieferung würde innerhalb von zwei Wochen erfolgen.} « La livraison serait effectuée dans un délai de deux semaines. »
\item \all{Könnten Sie uns mitteilen, wann die Ware versandt wird?} « Pourriez-vous nous indiquer quand la marchandise sera expédiée ? »
\end{itemize}
\end{exemplebox}

\begin{attention}[Faux amis et pièges de traduction -- Lettre commerciale]
\begin{itemize}
\item \all{sich beziehen auf} + Accusatif = se référer à ; \textbf{ne pas} traduire par « se baser sur » (\all{sich basieren auf}).
\item \all{sich freuen auf} + Accusatif = attendre avec plaisir (futur/anticipation) ; \all{sich freuen über} + Accusatif = se réjouir de (passé/présent). Dans une lettre commerciale : toujours \all{auf Ihre Antwort / Bestellung / Rückmeldung}.
\item \all{die Bestellung} = la commande (commerciale), pas « l'ordre » (militaire : \all{der Befehl}).
\item \all{im Voraus} = d'avance ; \textbf{orthographe} : un seul \all{r} (pas *\all{im Vorraus}).
\item \all{der Preis} (le prix) vs \all{der Preise} (pluriel) ; \all{zum Preis von} = \all{zu dem Preis von} (Datif).
\item \all{liefern} (livrer) $\neq$ \all{beliefern} (approvisionner qqn : \all{wir beliefern Sie} = nous vous livrons / nous vous approvisionnons).
\end{itemize}
\end{attention}

\begin{methode}[Rédiger une lettre commerciale en quatre étapes]
\begin{enumerate}
\item \textbf{Respecter la mise en page} : placer les huit zones (Absender, Empfänger, Datum, Betreff, Anrede, Brieftext, Grußformel, Unterschrift) aux bons endroits. Utiliser les guillemets allemands „ “ pour le nom de l'entreprise si nécessaire.
\item \textbf{Utiliser les formules figées} : ouvrir avec \all{wir beziehen uns auf \dots} (réponse) ou \all{wir bitten Sie um \dots} (initiative) ; fermer avec \all{vielen Dank im Voraus} et \all{wir freuen uns auf Ihre Antwort}.
\item \textbf{Employer le conditionnel de politesse (Konjunktiv II)} : \all{könnten Sie \dots ?}, \all{wir würden uns freuen \dots}, \all{wir möchten \dots} — \textbf{jamais d'impératif brutal} envers un partenaire commercial (\all{Schicken Sie!} $\rightarrow$ \all{Könnten Sie schicken?}).
\item \textbf{Relire (Check-list BAC)} : 
      \begin{itemize}
      \item Majuscule à \textbf{tous les noms} ?
      \item Verbe conjugué en \textbf{2\up{e} position} (phrase principale) / à la \textbf{fin} (subordonnée) ?
      \item Infinitif / Participe II en \textbf{fin de phrase} ?
      \item Cas corrects après prépositions : \all{auf Ihre Anfrage} (Acc), \all{zum Preis von} (Dativ), \all{bei der Lieferung} (Dativ), \all{wegen des Wetters} (Génitif) ?
      \item Ponctuation : virgule après Anrede, \textbf{pas} de virgule après Grußformel.
      \end{itemize}
\end{enumerate}
\end{methode}

\courssub{Le compte rendu (der Bericht)}

Le compte rendu rapporte des faits — une réunion, une visite d'entreprise, un salon économique, une assemblée générale — de manière \cle{objective} et \cle{structurée}. À la différence de la lettre, il n'y a \textbf{ni destinataire nommé, ni formules de politesse} (Anrede, Grußformel).

\begin{propriete}[Structure du compte rendu en trois parties]
\begin{enumerate}
\item \all{die Einleitung} (introduction) : répondre aux quatre questions \all{wer? was? wann? wo?} (qui, quoi, quand, où). Souvent une seule phrase. Exemple : \all{Am Freitag, den 12. April, besuchte unsere Klasse die Kakaofabrik in Douala.}
\item \all{der Hauptteil} (développement) : les faits dans l'ordre chronologique ou thématique, reliés par des connecteurs logiques. Style télégraphique évité : phrases complètes.
\item \all{der Schluss} (conclusion) : les résultats, les décisions prises, éventuellement une appréciation brève et neutre (pas d'émotion forte).
\end{enumerate}
\end{propriete}

\begin{popfigure}
\begin{tikzpicture}[font=\small, node distance=0.8cm]
\node[draw=popblue, fill=popblueL, rounded corners, align=center, text width=5.5cm] (ein) {\textbf{Einleitung}\\Wer? Was? Wann? Wo?\\Präteritum};
\node[draw=poporange, fill=poporangeL, rounded corners, align=center, text width=5.5cm, below=of ein] (haupt) {\textbf{Hauptteil}\\Faits chronologiques\\Connecteurs : zuerst, dann, danach, außerdem, schließlich\\Inversion sujet-verbe !};
\node[draw=popgreen, fill=popgreenL, rounded corners, align=center, text width=5.5cm, below=of haupt] (schluss) {\textbf{Schluss}\\Résultats, décisions\\Connecteurs : deshalb, daher, insgesamt\\Passif fréquent};
\draw[-{Stealth}, thick, popdark] (ein) -- (haupt);
\draw[-{Stealth}, thick, popdark] (haupt) -- (schluss);
\end{tikzpicture}
\legende{Organigramme : la structure du compte rendu en trois parties.}
\end{popfigure}

\begin{propriete}[Les temps et le style du compte rendu]
\begin{itemize}
\item \textbf{Präteritum} (temps du récit écrit) pour les faits principaux : \all{Die Delegation \textbf{kam} um 9 Uhr \textbf{an}. Der Direktor \textbf{erklärte} die Produktion.}
\item \textbf{Perfekt} possible pour un compte rendu oral ou un style moins formel : \all{Wir \textbf{haben} die Fabrik \textbf{besucht}.}
\item \textbf{Style objectif} : pas de \all{ich finde}, pas d'exclamations, pas de \all{ich denke}. Phrases claires, verbe en 2\up{e} position.
\item \textbf{Passif bienvenu} pour l'objectivité : \all{Es \textbf{wurde beschlossen}, die Preise zu senken.} (Il a été décidé de baisser les prix.) -- Sujet impersonnel \all{es}.
\end{itemize}
\end{propriete}

\begin{definition}[Les connecteurs du compte rendu (Adverbes / Konjunktionaladverbien)]
Ces mots placent en \textbf{position 1} $\rightarrow$ \textbf{inversion sujet-verbe} (Verbe en position 2).
\begin{center}
\begin{tabular}{|l|l|l|}
\hline
\rowcolor{popblueL} \textbf{Connecteur} & \textbf{Sens} & \textbf{Exemple (inversion !)} \\
\hline
\all{zuerst} & d'abord & \all{Zuerst \textbf{begrüßte} der Direktor die Gäste.} \\
\hline
\all{dann} / \all{anschließend} & puis / ensuite & \all{Dann \textbf{führte} er uns durch die Halle.} \\
\hline
\all{danach} & après cela & \all{Danach \textbf{stellten} wir Fragen.} \\
\hline
\all{außerdem} & de plus & \all{Außerdem \textbf{erhielten} wir Proben.} \\
\hline
\all{schließlich} / \all{zum Schluss} & finalement / à la fin & \all{Schließlich \textbf{wurde} das Protokoll \textbf{unterschrieben}.} \\
\hline
\all{deshalb} / \all{deswegen} / \all{daher} & c'est pourquoi / par conséquent & \all{Deshalb \textbf{plant} die Firma eine Expansion.} \\
\hline
\all{trotzdem} & néanmoins & \all{Trotzdem \textbf{blieb} die Stimmung gut.} \\
\hline
\end{tabular}
\end{center}
\end{definition}

\begin{textebox}[Modèle de compte rendu : visite d'entreprise]
\textbf{Bericht über den Besuch der Kakaofabrik Chococam in Douala}\\
\\
Am Dienstag, den 8. April 2025, besuchte die Klasse Tle A2 des Gymnasiums Yaoundé die Kakaofabrik in Douala. Die Abfahrt war um 6 Uhr morgens.\\
\\
Zuerst wurden wir vom Direktor der Fabrik begrüßt. Er erklärte uns die Geschichte des Unternehmens und die Bedeutung des Kakaos für die Wirtschaft Kameruns. Danach besichtigten wir die Produktionshallen. Dort sahen wir, wie die Kakaobohnen gewaschen, getrocknet und gemahlen werden. Außerdem zeigte uns eine Ingenieurin die Maschinen, die aus Deutschland importiert wurden. Dann konnten wir Fragen stellen: Wir wollten wissen, wie viele Arbeiter in der Fabrik beschäftigt sind und wohin die Schokolade exportiert wird.\\
\\
Schließlich gab es eine Verkostung verschiedener Schokoladensorten. Trotzdem blieb genug Zeit für ein Gruppenfoto vor dem Fabrikgebäude.\\
\\
Der Besuch war sehr lehrreich. Wir haben gelernt, dass der Kakao aus Kamerun in viele Länder Europas exportiert wird. Deshalb möchte unsere Klasse im nächsten Jahr eine Kaffeeplantage in der Westregion besuchen.
\end{textebox}

\textbf{Glossaire :} \all{die Abfahrt, -en} (le départ) ; \all{die Bedeutung, -en} (l'importance) ; \all{die Produktionshalle, -n} (l'atelier de production) ; \all{die Kakaobohne, -n} (la fève de cacao) ; \all{mahlen} (moudre -- fort : mahlte, gemahlen) ; \all{beschäftigen} (employer) ; \all{die Verkostung, -en} (la dégustation) ; \all{lehrreich} (instructif) ; \all{die Westregion, -en} (la région de l'Ouest).

\textbf{Observation guidée :} Relevez dans le texte :
\begin{itemize}
\item 4 formes de \textbf{Vorgangspassiv} : \all{wurden begrüßt} (Präteritum), \all{werden gewaschen/getrocknet/gemahlen} (Présens), \all{wurden importiert} (Präteritum), \all{wird exportiert} (Présens).
\item Les temps : \textbf{Präteritum} dominant (\all{besuchte, war, erklärte, besichtigten, sahen, zeigte, konnten, wollten, gab, blieb, war}), un \textbf{Perfekt} final (\all{haben gelernt}) pour marquer la conséquence présente.
\item Les connecteurs en tête de phrase avec inversion : \all{Zuerst wurden...}, \all{Danach besichtigten...}, \all{Außerdem zeigte...}, \all{Schließlich gab...}, \all{Trotzdem blieb...}, \all{Deshalb möchte...}.
\end{itemize}

\begin{exoresolu}[Transformer des notes en compte rendu]
Énoncé : À partir de ces notes, rédigez le début d'un compte rendu (3 phrases, Präteritum, connecteurs) : \emph{vendredi 20 juin — réunion des commerçants du marché Mvog-Mbi — sujet : les prix — décision : baisser les prix du manioc.}\tcblower
Solution : \all{Am Freitag, den 20. Juni, fand auf dem Markt Mvog-Mbi in Yaoundé eine Versammlung der Händler statt. Zuerst diskutierte man über die Preise. Schließlich wurde beschlossen, die Preise für Maniok zu senken.}\\
Traduction : « Vendredi 20 juin, une réunion des commerçants a eu lieu au marché Mvog-Mbi à Yaoundé. D'abord, on a discuté des prix. Finalement, il a été décidé de baisser les prix du manioc. »\\
Points de méthode : 
\begin{enumerate}
\item L'introduction répond à \all{wann?} (20. Juni), \all{wo?} (Markt Mvog-Mbi), \all{wer?} (Händler), \all{was?} (Versammlung). Verbe \all{fand ... statt} (trennbares Verb).
\item \all{Zuerst} (Pos 1) $\rightarrow$ inversion \all{diskutierte man}. \all{man} = on (impersonnel).
\item \all{Schließlich} (Pos 1) $\rightarrow$ inversion \all{wurde beschlossen}. Passif Präteritum pour l'objectivité. \all{zu senken} (Infinitivgruppe) à la fin.
\end{enumerate}
\end{exoresolu}

\courssub{Traduction : thème et version}

\begin{methode}[Méthodologie de la traduction (Thème / Version)]
\begin{enumerate}
\item \textbf{Analyse} : Identifier le type de texte (lettre, rapport, article), le registre (formel), les temps, la structure.
\item \textbf{Vocabulaire} : Noter les termes techniques (\all{Anfrage, Liefertermin, Reklamation, Vorgangspassiv, Konjunktiv II}). Attention aux faux amis.
\item \textbf{Grammaire} : Repérer les structures clés (passif, conditionnel, comparatif, ordre des mots, cas prépositionnels).
\item \textbf{Brouillon} : Traduire phrase par phrase sans mot-à-mot. Respecter la syntaxe allemande (verbe en 2, infinitif à la fin, inversion après connecteur).
\item \textbf{Relecture} : Vérifier majuscules (Noms), Umlaute/ß, accords (genre/nombre/cas), terminaisons verbales, prépositions (\all{auf} + Acc, \all{von} + Dat, \all{wegen} + Gen).
\end{enumerate}
\end{methode}

\begin{exercice}[Thème : Français $\rightarrow$ Allemand]
Traduisez les phrases suivantes en allemand correct (soignez l'ordre des mots, les cas et les temps).
\begin{enumerate}
\item Notre entreprise (f) peut vous offrir un meilleur prix. (Utilisez \all{könnten} + comparatif).
\item La marchandise a été livrée hier par camion. (Passif Präteritum, agent \all{durch}).
\item Si nous avions le temps, nous visiterions l'usine. (Irréel présent : \all{wenn} + KII).
\item Nous nous référons à votre lettre du 5 mai. (\all{sich beziehen auf} + Acc).
\item Il a été décidé d'augmenter la production. (Passif impersonnel \all{Es wurde beschlossen} + \all{zu} + Inf).
\end{enumerate}
\end{exercice}

\begin{corrige}[Exercice : Thème]
\begin{enumerate}
\item \all{Unsere Firma könnte Ihnen einen \textbf{besseren} Preis anbieten.} (\all{besser} comparatif irrégulier de \all{gut} ; \all{einen besseren Preis} Accusatif).
\item \all{Die Ware wurde gestern \textbf{durch einen Lastwagen} geliefert.} (Passif Präteritum \all{wurde geliefert} ; agent inanimé \all{durch} + Acc).
\item \all{Wenn wir Zeit \textbf{hätten}, \textbf{würden} wir die Fabrik \textbf{besuchen}.} (\all{wenn} + KII \all{hätten}, principale KII \all{würden besuchen}).
\item \all{Wir beziehen uns auf Ihren Brief vom 5. Mai.} (\all{sich beziehen auf} + Accusatif \all{Ihren Brief}).
\item \all{Es wurde beschlossen, die Produktion \textbf{zu steigern}.} (Passif impersonnel Präteritum + \all{zu} + Infinitif \all{steigern} à la fin).
\end{enumerate}
\end{corrige}

\begin{exercice}[Version : Allemand $\rightarrow$ Français]
Traduisez en français naturel (style correspondance/rapport).
\begin{enumerate}
\item \all{Sehr geehrte Damen und Herren, wir möchten 50 Säcke Kaffee bestellen. Könnten Sie die Ware bis Ende Juni liefern?}
\item \all{Am Montag fand eine Besprechung statt. Zuerst wurde der Bericht vorgestellt. Danach diskutierte man über die Kosten. Schließlich einigte man sich auf einen Kompromiss.}
\item \all{Die Rechnung wurde noch nicht bezahlt. Wir bitten Sie, den Betrag innerhalb von 10 Tagen zu überweisen.}
\item \all{Wenn die Qualität besser wäre, würden wir mehr kaufen.}
\item \all{Der Direktor erklärte, dass die Maschinen aus Deutschland importiert worden seien.} (Konjunktiv I dans le discours rapporté -- niveau excellence).
\end{enumerate}
\end{exercice}

\begin{corrige}[Exercice : Version]
\begin{enumerate}
\item « Madame, Monsieur, nous aimerions commander 50 sacs de café. Pourriez-vous livrer la marchandise d'ici fin juin ? »
\item « Lundi a eu lieu une réunion. D'abord, le rapport a été présenté. Ensuite, on a discuté des coûts. Finalement, on s'est mis d'accord sur un compromis. »
\item « La facture n'a pas encore été payée. Nous vous prions de virer le montant dans un délai de 10 jours. »
\item « Si la qualité était meilleure, nous achèterions davantage. »
\item « Le directeur a expliqué que les machines avaient été importées d'Allemagne. » (Discours rapporté : \all{seien} = Konjunktiv I de \all{sind}).
\end{enumerate}
\end{corrige}

\courssub{Entraînement écrit : Lettre et Compte rendu}

\begin{exercice}[Production écrite guidée]
\begin{enumerate}
\item \textbf{Lettre (Bestellung)} : La librairie „Les Classiques“ de Yaoundé (B.P. 123) commande 50 manuels \all{„Deutsch für Tle“} aux éditions Klett (Rotebühlstraße 77, 70178 Stuttgart). Date : Yaoundé, 10. Mai 2025. Utilisez : \all{Sehr geehrte Damen und Herren}, \all{wir möchten ... bestellen}, \all{könnten Sie ... liefern?}, \all{Mit freundlichen Grüßen}. Respectez les 8 zones.
\item \textbf{Compte rendu (Schluss)} : Rédigez la conclusion (\all{Schluss}) d'un compte rendu sur la visite de la banque \all{Afriland First Bank} à Douala (2 phrases minimum). Utilisez \all{schließlich} et \all{deshalb} (avec inversion). Mentionnez les remerciements à la directrice et le souhait de visiter une autre agence l'an prochain.
\end{enumerate}
\end{exercice}

\begin{corrige}[Exercice : Production écrite]
\begin{enumerate}
\item \all{Librairie „Les Classiques“}\\\all{B.P. 123, Yaoundé}\\\all{}\\\all{Klett Verlag\\Rotebühlstraße 77\\70178 Stuttgart\\\\Yaoundé, den 10. Mai 2025\\\\Betreff: Bestellung von Lehrbüchern\\\\Sehr geehrte Damen und Herren,\\wir möchten 50 Exemplare des Lehrbuchs „Deutsch für Tle“ bestellen. Könnten Sie uns die Bücher bis Ende Mai liefern? Vielen Dank im Voraus. Wir freuen uns auf Ihre Antwort.\\\\Mit freundlichen Grüßen\\\\J. Ngo Bell}
\item \all{Schließlich bedankten sich die Schüler bei der Direktorin der Bank für die interessante Führung. Der Besuch war sehr lehrreich; deshalb möchte die Klasse nächstes Jahr eine andere Filiale in Douala besuchen.}
\end{enumerate}
\end{corrige}

\begin{aretenir}[L'essentiel de la leçon AL10 -- Vie économique]
\begin{itemize}
\item \textbf{Lexique} : Maîtriser les termes du commerce (Anfrage, Angebot, Bestellung, Reklamation, Lieferung, Rechnung, Vertrag, Zahlung) et de l'économie (Import, Export, Investition, Gewinn, Arbeitsmarkt).
\item \textbf{Comparatif/Superlatif} : \all{-er / als}, \all{so ... wie}, \all{am ... -sten} ; irréguliers : \all{besser, mehr, lieber, höher, nächster}.
\item \textbf{Konjunktiv II} : Formation (\all{würde} + Inf / forme simple \all{wäre, hätte, könnte, möchte, müsste, sollte, wollte, ginge, käme, stünde, ließe}). Emplois : Politesse (lettre), Irréel (\all{wenn}), Discours rapporté.
\item \textbf{Vorgangspassiv} : \all{werden} + Partizip II (tous temps). \all{von} (agent) / \all{durch} (moyen). Passif impersonnel \all{Es wird ...}.
\item \textbf{Lettre commerciale} : 8 zones fixes. Virgule après Anrede, minuscule au début du corps, PAS de virgule après Grußformel. Formules types par type (Anfrage, Angebot, Bestellung, Reklamation). Konjunktiv II obligatoire.
\item \textbf{Compte rendu} : Structure \all{Einleitung - Hauptteil - Schluss}. Temps : Präteritum (écrit). Style objectif. Connecteurs en position 1 $\rightarrow$ inversion. Passif fréquent.
\item \textbf{Traduction} : Méthode rigoureuse (Analyse $\rightarrow$ Vocab $\rightarrow$ Grammaire $\rightarrow$ Brouillon $\rightarrow$ Relecture). Attention aux faux amis (\all{beziehen auf, freuen auf/über, Bestellung, im Voraus}).
\end{itemize}
\end{aretenir}

\coursec{Traduction : thème et version}

Après avoir appris à rédiger une lettre commerciale et un compte rendu, il nous reste à maîtriser la compétence reine de la Terminale : la \cle{traduction}. Au baccalauréat, on te demandera soit de traduire un texte allemand vers le français (la \cle{version}), soit de traduire des phrases françaises vers l'allemand (le \cle{thème}). Les deux épreuves exigent la même rigueur : comprendre le sens, respecter la grammaire, relire. Cette section te donne une méthode en cinq étapes, des exemples guidés sur la vie économique, et un entraînement au format de l'examen.

\courssub{Les deux directions de la traduction}

\begin{definition}[Thème et version]
La \cle{version} consiste à traduire un texte de l'\textbf{allemand vers le français}. La difficulté majeure réside dans la compréhension de la phrase allemande (verbe en fin de phrase, déclinaisons) et sa reformulation en français naturel. Le \cle{thème} consiste à traduire du \textbf{français vers l'allemand}. La difficulté majeure est de produire une phrase allemande correcte : place du verbe, cas, majuscules, particules séparables.
\end{definition}

\begin{center}
\begin{tabularx}{\linewidth}{|X|X|X|}
\hline
\rowcolor{popblueL} & Version (DE $\rightarrow$ FR) & Thème (FR $\rightarrow$ DE) \\
\hline
Point de départ & Texte allemand & Phrases françaises \\
\hline
Compétence clé & Comprendre et reformuler & Produire une grammaire correcte \\
\hline
Piège majeur & Traduction mot à mot & Calquer l'ordre des mots français \\
\hline
Vérification finale & Le français est-il naturel ? & Verbe en position 2 ? Cas ? Majuscules ? \\
\hline
\end{tabularx}
\end{center}

\courssub{Méthode de la version : de l'allemand vers le français}

Observe d'abord cette phrase tirée de notre texte support :

\all{Der kamerunische Kaffee, der in den westlichen Regionen angebaut wird, erobert seit einigen Jahren den deutschen Markt.}

Un élève pressé risque une traduction mot à mot, produisant un français bancal. Le traducteur averti procède par étapes : il repère le verbe conjugué (\all{erobert}), le sujet (\all{Der kamerunische Kaffee}), il isole la relative (\all{der ... angebaut wird} = « qui est cultivé dans les régions de l'Ouest »), puis il assemble : « Le café camerounais, qui est cultivé dans les régions de l'Ouest, conquiert depuis quelques années le marché allemand. »

\begin{methode}[Traduire une phrase en cinq étapes]
\begin{enumerate}
\item \textbf{Lire} toute la phrase (ou tout le texte) une première fois, sans rien écrire, pour saisir le sens global.
\item \textbf{Trouver le verbe conjugué} : il est en position 2 dans la principale, en \textbf{fin de phrase} dans la subordonnée. C'est le pivot de la phrase.
\item \textbf{Trouver le sujet} : il est au Nominativ. Puis repérer les compléments et leurs cas (Akkusativ = complément direct, Dativ = complément indirect).
\item \textbf{Traduire par le sens}, jamais mot à mot : chercher l'équivalent français naturel de chaque groupe de mots.
\item \textbf{Reformuler et relire} : la phrase française doit être correcte, fluide, et fidèle au sens (temps, négation, chiffres).
\end{enumerate}
\end{methode}

\begin{popfigure}
\begin{center}
\begin{tikzpicture}[node distance=0.55cm,
etape/.style={rectangle, rounded corners=3pt, draw=popblue, fill=popblueL, thick, minimum width=2.1cm, minimum height=0.9cm, align=center, font=\small\bfseries},
fleche/.style={-{Stealth[length=3mm]}, thick, poporange}]
\node[etape] (lire) {1. LIRE};
\node[etape, right=of lire, align=center] (analyser){2. ANALYSER\\ \footnotesize verbe + sujet};
\node[etape, right=of analyser, align=center] (traduire){3. TRADUIRE\\ \footnotesize par le sens};
\node[etape, right=of traduire] (reformuler) {4. REFORMULER};
\node[etape, right=of reformuler] (relire) {5. RELIRE};
\draw[fleche] (lire) -- (analyser);
\draw[fleche] (analyser) -- (traduire);
\draw[fleche] (traduire) -- (reformuler);
\draw[fleche] (reformuler) -- (relire);
\end{tikzpicture}
\end{center}
\legende{Les cinq étapes de la traduction : aucune ne peut être sautée.}
\end{popfigure}

\begin{attention}[Les cadres verbaux : le piège n° 1 de la version]
En allemand, le verbe se « casse » souvent en deux morceaux qui encadrent la phrase : l'auxiliaire en position 2 et le participe ou l'infinitif \textbf{en toute fin}. Ne propose jamais de traduction avant d'avoir lu la phrase jusqu'au dernier mot !

\all{Die Nachfrage nach Kaffee ist in den letzten Jahren größer \textbf{geworden}.} — « La demande de café s'est intensifiée ces dernières années. » (et non « est devenue plus grande » : préfère un français naturel.)

\all{Diese Ware wurde zu einem guten Preis \textbf{verkauft}.} — « Cette marchandise a été vendue à un bon prix. »
\end{attention}

\courssub{Version guidée : extraits du texte support}

Reprenons des extraits du texte « \all{Der kamerunische Kaffee erobert den deutschen Markt} » étudié dans la première section, et traduisons-les en appliquant la méthode.

\begin{exemplebox}[Version guidée, phrase 1]
\all{In Kamerun arbeiten viele Bauern auf kleinen Kaffeeplantagen.}

\textbf{Analyse} : verbe conjugué \all{arbeiten} (position 2) ; le sujet est \all{viele Bauern} (Nominativ pluriel) ; \all{In Kamerun} et \all{auf kleinen Kaffeeplantagen} (Dativ après \all{auf} pour le lieu) sont des compléments de lieu.

\textbf{Traduction} : « Au Cameroun, de nombreux paysans travaillent sur de petites plantations de café. »
\end{exemplebox}

\begin{exemplebox}[Version guidée, phrase 2]
\all{Weil die Qualität des Kaffees besser geworden ist, kaufen deutsche Firmen jetzt mehr Bohnen.}

\textbf{Analyse} : subordonnée avec \all{weil} : le verbe conjugué \all{ist} est en fin, le participe \all{geworden} juste avant (« parce que la qualité du café est devenue meilleure » = « s'est améliorée »). Principale : verbe \all{kaufen} en position 2, sujet \all{deutsche Firmen}.

\textbf{Traduction} : « Parce que la qualité du café s'est améliorée, les entreprises allemandes achètent maintenant davantage de grains. »
\end{exemplebox}

\begin{exemplebox}[Version guidée, phrase 3]
\all{Der Kaffee wird von Händlern in den Hafen von Douala gebracht und von dort nach Deutschland exportiert.}

\textbf{Analyse} : deux passifs (\all{wird ... gebracht}, \all{exportiert}) ; \all{von Händlern} = complément d'agent (« par des négociants ») ; \all{in den Hafen} (Akkusativ : direction) vs \all{von dort} (origine).

\textbf{Traduction} : « Le café est transporté par des négociants jusqu'au port de Douala, puis exporté de là vers l'Allemagne. »
\end{exemplebox}

\courssub{Méthode du thème : du français vers l'allemand}

Pour le thème, le raisonnement est inversé : tu pars du français et tu dois \textbf{construire} la phrase allemande. Trois questions à te poser dans l'ordre :

\begin{enumerate}
\item \textbf{Quel temps ?} Le français « a vendu » peut être un Perfekt (\all{hat verkauft}, langue parlée) ou un Präteritum (\all{verkaufte}, langue écrite). Le conditionnel français se rend par le Konjunktiv II.
\item \textbf{Quelle structure ?} Principale : verbe conjugué en position 2. Subordonnée (\all{weil}, \all{dass}, \all{wenn}, \all{obwohl}) : verbe conjugué en fin.
\item \textbf{Quels cas ?} Sujet au Nominativ, complément direct à l'Akkusativ, complément indirect au Dativ. Et n'oublie pas la majuscule à tous les noms !
\end{enumerate}

\begin{center}
\begin{tabularx}{\linewidth}{|X|X|X|}
\hline
\rowcolor{popblueL} Français & Deutsch & Remarque \\
\hline
L'entreprise vend du café. & \all{Die Firma verkauft Kaffee.} & Verbe en position 2. \\
\hline
Aujourd'hui, l'entreprise vend du café. & \all{Heute verkauft die Firma Kaffee.} & Le verbe reste en position 2 : le sujet passe après. \\
\hline
Je pense que le prix est trop élevé. & \all{Ich denke, dass der Preis zu hoch ist.} & Subordonnée : verbe \all{ist} en fin. \\
\hline
Parce que le café est cher, j'achète du thé. & \all{Weil der Kaffee teuer ist, kaufe ich Tee.} & Subordonnée en tête : le verbe de la principale suit immédiatement la virgule. \\
\hline
\end{tabularx}
\end{center}

\begin{attention}[Ne calque jamais l'ordre des mots français !]
C'est l'erreur numéro un des francophones. « Aujourd'hui l'entreprise vend » ne se traduit PAS par \all{Heute die Firma verkauft} mais par \all{Heute verkauft die Firma}. En allemand, le verbe conjugué est \textbf{toujours} le deuxième élément de la phrase déclarative, quoi qu'il y ait devant. De même, « parce que le café est cher » devient \all{weil der Kaffee teuer \textbf{ist}}, avec le verbe à la fin.
\end{attention}

\courssub{Thème guidé : cinq phrases économiques}

\begin{exemplebox}[Thème guidé]
\begin{enumerate}
\item « Si l'entreprise avait plus de clients, elle embaucherait plus d'ouvriers. » $\rightarrow$ \all{Wenn die Firma mehr Kunden hätte, würde sie mehr Arbeiter einstellen.} (Konjunktiv II : \all{hätte} en fin de subordonnée ; \all{würde einstellen}, participe séparable \all{ein-} en fin.)
\item « Le cacao est exporté par de grandes entreprises. » $\rightarrow$ \all{Der Kakao wird von großen Unternehmen exportiert.} (Passif : \all{wird} + participe en fin ; agent introduit par \all{von} + Dativ.)
\item « Le café camerounais est meilleur que le café robusta ordinaire. » $\rightarrow$ \all{Der kamerunische Kaffee ist besser als der gewöhnliche Robusta-Kaffee.} (Comparatif : \all{besser als}.)
\item « Les prix ont augmenté parce que la demande est devenue plus grande. » $\rightarrow$ \all{Die Preise sind gestiegen, weil die Nachfrage größer geworden ist.} (Perfekt \all{sind gestiegen} ; subordonnée \all{weil} avec verbe en fin.)
\item « Si j'avais plus d'argent, j'ouvrirais un magasin à Yaoundé. » $\rightarrow$ \all{Wenn ich mehr Geld hätte, würde ich ein Geschäft in Yaoundé eröffnen.} (Konjunktiv II des deux côtés.)
\end{enumerate}
\end{exemplebox}

\courssub{Pièges de traduction et faux amis}

\begin{attention}[Le passif allemand rendu par « on »]
Le passif impersonnel allemand se traduit très souvent en français par « on » + verbe actif : \all{Hier wird viel gearbeitet.} — « Ici, on travaille beaucoup. » \all{In Deutschland wird viel Kaffee getrunken.} — « En Allemagne, on boit beaucoup de café. » Inversement, au thème, le « on » français appelle naturellement un passif allemand : « On exporte le cacao » $\rightarrow$ \all{Der Kakao wird exportiert.}
\end{attention}

\begin{center}
\begin{tabularx}{\linewidth}{|X|X|X|}
\hline
\rowcolor{popblueL} Faux ami ou piège & Sens réel & Exemple \\
\hline
\all{aktual} & actuel (pas « actuel » au sens de « réel ») & \all{die aktuellen Preise} = les prix actuels \\
\hline
\all{die Fabrik} & l'usine (pas « la fabrique » au sens abstrait) & \all{Er arbeitet in einer Fabrik.} \\
\hline
\all{bekommen} & recevoir (pas « devenir » = \all{werden}) & \all{Die Firma bekommt neue Aufträge.} = L'entreprise reçoit de nouvelles commandes. \\
\hline
\all{die Mappe} & la chemise, le porte-documents (pas « la mappe ») & \all{die Bewerbungsmappe} = le dossier de candidature \\
\hline
\all{einstellen} & embaucher (particule séparable !) & \all{Die Firma stellt Arbeiter ein.} \\
\hline
\end{tabularx}
\end{center}

\begin{attention}[Particules séparables au thème]
Quand tu traduis « L'entreprise embauche des ouvriers », n'oublie pas que \all{einstellen} est séparable : \all{Die Firma \textbf{stellt} Arbeiter \textbf{ein}.} La particule part en fin de phrase. Au Perfekt : \all{Die Firma hat Arbeiter \textbf{eingestellt}.}
\end{attention}

\courssub{Exercice résolu}

\begin{exoresolu}[Version : un extrait économique]
Traduis en français : \all{Seit zehn Jahren wird der kamerunische Kakao immer bekannter. Viele europäische Firmen kaufen ihn, weil seine Qualität ausgezeichnet ist. Wenn die Bauern besser bezahlt würden, könnten sie mehr produzieren.}
\tcblower
\textbf{Solution détaillée.}

Phrase 1 : verbe conjugué \all{wird} (position 2), \all{bekannter} est un comparatif adjectival (pas un participe) : \all{werden} + comparatif = « devenir de plus en plus... ». Traduction : « Depuis dix ans, le cacao camerounais est de plus en plus connu. »

Phrase 2 : verbe \all{kaufen} (position 2), sujet \all{Viele europäische Firmen} ; subordonnée \all{weil} avec verbe \all{ist} en fin. Traduction : « De nombreuses entreprises européennes l'achètent parce que sa qualité est excellente. »

Phrase 3 : subordonnée \all{wenn} au Konjunktiv II passif (\all{bezahlt würden}, verbe en fin), principale au Konjunktiv II (\all{könnten ... produzieren}). Traduction : « Si les paysans étaient mieux payés, ils pourraient produire davantage. »
\end{exoresolu}

\courssub{La grille de relecture}

Avant de rendre ta copie, passe chaque phrase au crible de cette grille. Elle vaut pour le thème comme pour la version.

\begin{aretenir}[Grille de relecture en six points]
\begin{enumerate}
\item \textbf{Place du verbe} : position 2 dans la principale, fin de phrase dans la subordonnée.
\item \textbf{Majuscules} : tous les noms allemands prennent une majuscule (\all{der Markt, die Firma, das Geld}).
\item \textbf{Cas} : sujet au Nominativ, complément direct à l'Akkusativ, Dativ après \all{mit, von, zu, bei, nach}.
\item \textbf{Temps} : le temps allemand correspond-il au temps français (conditionnel = Konjunktiv II) ?
\item \textbf{Orthographe} : Umlaute (\all{ä, ö, ü}), \all{ß}, particules séparables détachées.
\item \textbf{Sens} : la traduction est-elle fidèle (négations, chiffres, « plus/moins ») et le français est-il naturel ?
\end{enumerate}
\end{aretenir}

\begin{savaistu}[Comment la traduction est-elle notée au Bac ?]
Au baccalauréat, l'épreuve de traduction est notée sur deux critères indissociables : la \textbf{fidélité au sens} (as-tu compris et rendu toutes les informations ?) et la \textbf{correction grammaticale} (place du verbe, cas, temps, orthographe). Une erreur de place du verbe coûte des points même si le sens global est compris ! C'est pourquoi la relecture grammaticale n'est pas un luxe : c'est une étape de l'épreuve à part entière. Prévois toujours trois à quatre minutes pour relire ta traduction.
\end{savaistu}

\courssub{Entraînement au format Bac}

\begin{textebox}[Texte de version — 20 minutes]
\all{Der Handel zwischen Kamerun und Deutschland ist in den letzten Jahren stärker geworden. Kamerun exportiert vor allem Kakao, Kaffee und Bananen, während Deutschland Maschinen und Autos liefert. Viele kamerunische Produkte werden im Hafen von Douala auf Schiffe verladen und brauchen dann etwa drei Wochen bis Hamburg. Wenn der Hafen moderner wäre, könnten die Waren schneller transportiert werden. Einige Experten meinen, dass Kamerun seine Produkte mehr verarbeiten sollte, denn verarbeitete Waren werden zu höheren Preisen verkauft. So könnten mehr Arbeitsplätze im Land geschaffen werden.}

\textbf{Glossaire} : der Handel = le commerce ; während = tandis que ; die Maschine, -n = la machine ; liefern = livrer ; verladen (verlädt, verlud, hat verladen) = charger ; etwa = environ ; verarbeiten = transformer ; der Arbeitsplatz, -plätze = l'emploi ; schaffen = créer.
\end{textebox}

\begin{exercice}[Version — texte ci-dessus]
Traduis le texte « \all{Der Handel zwischen Kamerun und Deutschland} » en français soigné. Temps conseillé : 20 minutes. Applique les cinq étapes et termine par la grille de relecture.
\end{exercice}

\begin{corrige}[Version — texte « Der Handel »]
« Le commerce entre le Cameroun et l'Allemagne s'est intensifié ces dernières années. Le Cameroun exporte surtout du cacao, du café et des bananes, tandis que l'Allemagne livre des machines et des voitures. De nombreux produits camerounais sont chargés sur des navires dans le port de Douala et mettent ensuite environ trois semaines jusqu'à Hambourg. Si le port était plus moderne, les marchandises pourraient être transportées plus rapidement. Certains experts estiment que le Cameroun devrait davantage transformer ses produits, car les marchandises transformées sont vendues à des prix plus élevés. Ainsi, davantage d'emplois pourraient être créés dans le pays. »

Points de vigilance : \all{ist ... stärker geworden} = « s'est intensifié » (pas « est devenu plus fort ») ; \all{brauchen ... drei Wochen} = « mettent trois semaines » ; \all{während} introduit une subordonnée (verbe \all{liefert} en fin) ; \all{verladen}, \all{verkauft}, \all{geschaffen} sont des passifs à repérer jusqu'en fin de phrase.
\end{corrige}

\begin{exercice}[Thème — cinq phrases]
Traduis en allemand :
\begin{enumerate}
\item Le marché de Douala est plus grand que le marché de Bafoussam.
\item Si les prix baissaient, les clients achèteraient plus de marchandises.
\item Le café est cultivé par de petits paysans dans l'Ouest du pays.
\item On travaille beaucoup dans cette entreprise.
\item Je crois que le cacao camerounais est le meilleur du monde.
\end{enumerate}
\end{exercice}

\begin{corrige}[Thème — cinq phrases]
\begin{enumerate}
\item \all{Der Markt von Douala ist größer als der Markt von Bafoussam.} (comparatif \all{größer als}, Umlaut sur \all{größer}.)
\item \all{Wenn die Preise sänken, würden die Kunden mehr Waren kaufen.} (Konjunktiv II : \all{sänken} forme préteritale du verbe fort \all{sinken} ; \all{würden kaufen}.)
\item \all{Der Kaffee wird von kleinen Bauern im Westen des Landes angebaut.} (passif ; particule séparable \all{an-} soudée au participe : \all{angebaut}.)
\item \all{In dieser Firma wird viel gearbeitet.} (le « on » français rendu par le passif impersonnel allemand.)
\item \all{Ich glaube, dass der kamerunische Kakao der beste der Welt ist.} (subordonnée \all{dass} : verbe \all{ist} en fin ; superlatif \all{der beste}.)
\end{enumerate}
\end{corrige}

\begin{experience}[Atelier de correction collective]
Par groupes de trois : chacun traduit le texte « \all{Der Handel} » seul en 20 minutes, puis le groupe compare les trois versions. Pour chaque phrase, désignez la meilleure formulation et justifiez votre choix avec la grille de relecture (place du verbe, cas, fidélité au sens). Le groupe présente ensuite sa version commune à la classe.
\end{experience}

\begin{aretenir}[L'essentiel de la section]
La traduction réussie repose sur une méthode en cinq étapes : lire, analyser (verbe, sujet, cas), traduire par le sens, reformuler, relire. En version, lis toujours la phrase allemande jusqu'au bout pour attraper le verbe en fin. En thème, impose-toi les trois réflexes allemands : verbe en position 2 (ou en fin de subordonnée), majuscules aux noms, cas corrects. Le passif allemand dialogue avec le « on » français, et le conditionnel français appelle le Konjunktiv II. Enfin, la relecture n'est pas optionnelle : au Bac, chaque erreur de grammaire coûte des points, même quand le sens est bon.
\end{aretenir}

\newpage

\coursec{Activité d'intégration}

\coursec{Leçon AL10 : Production et traduction de textes liés à la vie économique}

\courssub{Objectifs de l'activité}

À l'issue de cette activité d'intégration, tu seras capable de :

\begin{itemize}
\item rédiger une lettre commerciale complète en allemand (en-tête, formule d'appel, corps, formule de politesse) à partir d'une situation de communication authentique ;
\item réinvestir le lexique de la vie économique (l'exportation, les prix, la livraison, la commande) dans un texte cohérent ;
\item employer correctement au moins un \cle{comparatif}, un \cle{superlatif}, un \cle{conditionnel de politesse} (Konjunktiv II) et une phrase au \cle{passif} (Vorgangspassiv) ;
\item traduire une lettre commerciale de l'allemand vers le français (version) et du français vers l'allemand (thème) ;
\item relire et corriger une production écrite à l'aide d'une grille (orthographe, majuscules, ordre des mots, cas, place du verbe).
\end{itemize}

\courssub{Situation}

Tu es directeur ou directrice commercial(e) de la coopérative \all{Kamerun Gold Kakao}, installée à Obala, dans la région du Centre. Ta coopérative regroupe 240 planteurs et exporte du cacao, du café arabica et des ananas. La semaine dernière, tu as participé au salon international de l'alimentation à Douala. Un importateur allemand de Hambourg, la société \all{Nordhandel GmbH}, a visité ton stand et t'a laissé un message électronique. Le voici :

\begin{textebox}[E-Mail des Importeurs — Nordhandel GmbH, Hamburg]
Sehr geehrte Damen und Herren,\par
wir haben Ihren Stand auf der Messe in Douala besucht und waren sehr beeindruckt von der Qualität Ihres Kakaos. Wir möchten gern Kakao und eventuell auch Ananas aus Kamerun importieren.\par
Könnten Sie uns bitte mehr Informationen schicken? Wir interessieren uns für Ihre Preise, die Lieferbedingungen und die Mindestmenge pro Bestellung. Wird der Kakao von Ihrer Kooperative direkt exportiert, oder arbeiten Sie mit einem Zwischenhändler?\par
Wir würden uns über ein Angebot freuen.\par
Mit freundlichen Grüßen\par
Jens Berger\par
Nordhandel GmbH, Hamburg
\end{textebox}

\textbf{Glossaire :} \all{die Messe, -n} : le salon, la foire ; \all{beeindruckt} : impressionné ; \all{die Lieferbedingungen (pl.)} : les conditions de livraison ; \all{die Mindestmenge, -n} : la quantité minimale ; \all{die Bestellung, -en} : la commande ; \all{der Zwischenhändler, -} : l'intermédiaire ; \all{das Angebot, -e} : l'offre.

En binôme, tu réponds à ce message par une lettre commerciale formelle de 15 à 20 lignes. Puis tu échanges ta lettre avec un autre groupe, qui la traduit en français et la corrige.

\courssub{Consignes}

\textbf{Tâche 1 — Compréhension du message.} Lis le courriel de M. Berger et relève, par écrit, les trois questions précises auxquelles ta lettre devra répondre. Souligne dans le texte les mots du champ lexical de l'économie.

\textbf{Tâche 2 — Préparation du lexique.} Avec ton binôme, construis une liste de dix mots ou expressions utiles pour la réponse (par exemple : \all{das Angebot}, \all{der Preis}, \all{liefern}, \all{der Hafen}, \all{die Qualität}, \all{exportieren}, \all{die Zahlung}). Vérifie l'article et le pluriel de chaque nom dans le lexique de la leçon.

\textbf{Tâche 3 — Plan de la lettre.} Rédige le plan de ta lettre en respectant la structure de la lettre commerciale allemande : en-tête avec expéditeur et destinataire, lieu et date, objet (\all{Betreff}), formule d'appel, corps en trois paragraphes (remerciement et présentation ; réponses aux questions ; proposition et conclusion), formule de politesse, signature.

\textbf{Tâche 4 — Rédaction.} Écris ta lettre (15 à 20 lignes). Elle doit contenir obligatoirement :
\begin{enumerate}
\item au moins un \cle{comparatif} (par exemple : notre cacao est de meilleure qualité que...) ;
\item au moins un \cle{superlatif} (par exemple : \all{unser Kakao ist der beste}) ;
\item au moins un \cle{conditionnel de politesse} (par exemple : \all{Wir würden uns freuen...}, \all{Könnten Sie...}) ;
\item au moins une phrase au \cle{passif} (par exemple : le cacao est exporté directement par notre coopérative) ;
\item au moins cinq mots du lexique économique du chapitre.
\end{enumerate}

\textbf{Tâche 5 — Échange et version.} Échange ta lettre avec un autre binôme. Traduis la lettre reçue en français (version) sur ta copie, en respectant le sens sans traduire mot à mot.

\textbf{Tâche 6 — Relecture et correction.} Corrige la lettre que tu as traduite à l'aide de la grille de relecture ci-dessous. Note chaque erreur dans la marge et propose une correction.

\begin{center}
\begin{tabularx}{\linewidth}{|X|X|}\hline
\rowcolor{popblueL} \textbf{Point de contrôle} & \textbf{À vérifier} \\ \hline
Majuscules & Tous les noms prennent une majuscule. \\ \hline
Ordre des mots & Verbe conjugué en 2\up{e} position ; verbe à la fin dans les subordonnées. \\ \hline
Cas & Datif après \all{mit}, \all{von}, \all{zu}, \all{bei} ; accusatif après \all{für}, \all{über}, \all{ohne}. \\ \hline
Konjunktiv II & \all{würden} + infinitif ; \all{könnten}, \all{hätten}, \all{wären} correctement formés. \\ \hline
Passif & \all{werden} + Partizip II ; \all{von} + datif pour l'agent. \\ \hline
Formules & \all{Sehr geehrter Herr Berger,} ... \all{Mit freundlichen Grüßen} \\ \hline
\end{tabularx}
\end{center}

\textbf{Tâche 7 — Restitution orale.} Les deux meilleures lettres sont lues à voix haute devant la classe. La classe vote pour la lettre la plus convaincante du point de vue commercial.

\courssub{Ressources à mobiliser}

\begin{aretenir}[Rappels grammaticaux : Comparatif et Superlatif]
\textbf{Comparatif :} adjectif/adverbe + \all{-er} (+ \all{als}) : \all{unser Kakao ist besser als...} ; \all{wir liefern schneller als...}.\par
\textbf{Formes irrégulières :} \all{gut} $\rightarrow$ \all{besser} ; \all{viel} $\rightarrow$ \all{mehr} ; \all{gern} $\rightarrow$ \all{lieber} ; \all{hoch} $\rightarrow$ \all{höher} ; \all{nah} $\rightarrow$ \all{näher}.\par
\textbf{Superlatif :} \all{am} + adjectif + \all{-sten} (prédicatif) : \all{unser Kakao ist am besten} ; article défini + adjectif + \all{-ste} (attribut) : \all{der beste Kakao}. \textbf{Attention :} \all{am meisten}, \all{am liebsten}, \all{am höchsten}.
\end{aretenir}

\begin{aretenir}[Rappels grammaticaux : Conditionnel de politesse (Konjunktiv II)]
\textbf{Forme générale :} \all{würden} (conjugué) + infinitif à la fin de la phrase : \all{Wir würden uns über eine Zusammenarbeit freuen.} \all{Könnten Sie uns bitte mitteilen...?}\par
\textbf{Formes spéciales (fréquentes à l'écrit formel) :} \all{ich könnte, du könntest, er könnte / wir könnten} ; \all{ich hätte, wir hätten} ; \all{ich wäre, wir wären} ; \all{es gäbe}.\par
\textbf{Emploi :} politesse, hypothèse, souhait : \all{Wir hätten gern ein Angebot.} \all{Wäre das möglich?}
\end{aretenir}

\begin{aretenir}[Rappels grammaticaux : Passif (Vorgangspassiv)]
\textbf{Structure :} Sujet + \all{werden} (conjugué au temps voulu) + \all{Partizip II} à la fin.\par
\textbf{Présent :} \all{Der Kakao wird exportiert.} \textbf{Prétérit :} \all{Der Kakao wurde exportiert.} \textbf{Parfait :} \all{Der Kakao ist exportiert worden.}\par
\textbf{Agent :} introduit par \all{von} + \textbf{Datif} : \all{von unserer Kooperative} ; \all{von den Bauern}.\par
\textbf{Passif sans agent (impersonnel) :} \all{Hier wird Kakao angebaut.} (Sujet \all{es} élidé).
\end{aretenir}

\begin{aretenir}[Structure de la lettre commerciale (DIN 5008)]
\begin{enumerate}
\item \textbf{Expéditeur} (haut gauche) : Nom, adresse, téléphone, e-mail.
\item \textbf{Destinataire} (sous l'expéditeur) : Nom de l'entreprise, \all{z. Hd.} (\all{zu Händen} = à l'attention de) + \textbf{Prénom Nom} (accusatif : \all{Herrn Jens Berger}), adresse.
\item \textbf{Lieu et date} (alignés à droite) : \all{Obala, den 12. März 2025} (\all{den} = accusatif de temps).
\item \textbf{Objet} : \all{Betreff: Ihre Anfrage — Angebot Kakao und Ananas} (sans deux-points après \all{Betreff} selon la norme actuelle, mais toléré).
\item \textbf{Formule d'appel} : \all{Sehr geehrter Herr Berger,} (virgule obligatoire, ligne suivante en \textbf{minuscule}).
\item \textbf{Corps de la lettre} : 3 paragraphes standards.
\item \textbf{Formule de politesse} : \all{Mit freundlichen Grüßen} (\textbf{sans virgule}).
\item \textbf{Signature} : Prénom Nom, fonction.
\item \textbf{Pièces jointes} : \all{Anlage(n)} (en bas à gauche).
\end{enumerate}
\end{aretenir}

\begin{aretenir}[Structure d'un compte rendu (Protokoll / Geschäftsbericht)]
Un \cle{compte rendu} (procès-verbal, rapport d'activité) suit une structure fixe :
\begin{enumerate}
\item \textbf{En-tête :} \all{Protokoll über...} / \all{Bericht über...}, \all{Datum}, \all{Ort}, \all{Teilnehmer} (liste), \all{Verteiler}.
\item \textbf{Corps :} \all{1. Tagesordnungspunkte} (points à l'ordre du jour), \all{2. Beschlüsse} (décisions : \all{Es wurde beschlossen, dass...} + passif), \all{3. Maßnahmen / To-Dos} (actions : \all{Herr X kümmert sich um...} / \all{Die Lieferung wird bis zum... organisiert}).
\item \textbf{Pied :} \all{Unterschrift} (rédacteur + président), \all{Ort, Datum}.
\end{enumerate}
\textbf{Style :} impersonnel, passif fréquent (\all{wurde entschieden, wird erledigt}), présent ou prétérit, vocabulaire précis (\all{die Frist}, \all{die Zuständigkeit}, \all{die Rückmeldung}).
\end{aretenir}

\begin{definition}[Lexique utile pour la lettre et le compte rendu]
\all{das Angebot, -e} : l'offre ; \all{der Preis, -e} : le prix ; \all{die Qualität, -en} : la qualité ; \all{die Lieferung, -en} : la livraison ; \all{die Bestellung, -en} : la commande ; \all{der Hafen, -Häfen} : le port ; \all{die Zahlung, -en} : le paiement ; \all{der Kunde, -n / die Kundin, -nen} : le/la client(e) ; \all{der Lieferant, -en} : le fournisseur ; \all{liefern} (reg.) : livrer ; \all{exportieren} (reg.) : exporter ; \all{importieren} (reg.) : importer ; \all{sich freuen über} + Akk. : se réjouir de ; \all{die Zusammenarbeit} : la collaboration ; \all{die Anfrage, -n} : la demande de renseignements ; \all{die Mindestmenge, -n} : la quantité minimale ; \all{die Anlage, -n} : la pièce jointe ; \all{der Vertrag, -\"{e}r} : le contrat ; \all{die Frist, -en} : le délai ; \all{verbindlich} : ferme, contraignant.
\end{definition}

\begin{attention}[Pièges fréquents des francophones]
\begin{itemize}
\item \textbf{Majuscules :} \all{Sehr geehrter Herr Berger,} (virgule) $\rightarrow$ ligne suivante \all{vielen Dank...} (minuscule). Jamais \all{Vielen Dank}.
\item \textbf{Faux ami « reconnaissant » :} Ne traduis pas « nous vous serions reconnaissants » par \all{*wir wären Ihnen erkennend}. Correct : \all{Wir wären Ihnen dankbar, wenn...} ou \all{Wir würden es schätzen, wenn...}.
\item \textbf{Passif :} \all{werden} se conjugue (2\up{e} position), Partizip II à la fin. \all{Die Ware wird nächste Woche geliefert.} (Pas \all{*wird geliefert die Ware}).
\item \textbf{Comparatif :} \all{als} (pas \all{wie}) pour l'inégalité. \all{besser als}, \all{teurer als}.
\item \textbf{Ordre des mots (subordonnée) :} \all{..., weil er von Bauern \textbf{geerntet wird}.} (verbe tout à la fin).
\item \textbf{Prépositions :} \all{mit} + Datif (\all{mit dem Hafen}), \all{für} + Accusatif (\all{für den Export}), \all{über} + Accusatif (\all{über den Preis}).
\item \textbf{ß vs ss :} \all{Straße, groß, weiß} (après voyelle longue/diphtongue) ; \all{dass, muss, Fluss} (après voyelle brève).
\end{itemize}
\end{attention}

\courssub{Proposition de production et corrigé commenté}

\begin{exoresolu}[Lettre commerciale modèle]
\textbf{Production modèle :}\par
Kamerun Gold Kakao, Genossenschaft\par
Postfach 124, Obala, Kamerun\par
Tel.: +237 2 22 00 00 00, E-Mail: export@kamgold.kr\par
\par
Nordhandel GmbH\par
z. Hd. Herrn Jens Berger\par
Ballindamm 15, 20095 Hamburg, Deutschland\par
\par
Obala, den 12. März 2025\par
Betreff: Ihre Anfrage — Angebot Kakao und Ananas\par
Sehr geehrter Herr Berger,\par
vielen Dank für Ihre Anfrage und für Ihr Interesse an unseren Produkten. Unsere Kooperative vereint 240 Bauern aus der Region Obala und produziert Kakao, Kaffee und Ananas von höchster Qualität.\par
Unser Kakao ist besser als viele andere Produkte auf dem Markt, weil er von unseren Bauern sorgfältig getrocknet und fermentiert wird. Er wird direkt von unserer Kooperative exportiert, ohne Zwischenhändler. Deshalb können wir Ihnen günstigere Preise anbieten als andere Lieferanten. Die Mindestmenge pro Bestellung beträgt eine Tonne. Die Lieferung erfolgt über den Hafen von Douala; die Zahlung per Akkreditiv ist für uns am sichersten.\par
Könnten Sie uns bitte mitteilen, welche Menge Sie pro Monat benötigen? Wir würden uns sehr über eine langfristige Zusammenarbeit mit Ihrer Firma freuen. In der Anlage finden Sie unser detailliertes Angebot mit allen Preisen und Lieferbedingungen.\par
Mit freundlichen Grüßen\par
Amina Mbarga\par
Handelsdirektorin
\tcblower
\textbf{Commentaire des choix de langue :}\par
\textbf{Structure :} la lettre respecte la norme DIN 5008 : expéditeur complet, destinataire avec \all{z. Hd.} + accusatif (\all{Herrn}), lieu/date avec \all{den} + accusatif, \all{Betreff} clair, appel avec virgule et minuscule suivante, trois paragraphes distincts, formule finale sans virgule, signature avec fonction.\par
\textbf{Comparatif :} \all{besser als} (irrégulier \all{gut $\rightarrow$ besser}) ; \all{günstigere Preise ... als andere Lieferanten} (adjectif attribut \all{günstiger} + désinence \all{-e} pour nominatif/accusatif pluriel sans article).\par
\textbf{Superlatif :} \all{von höchster Qualität} (génitif de qualité : \all{höchster} = \all{hoch} superlatif attributif féminin génitif) ; \all{am sichersten} (prédicatif \all{am + adj. + -sten}, irrégulier \all{sicher $\rightarrow$ am sichersten}).\par
\textbf{Conditionnel de politesse (Konjunktiv II) :} \all{Könnten Sie uns bitte mitteilen...} (verbe modal \all{könnten} en position 1 question, infinitif \all{mitteilen} à la fin) ; \all{Wir würden uns ... freuen} (\all{würden} position 2, infinitif \all{freuen} à la fin).\par
\textbf{Passif (Vorgangspassiv) :} \all{er wird ... getrocknet und fermentiert} (présent \all{wird} + Partizip II) ; \all{Er wird direkt ... exportiert} ; agent \all{von unserer Kooperative} (Datif féminin), \all{von unseren Bauern} (Datif pluriel \all{-n}).\par
\textbf{Lexique économique (9 occurrences) :} \all{die Anfrage}, \all{das Angebot}, \all{die Mindestmenge}, \all{die Bestellung}, \all{die Lieferung}, \all{der Hafen}, \all{der Lieferant}, \all{die Zusammenarbeit}, \all{die Anlage}, \all{die Zahlung}, \all{das Akkreditiv}, \all{die Lieferbedingungen}.\par
\textbf{Traduction (version) du dernier paragraphe :} « Pourriez-vous nous indiquer de quelle quantité vous avez besoin par mois ? Nous serions très heureux d'une collaboration à long terme avec votre entreprise. Vous trouverez en pièce jointe notre offre détaillée avec tous les prix et conditions de livraison. »
\end{exoresolu}

\courssub{Bilan de l'activité}

Cette activité t'a permis de mobiliser en situation réelle l'ensemble des acquis du chapitre :

\begin{itemize}
\item \textbf{Compréhension de l'écrit :} tu as identifié les attentes précises d'un correspondant professionnel germanophone (prix, conditions, quantité, circuit d'exportation) ;
\item \textbf{Lexique :} tu as réemployé le vocabulaire de la vie économique (offre, prix, livraison, commande, exportation, paiement, délai) dans un contexte camerounais authentique (cacao, Obala, Douala) ;
\item \textbf{Grammaire :} tu as produit un comparatif, un superlatif, un conditionnel de politesse et une phrase au passif, quatre structures essentielles de la correspondance commerciale ;
\item \textbf{Traduction :} en traduisant la lettre d'un autre groupe, tu as pratiqué la version et vérifié ta compréhension fine de l'allemand (passif, subordonnées, vocabulaire spécialisé) ;
\item \textbf{Relecture :} la grille de correction t'a rendu autonome dans la vérification des majuscules, de l'ordre des mots, des cas et des temps — compétence décisive à l'examen du Baccalauréat.
\end{itemize}

\begin{methode}[Pour aller plus loin]
Garde ta lettre corrigée dans ton portfolio : elle constitue un modèle réutilisable. Pour t'entraîner, remplace le cacao par le café ou l'ananas, change le destinataire (un importateur autrichien \all{Alpenfrucht GmbH, Wien} ou suisse \all{Tropical Trade AG, Zürich}) et adapte les prix, les quantités et le port d'embarquement (Douala, Kribi) : la structure de la lettre reste identique, seul le contenu varie. Entraîne-toi aussi à rédiger un \cle{compte rendu} (\all{Protokoll}) de la rencontre au salon de Douala en utilisant le passif et le style impersonnel.
\end{methode}

\newpage

\coursec{Méthodes et exercices résolus}

\courssub{Méthodes et exercices résolus}

\begin{methode}[Analyser un texte économique : lecture active et prise de notes]
\textbf{Objectif :} Comprendre un article de presse ou un rapport économique (vocabulaire spécialisé, chiffres, tendances) et en restituer les idées principales en allemand ou en français.
\begin{enumerate}
\item \textbf{Repérer la source et le type de texte :} Titre, auteur, date, journal (z.\,B. \all{Handelsblatt}, \all{Wirtschaftswoche}, \all{Le Monde} version allemande). Est-ce un article d'actualité, une analyse, un communiqué ?
\item \textbf{Lecture globale (Skimming) :} Identifier le \cle{sujet principal} (Thema), les \cle{acteurs} (Unternehmen, Staat, Verbraucher), le \cle{secteur} (Branche : Automobil, IT, Landwirtschaft, Tourismus).
\item \textbf{Lecture détaillée (Scanning) :} Souligner les \cle{mots-clés} (Wachstum, Inflation, Export, Import, Arbeitslosigkeit, Investition, Zinssatz). Noter les \cle{chiffres clés} (pourcentages, milliards, dates).
\item \textbf{Structurer l'information :} Utiliser un tableau « Qui ? Quoi ? Où ? Quand ? Combien ? Pourquoi ? » ou une carte mentale (Ursache $\rightarrow$ Wirkung).
\item \textbf{Rédaction de la synthèse :} En allemand (Zusammenfassung) ou en français. Utiliser le \cle{Präteritum} (récit journalistique) ou le \cle{Perfekt} (résultat actuel). Phrases courtes, style objectif, pas de « je ».
\end{enumerate}
\end{methode}

\begin{exoresolu}[Compréhension de texte : La croissance camerounaise et les investissements allemands]
\begin{textebox}[Article simulé \textit{WirtschaftsWoche Afrika}, mars 2025]
Deutschlands Wirtschaft zeigt wachsendes Interesse an Kamerun. Laut einer Studie der KfW-Bank stiegen die deutschen Direktinvestitionen (Foreign Direct Investment) im Jahr 2024 um 12\,\% im Vergleich zum Vorjahr. Besonders der Energiesektor (Wasserkraft, Solarenergie) und die Agrarindustrie (Kakao-Verarbeitung, Holz) profitieren. Der kamerunische Finanzminister betonte: „Wir brauchen Technologietransfer, nicht nur Kapital.“ Die Arbeitslosenquote junger Akademiker bleibt jedoch hoch, da die neuen Fabriken oft hochautomatisiert sind. Ein Freihandelsabkommen mit der EU (EPA) soll den Export kamerunischer Produkte nach Europa erleichtern, doch bürokratische Hürden und Korruption bremsen den Handel noch immer.
\end{textebox}
\textbf{Questions :}
\begin{enumerate}
\item Um wie viel Prozent stiegen die deutschen Direktinvestitionen 2024?
\item Nennen Sie zwei Sektoren, die besonders profitieren.
\item Was fordert der Finanzminister? (Antwort auf Deutsch)
\item Warum bleibt die Arbeitslosigkeit junger Akademiker hoch?
\item Nennen Sie ein Hindernis für den Handel trotz EPA.
\end{enumerate}
\tcblower
\textbf{Solution détaillée :}
\begin{enumerate}
\item \textbf{12\,\%}. (Texte : \all{stiegen ... um 12\,\% im Vergleich zum Vorjahr}).
\item \textbf{Energiesektor} (Wasserkraft, Solarenergie) und \textbf{Agrarindustrie} (Kakao-Verarbeitung, Holz). (Texte : \all{besonders der Energiesektor ... und die Agrarindustrie profitieren}).
\item \all{Er fordert Technologietransfer, nicht nur Kapital.} (Discours indirect : \all{betonte} + Konjunktiv I \all{brauche} $\rightarrow$ citation directe „Wir brauchen“ rapportée). Réponse synthétique : \all{Der Minister betont, dass Kamerun Technologietransfer brauche, nicht nur Kapital.}
\item Parce que les nouvelles usines sont \all{hochautomatisiert} (hautement automatisées) $\rightarrow$ peu d'embauche de main-d'œuvre qualifiée.
\item \textbf{Bürokratische Hürden und Korruption} (obstacles bureaucratiques et corruption). (Texte : \all{bremsen den Handel noch immer}).
\end{enumerate}
\textbf{Conclusion :} Repérer les chiffres (\cle{Zahlen}), les secteurs (\cle{Branchen}) et les relations de cause à effet (\cle{Ursache/Wirkung : automatisiert $\rightarrow$ Arbeitslosigkeit}) est essentiel pour le Bac.
\end{exoresolu}

\begin{methode}[Maîtriser le comparatif et le superlatif en contexte économique]
\textbf{Objectif :} Comparer des indicateurs économiques (PIB, inflation, taux de change, coûts de production) ou des entreprises/marchés.
\begin{enumerate}
\item \textbf{Identifier les deux termes de comparaison} (A vs B). Ex: \all{Kamerun} vs \all{Deutschland} ; \all{Investition A} vs \all{Investition B}.
\item \textbf{Choisir la structure correcte} :
\begin{itemize}
\item \textbf{Comparatif d'égalité} : \all{so + Adjektiv/Adverb + wie} (aussi ... que). \emph{Ex:} \all{Die Inflation in Kamerun ist so hoch wie in Nigeria.}
\item \textbf{Comparatif de supériorité} : \all{Adjektiv-Komparativ + als} (plus ... que). \emph{Ex:} \all{Das Bruttoinlandsprodukt Deutschlands ist höher als das Kameruns.}
\item \textbf{Comparatif d'infériorité} : \all{nicht so + Adjektiv + wie} (moins ... que) OU \all{weniger + Adjektiv + als}. \emph{Ex:} \all{Die Löhne sind nicht so hoch wie in Europa.}
\item \textbf{Superlatif} : \all{am + Adjektiv-Superlativ} (prédicat) ou \all{der/die/das ... -ste} (attribut). \emph{Ex:} \all{Die USA sind die größte Volkswirtschaft. / Der Dollar ist am stärksten.}
\end{itemize}
\item \textbf{Attention aux formes irrégulières fréquentes en économie} :
\all{hoch -- höher -- (am) höchsten} | \all{viel -- mehr -- (am) meisten} | \all{wenig -- weniger -- (am) wenigsten} | \all{gut -- besser -- (am) besten} | \all{nah -- näher -- (am) nächsten}.
\item \textbf{Place du verbe} : En proposition principale, le verbe conjugué est en \cle{position 2}. La proposition comparative (\all{als/ wie ...}) se place généralement en \textbf{champ final} (après le participe, l'infinitif ou la particule).
\item \textbf{Accord (attribut)} : Si adjectif attribut (devant nom), déclinaison forte/faible/mixte selon l'article. \emph{Ex:} \all{ein höheres Wachstum} (forte), \all{das höchste Wachstum} (faible).
\end{enumerate}
\end{methode}

\begin{exoresolu}[Transformation et thème : Comparatifs économiques]
\textbf{Exercice A (Transformation) :} Mettez l'adjectif entre parenthèses au comparatif de supériorité ou d'infériorité selon le sens.
\begin{enumerate}
\item Die Exporte Kameruns sind \_\_\_\_\_\_\_\_\_\_ (hoch) die Importe. (Réalité : souvent faux, mais structure demandée).
\item Die Arbeitslosigkeit in Douala ist \_\_\_\_\_\_\_\_\_\_ (niedrig) in Yaoundé. (Hypothèse).
\item Der Transport per Schiff ist \_\_\_\_\_\_\_\_\_\_ (billig) per Flugzeug.
\item Die Zinsen in der Eurozone sind \_\_\_\_\_\_\_\_\_\_ (niedrig) in den USA (Stand 2024).
\end{enumerate}
\textbf{Exercice B (Thème) :} Traduisez en allemand.
\begin{enumerate}
\item Le cacao camerounais est meilleur que le cacao ivoirien. (Qualité subjective).
\item Nous avons besoin de plus d'investissements (Investitionen) que l'année dernière.
\item C'est l'usine la plus moderne de la région.
\end{enumerate}
\tcblower
\textbf{Solution détaillée :}
\textbf{Exercice A :}
\begin{enumerate}
\item \all{sind höher als} (Comparatif supériorité : \all{hoch -- höher}). Attention : \all{als} déclenche le même cas que l'élément comparé (\all{sind} $\rightarrow$ Nominatif : \all{die Importe}).
\item \all{ist niedriger als} (\all{niedrig -- niedriger}). \emph{e} final $\rightarrow$ \all{er} ajouté.
\item \all{ist billiger als} (\all{billig -- billiger}). Monosyllabe $\rightarrow$ pas d'Umlaut pour \all{billig}. \all{als} suit.
\item \all{sind niedriger als} (\all{niedrig -- niedriger}). Pluriel \all{Zinsen} $\rightarrow$ \all{sind}.
\end{enumerate}
\textbf{Exercice B :}
\begin{enumerate}
\item \all{Kamerunischer Kakao ist besser als ivorischer Kakao.} (Irrégulier \all{gut -- besser}). \emph{Note :} \all{ivorischer} (adjectif dérivé de \all{Elfenbeinküste}, déclinaison forte nominatif masculin/neutre).
\item \all{Wir brauchen mehr Investitionen als letztes Jahr.} (\all{viel -- mehr} pour quantité non comptable/pluriel). \all{letztes Jahr} (accusatif temps) ou \all{als im letzten Jahr}.
\item \all{Das ist die modernste Fabrik der Region.} (Superlatif attributif : \all{die} + \all{modernste} + \all{Fabrik}). \all{modern} $\rightarrow$ \all{modernste} (\emph{e} muet conservé). Génitif \all{der Region} (féminin).
\end{enumerate}
\textbf{Point de vigilance :} Ne pas confondre \all{als} (comparatif) et \all{wie} (égalité). Ne pas oublier l'Umlaut sur les monosyllabes \all{a, o, u} au comparatif (\all{alt -- älter, groß -- größer, kurz -- kürzer}).
\end{exoresolu}

\begin{methode}[Utiliser le Konjunktiv II (Conditionnel) : politesse, hypothèses, conseils]
\textbf{Objectif :} Exprimer une hypothèse économique (« Si j'étais ministre... »), formuler une demande polie dans une lettre commerciale (« Nous souhaiterions... »), donner un avis prudent (« Cela serait risqué »).
\begin{enumerate}
\item \textbf{Deux formes principales} :
\begin{itemize}
\item \textbf{Forme simple (Präteritum-Konjunktiv)} : Verbes forts/irréguliers + Umlaut souvent. \all{ich wäre, ich hätte, ich käme, ich gäbe, ich wüsste, ich täte}. \emph{Usage :} Verbes auxiliaires/modaux/fréquents (\all{sein, haben, werden, können, müssen, wollen, sollen, dürfen, mögen, geben, wissen}).
\item \textbf{Forme avec \all{würde} + Infinitif} : \all{ich würde machen, ich würde investieren, ich würde exportieren}. \emph{Usage :} Tous les autres verbes (faibles, forts peu usuels). \textbf{Standard à l'écrit et à l'oral moderne.}
\end{itemize}
\item \textbf{Structure conditionnelle (Si-clause) :} \all{Wenn} + \textbf{Konjunktiv II} (dans les deux propositions pour l'irréel du présent/futur). \emph{Ex:} \all{Wenn die Zinsen sänken, würden wir mehr investieren.} (Si les taux baissaient, nous investirions plus).
\item \textbf{Formules de politesse commerciales (Feststehende Formeln)} : Les apprendre par cœur.
\begin{itemize}
\item \all{Wir \textbf{würden} Sie bitten, ...} (Nous vous prierions de...)
\item \all{Könnten Sie uns bitte ... mitteilen?} (Pourriez-vous nous indiquer...?)
\item \all{Wir \textbf{wären} Ihnen dankbar, wenn Sie ...} (Nous vous saurions gré si vous...)
\item \all{Wir \textbf{möchten} (Präteritum Konj. II de mögen) Ihnen unser Angebot unterbreiten.}
\end{itemize}
\item \textbf{Conseil (Rat) :} \all{An Ihrer Stelle \textbf{würde} ich ...} (À votre place, je...). \all{Man \textbf{sollte} (sollen) ...} (On devrait...).
\end{enumerate}
\end{methode}

\begin{exoresolu}[Konjunktiv II : Lettre de réclamation et hypothèses]
\textbf{Exercice A (Conjugaison) :} Mettez les verbes au Konjunktiv II (forme simple si possible, sinon \all{würde}-Form).
\begin{enumerate}
\item Wenn wir mehr Kapital \_\_\_\_\_\_\_\_\_\_ (haben), wir \_\_\_\_\_\_\_\_\_\_ (können) expandieren.
\item Der Manager \_\_\_\_\_\_\_\_\_\_ (vorschlagen), dass wir die Preise \_\_\_\_\_\_\_\_\_\_ (senken).
\item Ich \_\_\_\_\_\_\_\_\_\_ (möchten) einen Termin mit dem Direktor \_\_\_\_\_\_\_\_\_\_ (vereinbaren).
\end{enumerate}
\textbf{Exercice B (Thème - Lettre formelle) :} Traduisez en allemand (style commercial).
« Monsieur le Directeur, nous souhaiterions commander 500 tonnes de cacao. Pourriez-vous nous envoyer une offre ? Nous serions reconnaissants d'une réponse rapide. »
\tcblower
\textbf{Solution détaillée :}
\textbf{Exercice A :}
\begin{enumerate}
\item \all{Wenn wir mehr Kapital \textbf{hätten}, \textbf{könnten} wir expandieren.}
(\all{haben $\rightarrow$ hätte} ; \all{können $\rightarrow$ könnte}. Formes simples obligatoires pour auxiliaires/modaux. Verbe à la fin dans le \all{Wenn}-Satz : \all{hätten}. Verbe en position 2 dans la principale : \all{könnten}).
\item \all{Der Manager \textbf{würde vorschlagen}, dass wir die Preise \textbf{senkten}.}
(\all{vorschlagen} : verbe fort mais \all{würde}-forme préférée ; \all{schlagen $\rightarrow$ schlüge} possible mais soutenu). Subordonnée \all{dass} $\rightarrow$ verbe à la fin : \all{senkten} (verbe faible \all{senken} = baisser [transitif], pas \all{sinken} = baisser [intransitif]). \emph{Note :} \all{vorschlagen} $\rightarrow$ \all{schlug vor} (Prät) $\rightarrow$ \all{schlüge vor} (KII). \all{würde vorschlagen} est plus sûr.
\item \all{Ich \textbf{möchte} einen Termin mit dem Direktor \textbf{vereinbaren}.}
(\all{mögen $\rightarrow$ möchte} KII simple. \all{vereinbaren} infinitif à la fin). \emph{Attention :} \all{Ich würde gerne vereinbaren} possible mais \all{möchte} est la forme standard de politesse pour « je voudrais ».
\end{enumerate}
\textbf{Exercice B :}
\begin{textebox}[Lettre formelle type]
Sehr geehrter Herr Direktor, (ou Frau Direktorin)

wir \textbf{möchten} 500 Tonnen Kakao bestellen.

\textbf{Könnten} Sie uns bitte ein Angebot \textbf{zusenden}? (ou \all{schicken / senden})

Wir \textbf{wären} Ihnen für eine schnelle Antwort \textbf{dankbar}.

Mit freundlichen Grüßen

[Nom]
\end{textebox}
\textbf{Justifications :}
\begin{itemize}
\item \all{möchten} (KII de \all{mögen}) pour « souhaiter » (politesse standard).
\item \all{Könnten Sie ... zusenden?} (KII de \all{können} + infinitif). \all{zusenden} (verbe à particule séparable : particule \all{zu-} à la fin en infinitif).
\item \all{wären ... dankbar} (KII de \all{sein} + adjectif \all{dankbar} + \all{für} + accusatif). Structure figée.
\end{itemize}
\end{exoresolu}

\begin{methode}[Construire le Vorgangspassiv (Passif de procès) : Rapports et comptes rendus]
\textbf{Objectif :} Déplacer le focus de l'acteur (souvent inconnu ou indifférent : « on », « l'entreprise ») vers l'action/le résultat (« le contrat est signé », « l'usine a été construite »). Indispensable pour les \cle{Protokolle} (comptes rendus) et \cle{Geschäftsberichte} (rapports d'activité).
\begin{enumerate}
\item \textbf{Formule :} \textbf{werden} (conjugué) + \textbf{Partizip II} (à la fin de la phrase).
\item \textbf{Transformation Active $\rightarrow$ Passif :}
\begin{itemize}
\item Sujet Actif (Nominatif) $\rightarrow$ \all{von} + Datif (agent) \textbf{OU} suppression (si « man/on », gens, ils).
\item Objet Actif (Accusatif) $\rightarrow$ \textbf{Sujet Passif (Nominatif)} $\rightarrow$ Accord du verbe \all{werden} avec CE nouveau sujet.
\item Verbe Actif $\rightarrow$ \all{werden} (même temps) + \all{Partizip II}.
\end{itemize}
\item \textbf{Les temps principaux au Bac :}
\begin{center}
\begin{tabular}{|l|l|l|}\hline
\textbf{Temps} & \textbf{Actif (Ex: bauen)} & \textbf{Passiv (werden gebaut)} \\ \hline
Präsens & \all{man baut / die Firma baut} & \all{es wird gebaut / die Fabrik wird gebaut} \\ \hline
Präteritum & \all{man baute} & \all{es wurde gebaut} \\ \hline
Perfekt & \all{man hat gebaut} & \all{es ist gebaut worden} (\textbf{ist} + Partizip II \all{worden} -- \textbf{sans ge-}) \\ \hline
Futur I & \all{man wird bauen} & \all{es wird gebaut werden} \\ \hline
Konjunktiv II & \all{man würde bauen} & \all{es würde gebaut werden} \\ \hline
\end{tabular}
\end{center}
\item \textbf{Passif modal (Ersatzformen) :} \all{zu} + Infinitif (\all{das ist zu tun}), \all{sich lassen} (\all{das lässt sich machen}). Fréquent dans les recommandations : \all{Die Kosten sind zu senken.} (Les coûts doivent être réduits).
\item \textbf{Attention :} Verbes intransitifs (pas d'accusatif) $\rightarrow$ Pas de passif standard (sauf \all{es wird getanzt} -- impersonnel). Verbes avec double accusatif (\all{lehren, fragen, kosten, bitten}) $\rightarrow$ seul l'accusatif de la personne devient sujet passif. \emph{Ex:} \all{Man fragt ihn nach dem Preis} $\rightarrow$ \all{Er wird nach dem Preis gefragt.}
\end{enumerate}
\end{methode}

\begin{exoresolu}[Transformation Active $\rightarrow$ Passif \& Compte rendu]
\textbf{Exercice A :} Transformez au passif (temps indiqué). Supprimez l'agent « man / die Firma ».
\begin{enumerate}
\item (Präsens) \all{Die Firma baut eine neue Fabrik in Kribi.}
\item (Präteritum) \all{Man hat den Vertrag gestern unterschrieben.}
\item (Perfekt) \all{Die Bank hat dem Projekt 10 Millionen Euro bewilligt.} (Datif !)
\item (Futur I) \all{Man wird die Maschinen nächste Woche liefern.}
\end{enumerate}
\textbf{Exercice B (Rédaction - Compte rendu) :} Rédigez 3 phrases au passif (Präteritum/Perfekt) pour un compte rendu de réunion : « Décision : augmentation des salaires de 5\,\%. Signature du contrat le 12.06. Ouverture de la nouvelle succursale prévue pour août. »
\tcblower
\textbf{Solution détaillée :}
\textbf{Exercice A :}
\begin{enumerate}
\item \all{Eine neue Fabrik \textbf{wird} in Kribi \textbf{gebaut}.} (Sujet \all{die Fabrik} f $\rightarrow$ \all{wird}. \all{bauen} $\rightarrow$ \all{ge-baut}). Lieu \all{in Kribi} reste.
\item \all{Der Vertrag \textbf{wurde} gestern \textbf{unterschrieben}.} (Präteritum : \all{werden $\rightarrow$ wurde}. Sujet \all{der Vertrag} m. \all{unterschreiben} $\rightarrow$ \all{unterschrieben}). \all{gestern} reste.
\item \all{Dem Projekt \textbf{sind} 10 Millionen Euro \textbf{bewilligt worden}.} (\textbf{Piège Datif !} \all{bewilligen} prend Datif (\all{dem Projekt}) + Accusatif (\all{das Geld}). Seul l'accusatif devient sujet nominatif : \all{10 Millionen Euro} (Pluriel) $\rightarrow$ \all{sind}. \all{dem Projekt} RESTE DATIF. \all{Perfekt Passiv} : \all{ist/sind + Partizip II + worden}).
\item \all{Die Maschinen \textbf{werden} nächste Woche \textbf{geliefert werden}.} (Futur I : \all{werden} (aux futur) + \all{geliefert} (Partizip II) + \all{werden} (infinitif passif). Sujet \all{die Maschinen} Pl $\rightarrow$ \all{werden}).
\end{enumerate}
\textbf{Exercice B (Proposition de correction) :}
\begin{enumerate}
\item \all{Es \textbf{wurde beschlossen}, die Gehälter um 5\,\% \textbf{zu erhöhen}.} (Passif impersonnel \all{es} + infinitif \all{zu erhöhen}). Ou : \all{Die Gehälter \textbf{wurden} um 5\,\% \textbf{erhöht}.}
\item \all{Der Vertrag \textbf{wurde} am 12.\,Juni \textbf{unterzeichnet}.}
\item \all{Die neue Filiale \textbf{wird} im August \textbf{eröffnet werden}.} (Futur) ou \all{Die Eröffnung der neuen Filiale \textbf{ist} für August \textbf{geplant}.} (Zustandspassiv / \all{sein} + Partizip II -- très courant pour « prévu »).
\end{enumerate}
\textbf{Point clé Bac :} Le \textbf{Perfekt Passiv} (\all{ist ... worden}) et le \textbf{Futur Passif} (\all{wird ... werden}) sont des classiques. Ne pas oublier \all{worden} \textbf{sans ge-} !
\end{exoresolu}

\begin{methode}[Rédiger une lettre commerciale formelle (DIN 5008) et un compte rendu]
\textbf{Objectif :} Produire un texte structuré, normé, lexicalement précis (formules de politesse, objet, pièces jointes) pour le Bac (Expression écrite / Traduction).
\begin{enumerate}
\item \textbf{Structure de la lettre commerciale (Geschäftsbrief) :}
\begin{enumerate}
\item \textbf{Kopfzeile (En-tête) :} Expéditeur (Absender) haut gauche. Destinataire (Empfänger) avec \all{z. Hd. Herrn/Frau ...} (à l'attention de). Date (ligne vide au-dessus), Lieu + Date : \all{Yaoundé, den 15. Juni 2025}.
\item \textbf{Betreffzeile (Objet) :} Gras, sans « Betreff: » (norme moderne), résumé précis. \emph{Ex:} \all{Angebot für 500 t Kakao / Ihre Anfrage vom 10.06.2025}.
\item \textbf{Anrede (Formule d'appel) :} \all{Sehr geehrter Herr Müller,} / \all{Sehr geehrte Frau Müller,} / \all{Sehr geehrte Damen und Herren,} (si nom inconnu). \textbf{Virgule} puis minuscule à la ligne suivante.
\item \textbf{Textkörper (Corps) :}
\begin{itemize}
\item Référence : \all{Bezugnehmend auf Ihre Anfrage ... / Wir danken für Ihr Schreiben vom ...}
\item Objet principal : Clair, paragraphes distincts. \all{Konjunktiv II} pour politesse (\all{wir möchten, wir bitten, könnten Sie}).
\item Call to action / Prochaine étape : \all{Wir freuen uns auf Ihre Antwort. / Bitte senden Sie uns ...}
\end{itemize}
\item \textbf{Grußformel (Formule de politesse finale) :} \all{Mit freundlichen Grüßen} (Standard). \all{Mit freundlichen Grüßen aus Yaoundé}.
\item \textbf{Unterschrift + Namenszug (Signature + Nom imprimé) :} Signature manuscrite, dessous nom en clair. Fonction/titre optionnel.
\item \textbf{Anlagenvermerk (Pièces jointes) :} \all{Anlagen: 1. Preisliste 2. AGB} (en bas à gauche).
\end{enumerate}
\item \textbf{Structure du Compte rendu (Protokoll / Bericht) :}
\begin{enumerate}
\item \textbf{Titre :} \all{Protokoll der Sitzung vom ... / Geschäftsbericht Q2 2025}.
\item \textbf{En-tête :} Date, Lieu, Participants (Anwesende), Absents (Abwesende), Président (Leitung), Secrétaire (Protokollführung).
\item \textbf{Ordre du jour (Tagesordnung) :} Points numérotés.
\item \textbf{Corps :} Style \textbf{objectif, impersonnel, passif, Präteritum/Perfekt}. \all{Es wurde beschlossen, dass ... / Der Leiter berichtete über ... / Es wurde vereinbart, ... bis zum ... durchzuführen.}
\item \textbf{Décisions / Mesures (Beschlüsse / Maßnahmen) :} Tableau : \all{Was? Wer? Bis wann? Status}.
\item \textbf{Signature :} \all{Gezeichnet: [Präsident] [Secrétaire]}.
\end{enumerate}
\end{enumerate}
\end{methode}

\begin{exoresolu}[Rédaction : Lettre de commande / Offre]
\textbf{Consigne (Type Bac) :} Vous êtes le gérant de \all{AgroExport Cameroon SARL} (Douala). Écrivez une lettre formelle à \all{Maschinenbau GmbH, Berlin} pour commander une machine de séchage de cacao (Modell \all{TrockenPro 5000}), vue à la foire de Hanovre. Demandez confirmation de prix, délai de livraison à Douala port, et conditions de paiement (Lettre de crédit). Mentionnez PJ : Spécifications techniques.
\tcblower
\textbf{Solution détaillée (Modèle de lettre) :}
\begin{textebox}[Lettre commerciale type]
AgroExport Cameroon SARL\\
B.P. 1234 Akwa\\
Douala, Kamerun\\
Tel: +237 233 44 55 66\\
E-Mail: export@agrocam.cm\\

Maschinenbau GmbH\\
Industriestraße 45\\
10115 Berlin\\
Deutschland\\

z. Hd. Herrn Dr. Thomas Becker\\

Douala, den 20. Oktober 2025\\

\textbf{Bestellung: Trockenanlage Typ TrockenPro 5000 -- Ihre Messepräsentation Hannover}\\

Sehr geehrter Herr Dr. Becker,\\

bezugnehmend auf Ihr Angebot vom 15. September 2025 und die Gespräche auf der Hannover Messe möchten wir Ihnen folgende Bestellung übermitteln:\\

1 \texttimes\ Trockenanlage \textbf{TrockenPro 5000} (Kapazität: 5 t/Tag), gemäß beiliegender technischer Spezifikation (Anlage 1).\\

\textbf{Bitte bestätigen Sie uns schriftlich:}
\begin{itemize}
\item den endgültigen Preis inkl. Verpackung und Versicherung CIF Douala Hafen;
\item den voraussichtlichen Liefertermin (gewünscht: November 2025);
\item die Zahlungsbedingungen: Wir bevorzugen Zahlung per \textbf{Akkreditiv (Letter of Credit)} bei Versanddokumenten.
\end{itemize}

Wir wären Ihnen für eine zeitnahe Auftragsbestätigung dankbar.\\

Mit freundlichen Grüßen\\

\emph{(Unterschrift)}\\

Jean-Paul Mbarga\\
Geschäftsführer\\

Anlagen: 1. Technische Spezifikation TrockenPro 5000
\end{textebox}
\textbf{Analyse grammaticale et lexicale :}
\begin{itemize}
\item \all{bezugnehmend auf + Akk.} (Participe I prépositionnel, très formel).
\item \all{möchten ... übermitteln} (KII \all{mögen} + infinitif séparable \all{übermitteln}).
\item \all{gemäß + Dat.} (Préposition formelle « conformément à »).
\item \all{bitten, ... zu bestätigen} (Structure \all{bitten um} ou \all{bitten, zu + Inf.}).
\item \all{CIF} (Incoterm : Cost, Insurance, Freight). \all{Akkreditiv} (Terme bancaire précis).
\item \all{zeitnah} (Adjectif/Adverbe : rapide, dans les délais).
\item \all{Anlagen} (Pluriel de \all{die Anlage} -- Pièce jointe).
\end{itemize}
\end{exoresolu}

\begin{methode}[Stratégies de traduction : Thème (Fr $\rightarrow$ De) et Version (De $\rightarrow$ Fr)]
\textbf{Objectif :} Traduire fidèlement des phrases ou courts textes économiques en respectant la grammaire allemande (cas, ordre des mots, temps) et le registre (formel, technique).
\begin{enumerate}
\item \textbf{Analyse de la phrase source :} Identifier le \cle{sujet}, le \cle{verbe}, les \cle{compléments} (COD $\rightarrow$ Akkusativ, COI $\rightarrow$ Dativ, Circonstance $\rightarrow$ Préposition + Cas / Adverbe). Repérer le \cle{temps} et le \cle{mode} (Indicatif / Conditionnel / Subjonctif).
\item \textbf{Transfert lexical (Mots) :} Choisir le terme technique exact (\all{Wirtschaftswachstum}, \all{Handelsbilanz}, \all{Wechselkurs}, \all{Investitionsförderung}). \textbf{Attention aux Faux Amis :}
\begin{itemize}
\item \all{aktuell} = d'actualité, topical (pas « actuel » $\rightarrow$ \all{gegenwärtig, momentan}).
\item \all{eventuell} = éventuel, le cas échéant (pas « éventuel » $\rightarrow$ \all{möglich, potenziell}).
\item \all{konsequent} = cohérent, conséquent.
\item \all{sensibel} = sensible, délicat (pas « sensible » $\rightarrow$ \all{empfindlich, spürbar}).
\end{itemize}
\item \textbf{Restructuration syntaxique (Ordre des mots) :} \textbf{Règle d'or : Verbe conjugué en Position 2 (Principal) / Position Finale (Subordonnée).}
\begin{itemize}
\item \textbf{Champ initial (Position 1) :} Sujet \textbf{OU} Complément (Topicalisation). \emph{Ex:} \all{Gestern hat die Firma den Vertrag unterschrieben.}
\item \textbf{Champ moyen :} Sujet (si pas en P1), Adverbes (Temps - Manière - Lieu), Objets (Dativ avant Akkusativ \textbf{si pronoms}, sinon ordre libre selon focus), \textbf{Négation (\all{nicht})}.
\item \textbf{Champ final :} Participe II, Infinitif(s), Verbe à particule séparable (particule), Verbe conjugué (subordonnée).
\end{itemize}
\item \textbf{Gestion des cas (Kasus) :} \textbf{Le piège n°1 francophone.} \all{Ich helfe dem Mann} (Dativ !) pas \all{den Mann}. \all{Wir warten auf den Bus} (Akkusativ prépositionnel). \all{Der Preis des Öls} (Génitif).
\item \textbf{Révision finale (Relecture) :} \textbf{Majuscules aux noms} (\all{das Wachstum}), \textbf{Umlauts} (\all{Größe, Länder, Müller}), \textbf{ß vs ss} (\all{Straße, muss, Fluss}), \textbf{ponctuation} (virgules avant subordonnées \all{dass, weil, wenn, als}), \textbf{guillemets} \all{„...“}.
\end{enumerate}
\end{methode}

\begin{exoresolu}[Thème \& Version mixtes : Spécialité économique]
\textbf{Exercice A (Thème) :} Traduisez en allemand.
\begin{enumerate}
\item Bien que l'inflation soit haute, la banque centrale n'augmente pas les taux d'intérêt. (Concession + Présent).
\item Les exportations de cacao ont augmenté de 10\,\% l'année dernière. (Perfekt / Präteritum).
\item Si nous investissions dans l'énergie solaire, nous réduirions les coûts à long terme. (Hypothèse KII).
\item Le contrat doit être signé par le directeur demain. (Passif modal / Futur).
\end{enumerate}
\textbf{Exercice B (Version) :} Traduisez en français.
\all{„Trotz der globalen Rezession konnte Kamerun seine Exporteinnahmen steigern. Die Regierung plant, die Einnahmen in die Infrastruktur und Bildung zu investieren. Allerdings warnen Experten vor einer zu hohen Verschuldung in Fremdwährung.“}
\tcblower
\textbf{Solution détaillée :}
\textbf{Exercice A (Thème) :}
\begin{enumerate}
\item \all{Obwohl die Inflation hoch \textbf{ist}, erhöht die Zentralbank die Zinssätze \textbf{nicht}.}
(\all{Obwohl} $\rightarrow$ verbe à la fin \all{ist}. Principal : verbe P2 \all{erhöht}, sujet \all{die Zentralbank}, objet \all{die Zinssätze}, \all{nicht} en fin de phrase). \emph{Variante :} \all{Die Zentralbank erhöht die Zinssätze nicht, obwohl die Inflation hoch ist.}
\item \all{Die Kakaosexporte \textbf{sind} im letzten Jahr um 10\,\% \textbf{gestiegen}.} (Perfekt : \all{sein} verbe de mouvement/état pour \all{steigen}. \all{um 10\,\%} = de 10\,\%. \all{im letzten Jahr} = l'année dernière). \emph{Alternative Präteritum (style journalistique) :} \all{stiegen ... um 10\,\%}.
\item \all{Wenn wir in Solarenergie \textbf{investieren würden}, \textbf{würden} wir die Kosten langfristig \textbf{senken}.}
(\all{Wenn}-Satz : KII \all{würden investieren} (ou \all{investierten}) verbe fin. Principal : \all{würden} P2, sujet \all{wir}, objet \all{die Kosten}, adv \all{langfristig}, infinitif \all{senken} fin). \emph{Note :} \all{in} + Dativ (\all{Solarenergie} f).
\item \all{Der Vertrag \textbf{muss} morgen vom Direktor \textbf{unterschrieben werden}.} (Passif modal \all{müssen} + Infinitif Passif \all{unterschrieben werden}. Agent \all{vom} = \all{von dem} Direktor). \emph{Alternative Futur Passif :} \all{Der Vertrag wird morgen vom Direktor unterschrieben werden.}
\end{enumerate}
\textbf{Exercice B (Version) :}
« \textbf{Malgré la récession mondiale, le Cameroun a pu augmenter ses recettes d'exportation. Le gouvernement prévoit d'investir ces recettes dans les infrastructures et l'éducation. Toutefois, les experts mettent en garde contre un endettement excessif en devise étrangère.} »
\textbf{Justifications lexicales/grammaticales :}
\begin{itemize}
\item \all{Trotz + Genitiv} (\all{der globalen Rezession}) = Malgré. \all{können} (Präteritum \all{konnte}) + Infinitif \all{steigern} = a pu augmenter.
\item \all{planen, ... zu investieren} = prévoir de / avoir l'intention d'investir. \all{in + Akk} (\all{die Infrastruktur, die Bildung}).
\item \all{warnen vor + Dativ} (\all{einer zu hohen Verschuldung}) = mettre en garde contre. \all{Fremdwährung} = devise étrangère / monnaie étrangère. \all{Verschuldung} = endettement.
\end{itemize}
\end{exoresolu}

\begin{attention}[Les 5 erreurs qui coûtent des points au Bac (Vie économique)]
\begin{enumerate}
\item \textbf{Majuscules oubliées sur les noms :} \all{das wachstum, die inflation, der export, das bruttoinlandsprodukt} $\rightarrow$ \textbf{FAUX}. Toujours : \all{das Wachstum, die Inflation, der Export, das Bruttoinlandsprodukt}. \emph{Astuce :} Si on peut mettre \all{der/die/das} devant $\rightarrow$ MAJUSCULE.
\item \textbf{Confusion \all{als} / \all{wie} au comparatif :} \all{größer wie} (FAUX), \all{so groß als} (FAUX). \textbf{Règle :} Comparatif (\all{-er}) + \textbf{\all{als}} ; Égalité (\all{so ...}) + \textbf{\all{wie}}.
\item \textbf{Ordre des mots : Verbe pas en Position 2 (Principal) ou pas à la fin (Subordonnée).} \emph{Ex:} \all{*Gestern die Firma hat unterschrieben.} (FAUX). Correct : \all{Gestern \textbf{hat} die Firma unterschrieben.} / \all{..., \textbf{dass} die Firma gestern unterschrieben \textbf{hat}.}
\item \textbf{Cas (Kasus) : COD $\neq$ Accusatif systématique / Verbes à Dativ ignorés.} \emph{Principaux verbes économiques pièges :} \all{helfen} (Dativ), \all{danken} (Dativ), \all{gehorchen} (Dativ), \all{folgen} (Dativ), \all{vertrauen} (Dativ + \all{auf} + Akk), \all{warten} (Dativ + \all{auf} + Akk). \all{Ich helfe dem Unternehmen} (pas \all{das Unternehmen}).
\item \textbf{Faux-amis et calques lexicaux :}
\begin{itemize}
\item « Actuel » $\neq$ \all{aktuell} (actuel = \all{gegenwärtig, momentan} ; \all{aktuell} = d'actualité, topical).
\item « Éventuel » $\neq$ \all{eventuell} (éventuel = \all{möglich, potenziell} ; \all{eventuell} = au cas où, le cas échéant).
\item « Sensible » $\neq$ \all{sensibel} (sensible = \all{empfindlich, heikel} ; \all{sensibel} = sensible, délicat, perceptif).
\item « Contrôle » $\neq$ \all{Kontrolle} (souvent \all{Prüfung, Überwachung, Steuerung} selon contexte). \all{Kontrolle} = vérification, check.
\item « Demander » $\neq$ \all{demandieren} (exiger, réclamer fermement). « Demander » = \all{bitten um, fragen nach, fordern}.
\end{itemize}
\end{enumerate}
\textbf{Check-list relecture (5 min) :} 1. Majuscules Noms ? 2. Verbes Position 2 / Fin ? 3. Cas (der/den/dem) corrects ? 4. Umlauts / ß ? 5. Vocabulaire précis (pas de \all{machen/tun} vague) ?
\end{attention}

\newpage

\coursec{Exercices}

\courssub{Je vérifie mes connaissances}

\begin{exercice}[QCM : grammaire de la vie économique]
Pour chaque phrase, une seule réponse est correcte. Cochez la bonne case.
\begin{enumerate}
    \item Comment forme-t-on le comparatif de supériorité de l'adjectif \all{groß} ?
    \begin{itemize}
        \item[A)] \all{größer als}
        \item[B)] \all{mehr groß als}
        \item[C)] \all{am größten}
        \item[D)] \all{größer wie}
    \end{itemize}
    \item Quelle forme correspond au \emph{Konjunktiv II} (conditionnel) du verbe \all{können} à la 1\iere personne du pluriel ?
    \begin{itemize}
        \item[A)] \all{wir könnten}
        \item[B)] \all{wir können}
        \item[C)] \all{wir würden können}
        \item[D)] \all{wir gekonnt hätten}
    \end{itemize}
    \item Transformez la phrase active \all{„Der Direktor unterschreibt den Vertrag.“} en passif (Vorgangspassiv) au présent.
    \begin{itemize}
        \item[A)] \all{Der Vertrag wird vom Direktor unterschrieben.}
        \item[B)] \all{Der Vertrag wurde vom Direktor unterschrieben.}
        \item[C)] \all{Der Vertrag ist vom Direktor unterschrieben.}
        \item[D)] \all{Der Vertrag wird unterschreiben vom Direktor.}
    \end{itemize}
    \item Dans une lettre commerciale formelle, quelle formule de politesse de clôture est la plus appropriée ?
    \begin{itemize}
        \item[A)] \all{Liebe Grüße}
        \item[B)] \all{Mit freundlichen Grüßen}
        \item[C)] \all{Bis bald}
        \item[D)] \all{Dein}
    \end{itemize}
\end{enumerate}
\end{exercice}

\begin{exercice}[Vrai ou Faux ? Justifiez en allemand]
Déterminez si les affirmations suivantes sont vraies (\all{Richtig}) ou fausses (\all{Falsch}). Justifiez votre réponse par une phrase complète en allemand.
\begin{enumerate}
    \item Le superlatif de \all{viel} est \all{am meisten}.
    \item Le \emph{Konjunktiv II} du verbe \all{sein} au présent s'écrit \all{ich war}.
    \item Dans le passif (Vorgangspassiv), l'agent est toujours introduit par la préposition \all{durch}.
    \item L'objet d'une lettre commerciale (\all{Der Betreff}) se place après la formule d'appel (\all{Die Anrede}).
    \item \all{Die Lieferung} est un mot masculin.
\end{enumerate}
\end{exercice}

\begin{exercice}[Vocabulaire thématique : la vie économique]
Associez chaque terme allemand à sa définition ou à son équivalent français.
\begin{center}
\begin{tabularx}{\linewidth}{|X|X|}
\hline
\rowcolor{popblueL} \textbf{Deutsch} & \textbf{Français / Définition} \\
\hline
\all{die Lieferung} & A. Document listant des marchandises et leurs prix \\
\hline
\all{die Rechnung} & B. Action de transporter des marchandises au client \\
\hline
\all{der Angebot} & C. Réclamation formelle auprès d'un fournisseur \\
\hline
\all{die Beschwerde} & D. Proposition commerciale avec prix et conditions \\
\hline
\all{der Zahlungsaufschub} & E. Facture, demande de paiement \\
\hline
\all{der Vertrag} & F. Délai accordé pour payer une dette \\
\hline
\all{die Mahnung} & G. Accord juridiquement contraignant \\
\hline
\all{der Liefertermin} & H. Rappel de paiement (mise en demeure) \\
\hline
& I. Date convenue pour la livraison \\
\hline
\end{tabularx}
\end{center}
\textit{Écrivez les paires correctes (ex. 1--B).}
\end{exercice}

\begin{exercice}[Conjugaison et formes grammaticales en contexte]
Complétez le texte suivant en conjuguant les verbes entre parenthèses au temps et au mode indiqués, ou en mettant l'adjectif au degré demandé.
\begin{textebox}[Lettre commerciale incomplète]
Sehr geehrte Damen und Herren,

bezugnehmend auf Ihr Angebot vom 10. Mai \_\_\_\_\_\_\_\_\_\_\_\_ (wir / freuen uns / Konjunktiv II Präsens), Ihnen mitteilen zu können, dass wir den Auftrag erteilen möchten. Die Preise \_\_\_\_\_\_\_\_\_\_\_\_ (sein / Komparativ) als die der Konkurrenz, und die Qualität \_\_\_\_\_\_\_\_\_\_\_\_ (sein / Superlativ). 

Die Ware \_\_\_\_\_\_\_\_\_\_\_\_ (müssen / Passiv Präsens / liefern) bis zum 30. Juni. Sollte der Termin nicht eingehalten \_\_\_\_\_\_\_\_\_\_\_\_ (werden / Passiv Futur I), \_\_\_\_\_\_\_\_\_\_\_\_ (wir / vom Vertrag zurücktreten / Konjunktiv II).

Mit freundlichen Grüßen
\end{textebox}
\end{exercice}

\courssub{Je m'entraîne}

\begin{exercice}[Thème : phrases à traduire en allemand]
Traduisez les phrases suivantes en allemand en respectant la grammaire révisée (comparatif, conditionnel, passif, vocabulaire économique).
\begin{enumerate}
    \item Si les prix étaient plus bas, nous achèterions plus de produits.
    \item La lettre a été écrite hier par le directeur. (Präteritum Passiv)
    \item Le marché allemand est plus grand que le marché autrichien.
    \item Pourriez-vous nous envoyer votre catalogue, s'il vous plaît ? (Formule de politesse \emph{Konjunktiv II})
    \item Les marchandises doivent être livrées avant vendredi. (Passiv mit Modalverb)
    \item Nous vous prions de bien vouloir régler la facture numéro 456.
    \item Cette offre est la meilleure que nous ayons reçue. (Superlatif)
    \item Le paiement a été effectué la semaine dernière. (Perfekt Passiv)
\end{enumerate}
\end{exercice}

\begin{exercice}[Transformation : Voix active $\rightarrow$ Voix passive (Vorgangspassiv)]
Transformez les phrases actives suivantes en phrases passives. Conservez le temps verbal indiqué. N'oubliez pas la majuscule au nom de l'agent si vous le conservez (préposition \all{von}).
\begin{enumerate}
    \item \textbf{Présent :} \all{Die Firma Müller liefert die Maschinen.}
    \item \textbf{Präteritum :} \all{Der Kunde bezahlte die Rechnung nicht.}
    \item \textbf{Perfekt :} \all{Der Spediteur hat die Pakete abgeholt.}
    \item \textbf{Futur I :} \all{Wir werden den Vertrag nächste Woche unterschreiben.}
    \item \textbf{Avec modal (Présent) :} \all{Man muss die Waren sorgfältig prüfen.} (Sujet impersonnel \all{man} $\rightarrow$ passif sans agent)
    \item \textbf{Avec modal (Präteritum) :} \all{Der Chef konnte die Entscheidung nicht treffen.}
    \item \textbf{Présent :} \all{Man schreibt den Bericht auf Deutsch.}
    \item \textbf{Perfekt :} \all{Jemand hat das Fenster geöffnet.} (Agent indéfini $\rightarrow$ passif sans \all{von})
\end{enumerate}
\end{exercice}

\begin{exercice}[Texte à trous : Comparatif et Superlatif]
Complétez le rapport économique suivant en mettant les adjectifs/adverbes entre parenthèses au comparatif ou au superlatif (selon le sens). Attention aux formes irrégulières !
\begin{textebox}[Wirtschaftsbericht: Deutschland und Österreich]
Österreich und Deutschland sind wichtige Handelspartner. Der deutsche Markt ist \_\_\_\_\_\_\_\_\_\_\_\_ (groß) als der österreichische Markt, aber die österreichische Exportquote ist oft \_\_\_\_\_\_\_\_\_\_\_\_ (hoch). Schweizer Uhren sind weltweit \_\_\_\_\_\_\_\_\_\_\_\_ (bekannt) für ihre Präzision. 

Im vergangenen Jahr haben deutsche Unternehmen \_\_\_\_\_\_\_\_\_\_\_\_ (viel) in Österreich investiert als je zuvor. Die Arbeitslosenzahlen in Deutschland sind \_\_\_\_\_\_\_\_\_\_\_\_ (niedrig) als in vielen anderen EU-Ländern. Österreichische Unternehmen arbeiten \_\_\_\_\_\_\_\_\_\_\_\_ (gern) mit deutschen Partnern zusammen, weil die Sprache \_\_\_\_\_\_\_\_\_\_\_\_ (einfach) ist. 

Das \_\_\_\_\_\_\_\_\_\_\_\_ (gut) Ergebnis erzielt man jedoch durch gemeinsame Forschung. Die Zusammenarbeit ist \_\_\_\_\_\_\_\_\_\_\_\_ (wichtig) geworden in den letzten Jahren.
\end{textebox}
\end{exercice}

\begin{exercice}[Appariement : Structure de la lettre commerciale et du compte rendu]
Reliez chaque élément de structure à sa description ou à son emplacement correct.
\begin{center}
\begin{tabularx}{\linewidth}{|X|X|}
\hline
\rowcolor{popblueL} \textbf{Élément (Deutsch)} & \textbf{Description / Contenu} \\
\hline
1. \all{Der Briefkopf (Absender)} & A. Formules : \all{Sehr geehrte Damen und Herren,} \\
\hline
2. \all{Das Datum} & B. Objet de la lettre, en gras, sans « Betreff: » \\
\hline
3. \all{Die Anrede} & C. Nom, adresse, téléphone, email de l'expéditeur \\
\hline
4. \all{Der Betreffzeile} & D. Corps du texte : motifs, arguments, demandes \\
\hline
5. \all{Der Fließtext (Hauptteil)} & E. Lieu, date du jour (ex. \all{Yaoundé, den 15. Oktober 2025}) \\
\hline
6. \all{Die Grußformel} & F. \all{Mit freundlichen Grüßen} \\
\hline
7. \all{Die Unterschrift} & G. Nom manuscrit de l'auteur \\
\hline
8. \all{Der Anhang (Anlagen)} & H. Liste des pièces jointes : \all{Anlagen: Angebot, AGB} \\
\hline
\multicolumn{2}{|l|}{\textbf{Structure du compte rendu (\all{Der Bericht / Protokoll})}} \\
\hline
9. \all{Einleitung} & I. Résultats, décisions, mesures concrètes \\
\hline
10. \all{Hauptteil} & J. Sujet, date, lieu, participants, objectif de la visite/réunion \\
\hline
11. \all{Schluss} & K. Résumé, suites à donner, signature, date \\
\hline
\end{tabularx}
\end{center}
\end{exercice}

\courssub{Je me prépare au Bac}

\begin{exercice}[Compréhension et langue : Les exportations suisses]
Lisez le texte suivant et répondez aux questions en allemand (phrases complètes).
\begin{textebox}[Die Schweizer Exportwirtschaft 2024]
Die Schweiz ist ein Exportland par excellence. Ohne den Außenhandel wäre der Wohlstand nicht zu halten. Im Jahr 2023 beliefen sich die Ausfuhren auf über 270 Milliarden Franken. Die wichtigsten Güter sind chemisch-pharmazeutische Produkte, Uhren, Maschinen und Präzisionsinstrumente. Der wichtigste Handelspartner bleibt Deutschland, gefolgt von den USA und China. 

Interessant ist die Entwicklung bei den Uhren: Trotz globaler Krisen stiegen die Exporte um 7,6 \%. Die Nachfrage nach Luxusgütern bleibt stabil. Im Pharmasektor investieren Firmen wie Novartis oder Roche \_\_\_\_\_\_\_\_\_\_ (viel) in Forschung. Dies macht die Schweiz \_\_\_\_\_\_\_\_\_\_ (innovativ) als viele Konkurrenten. 

Ein Problem stellt der starke Franken dar. Er macht die Produkte \_\_\_\_\_\_\_\_\_\_ (teuer) für ausländische Käufer. Wenn der Franken \_\_\_\_\_\_\_\_\_\_ (nicht / so / stark / sein / Konjunktiv II), \_\_\_\_\_\_\_\_\_\_ (die Exporte / steigen / Konjunktiv II) wahrscheinlich weiter.
\end{textebox}

\begin{enumerate}
    \item \textbf{Compréhension :} Nommez les trois principaux biens d'exportation suisses cités dans le texte.
    \item \textbf{Compréhension :} Quel pays est le premier partenaire commercial de la Suisse ?
    \item \textbf{Grammaire (Comparatif/Superlatif) :} Complétez les quatre vides dans le deuxième paragraphe avec la forme correcte de l'adjectif/adverbe entre parenthèses.
    \item \textbf{Grammaire (Conditionnel) :} Complétez les deux vides de la dernière phrase (\emph{Si-Satz} + principale) avec la forme correcte du \emph{Konjunktiv II}.
    \item \textbf{Vocabulaire :} Expliquez en allemand le sens de \all{par excellence} dans le contexte de la première phrase.
    \item \textbf{Expression :} Selon le texte, pourquoi le franc fort est-il un problème ? Répondez en une phrase complète.
\end{enumerate}
\end{exercice}

\begin{exercice}[Thème et Version : Épreuves types Bac]

\textbf{Partie A : Thème (Français $\rightarrow$ Allemand)}
Traduisez les 5 phrases suivantes en allemand. Soignez l'ordre des mots, les cas, les majuscules et les temps verbaux.
\begin{enumerate}
    \item Si les prix étaient plus bas, nous achèterions plus de produits.
    \item La lettre a été écrite hier par le directeur.
    \item Le marché allemand est plus grand que le marché autrichien.
    \item Pourriez-vous nous envoyer votre catalogue, s'il vous plaît ?
    \item Les marchandises doivent être livrées avant vendredi.
\end{enumerate}

\textbf{Partie B : Version (Allemand $\rightarrow$ Français)}
Traduisez le texte suivant en français courant et correct.
\begin{textebox}[Deutschland und Kamerun: Wirtschaftliche Beziehungen]
Die wirtschaftlichen Beziehungen zwischen Deutschland und Kamerun haben eine lange Tradition. Deutschland importiert vor allem Kakao, Holz und Aluminium aus Kamerun. Im Gegenzug liefert Deutschland Maschinen, Fahrzeuge und chemische Erzeugnisse. 

Ein neuer Schwerpunkt ist die Förderung von erneuerbaren Energien. Deutsche Firmen bauen Solaranlagen in ländlichen Regionen. Das schafft Arbeitsplätze und verbessert die Stromversorgung. Die kamerunische Regierung hofft, dass die Investitionen in den nächsten Jahren \_\_\_\_\_\_\_\_\_\_ (steigen / Futur I). Wenn die Infrastruktur \_\_\_\_\_\_\_\_\_\_ (besser / sein / Konjunktiv II), \_\_\_\_\_\_\_\_\_\_ (mehr / Firmen / kommen / Konjunktiv II).
\end{textebox}
\end{exercice}

\begin{exercice}[Expression écrite : Lettre de réclamation et Compte rendu]
Traitez \textbf{les deux} sujets. Respectez le nombre de lignes indiqué (environ 60--70 mots par sujet) et les conventions de rédaction (structure, formules, grammaire).

\textbf{Sujet 1 : Lettre de réclamation (\all{Beschwerdebrief})}
Votre entreprise (vous êtes le gérant de \all{„Nkam Business SARL“}, B.P. 123, Douala) a commandé 50 ordinateurs portables chez \all{„TechImport GmbH“}, Musterstraße 10, 10115 Berlin. La livraison est arrivée avec 10 jours de retard et 5 appareils sont endommagés. Écrivez une lettre formelle de réclamation (\all{Beschwerdebrief}) en demandant le remplacement des appareils défectueux et une réduction de prix pour le retard. Structure : \all{Briefkopf, Datum, Anrede, Betreff, Fließtext, Grußformel, Unterschrift}.

\textbf{Sujet 2 : Compte rendu de visite d'entreprise (\all{Bericht über einen Betriebsbesuch})}
Votre classe de Terminale A a visité l'usine \all{„CacaoTransform SA“} à Douala. Rédigez le compte rendu (\all{Bericht}) de cette visite pour le journal du lycée. Structure obligatoire : \all{Einleitung} (Quand ? Qui ? Où ? But), \all{Hauptteil} (Déroulement : accueil, visite des chaînes de production, transformation du cacao, contrôle qualité, discussion avec le directeur), \all{Schluss} (Bilan personnel, remerciements, signature/date). Utilisez le passé (\emph{Präteritum} ou \emph{Perfekt}) et le style indirect / passif là où c'est approprié.
\end{exercice}

\newpage

\coursec{Corrigés des exercices}

\begin{corrige}[Exercice 1 — QCM : grammaire de la vie économique]
\begin{enumerate}
    \item \textbf{Réponse A} : \all{größer als}. Le \cle{comparatif de supériorité} se forme avec le suffixe \all{-er} + \all{als}. L'adjectif \all{groß} prend l'\cle{Umlaut} : \all{groß $\rightarrow$ größer}. B est faux (pas de \all{mehr} + adjectif en allemand), C est un \cle{superlatif}, D est un calque du français « plus grand \emph{que} » : on dit \all{als}, jamais \all{wie} après un comparatif.
    \item \textbf{Réponse A} : \all{wir könnten}. Le \cle{Konjunktiv II} de \all{können} se construit sur le \cle{Präteritum} \all{konnt-} avec Umlaut : \all{ich könnte, wir könnten}. B est l'\cle{Indikatif présent} ; C (\all{würden können}) est un tour lourd à éviter avec les verbes modaux ; D est un \cle{Plus-que-parfait} (Konjunktiv II Perfekt).
    \item \textbf{Réponse A} : \all{Der Vertrag wird vom Direktor unterschrieben.} \cle{Passif} (Vorgangspassiv) au présent : \all{werden} (conjugué) + \cle{Participe II} en fin de phrase ; l'agent est introduit par \all{von} + \cle{Datif} (\all{vom} = \all{von dem}). B est au \cle{Präteritum} (\all{wurde}), C est un \cle{Zustandspassiv} (état : \all{ist unterschrieben}), D a un ordre des mots fautif (le participe II doit être en position finale).
    \item \textbf{Réponse B} : \all{Mit freundlichen Grüßen}. C'est la \cle{formule de clôture} standard d'une \cle{lettre commerciale} formelle. A (\all{Liebe Grüße}), C (\all{Bis bald}) et D (\all{Dein}) sont réservées aux lettres amicales.
\end{enumerate}
\end{corrige}

\begin{corrige}[Exercice 2 — Vrai ou Faux ? Justifiez en allemand]
\begin{enumerate}
    \item \textbf{Richtig.} \all{„Viel“ hat den unregelmäßigen Superlativ „am meisten“ (viel -- mehr -- am meisten).}
    \item \textbf{Falsch.} \all{„Ich war“ ist das Präteritum; der Konjunktiv II von „sein“ lautet „ich wäre“ (mit Umlaut).}
    \item \textbf{Falsch.} \all{Im Passiv wird der Täter (Agens) normalerweise mit „von“ + Dativ eingeleitet; „durch“ drückt das Mittel oder die Vermittlung aus.}
    \item \textbf{Falsch.} \all{Die Betreffzeile steht vor der Anrede, direkt unter der Adresse des Empfängers.}
    \item \textbf{Falsch.} \all{„Die Lieferung“ ist feminin; Nomen auf „-ung“ sind immer feminin.}
\end{enumerate}
\end{corrige}

\begin{corrige}[Exercice 3 — Vocabulaire thématique : la vie économique]
Paires correctes :
\begin{center}
\begin{tabular}{|l|l|}
\hline
\rowcolor{popblueL} \textbf{Deutsch} & \textbf{Réponse} \\
\hline
\all{die Lieferung} & B (action de livrer les marchandises) \\
\hline
\all{die Rechnung} & E (facture) \\
\hline
\all{das Angebot} & D (proposition commerciale) \\
\hline
\all{die Beschwerde} & C (réclamation) \\
\hline
\all{der Zahlungsaufschub} & F (délai de paiement) \\
\hline
\all{der Vertrag} & G (contrat) \\
\hline
\all{die Mahnung} & H (rappel de paiement) \\
\hline
\all{der Liefertermin} & I (date de livraison) \\
\hline
\end{tabular}
\end{center}
\textbf{Point de vigilance :} l'énoncé donnait \all{der Angebot} ; le genre correct est \all{das Angebot} (nom neutre, pluriel \all{die Angebote}). Retenez les genres et pluriels :
\begin{center}
\begin{tabular}{|l|l|}
\hline
\rowcolor{popblueL} \textbf{Nom (singulier)} & \textbf{Pluriel} \\
\hline
\all{die Lieferung} & \all{die Lieferungen} \\
\hline
\all{die Rechnung} & \all{die Rechnungen} \\
\hline
\all{das Angebot} & \all{die Angebote} \\
\hline
\all{die Beschwerde} & \all{die Beschwerden} \\
\hline
\all{der Zahlungsaufschub} & \all{die Zahlungsaufschübe} (Umlaut !) \\
\hline
\all{der Vertrag} & \all{die Verträge} (Umlaut !) \\
\hline
\all{die Mahnung} & \all{die Mahnungen} \\
\hline
\all{der Liefertermin} & \all{die Liefertermine} \\
\hline
\end{tabular}
\end{center}
\end{corrige}

\begin{corrige}[Exercice 4 — Conjugaison et formes grammaticales en contexte]
Texte complété :
\begin{quote}
\all{Sehr geehrte Damen und Herren,}

\all{bezugnehmend auf Ihr Angebot vom 10. Mai \textbf{freuten wir uns}, Ihnen mitteilen zu können, dass wir den Auftrag erteilen möchten. Die Preise \textbf{sind niedriger} als die der Konkurrenz, und die Qualität \textbf{ist am besten}.}

\all{Die Ware \textbf{muss geliefert werden} bis zum 30. Juni. Sollte der Termin nicht eingehalten \textbf{werden}, \textbf{würden wir vom Vertrag zurücktreten}.}

\all{Mit freundlichen Grüßen}
\end{quote}
Justifications :
\begin{itemize}
    \item \all{freuten wir uns} : \cle{Konjunktiv II} présent de \all{sich freuen}. \textbf{Attention :} \all{freuen} est un verbe \emph{faible} (régulier) ; son Präteritum est \all{freute}, et le Konjunktiv II est identique (\all{freute}) \textbf{sans Umlaut} (l'Umlaut ne concerne que les verbes forts et les modaux). C'est une formule de politesse atténuée.
    \item \all{sind niedriger} : \cle{comparatif} de \all{niedrig} + \all{als}.
    \item \all{ist am besten} : \cle{superlatif} irrégulier de \all{gut} (\all{gut -- besser -- am besten}).
    \item \all{muss geliefert werden} : \cle{passif avec modal} au présent : modal conjugué (\all{muss}) + \cle{Participe II} (\all{geliefert}) + \all{werden} (infinitif) en fin de phrase.
    \item \all{eingehalten werden} : \cle{Participe II} du \cle{verbe à particule séparable} \all{einhalten} (\all{eingehalten}) + infinitif \all{werden} (passif au \cle{Futur I} : \all{wird eingehalten werden}, ici dans une \cle{subordonnée conditionnelle sans \all{wenn}} : le verbe conjugué \all{sollte} est en tête).
    \item \all{würden wir ... zurücktreten} : \cle{Konjunktiv II} avec \all{würden} + infinitif ; \cle{inversion du sujet} après la subordonnée conditionnelle placée en tête ; la \cle{particule séparable} \all{zurück} se détache en fin de phrase.
\end{itemize}
\end{corrige}

\begin{corrige}[Exercice 5 — Thème : phrases à traduire en allemand]
\begin{enumerate}
    \item \all{Wenn die Preise niedriger wären, würden wir mehr Produkte kaufen.} — \cle{Conditionnel irréel} : \cle{Konjunktiv II} dans les deux propositions (\all{wären}, \all{würden kaufen}) ; \cle{comparatif} \all{niedriger}. (Variante formelle rare : \all{käufen wir} — incorrect car \all{kaufen} est faible, pas d'Umlaut ; on évite \all{kauften wir} qui sonne comme un passé).
    \item \all{Der Brief wurde gestern vom Direktor geschrieben.} — \cle{Passif} au \cle{Präteritum} : \all{wurde} + \cle{Participe II} ; agent avec \all{von} + \cle{Datif} (\all{vom}).
    \item \all{Der deutsche Markt ist größer als der österreichische Markt.} — \cle{Comparatif} avec Umlaut : \all{größer als}.
    \item \all{Könnten Sie uns bitte Ihren Katalog schicken?} — \cle{Konjunktiv II} de \cle{politesse} (\all{könnten}) ; \all{Ihren Katalog} : \cle{accusatif} masculin.
    \item \all{Die Waren müssen vor Freitag geliefert werden.} — \cle{Passif avec modal} : \all{müssen} + \cle{Participe II} + \all{werden} en position finale.
    \item \all{Wir bitten Sie, die Rechnung Nummer 456 zu begleichen.} (ou : \all{... die Rechnung Nr. 456 zu bezahlen.}) — Tournure commerciale : \all{Wir bitten Sie, ... zu + Infinitiv}.
    \item \all{Dieses Angebot ist das beste, das wir erhalten haben.} — \cle{Superlatif attributif} : \all{das beste} (déclinaison faible après \all{das}) ; \cle{subordonnée relative} au \cle{Perfekt} (le français « ayons reçu » se rend naturellement par le Perfekt).
    \item \all{Die Zahlung ist letzte Woche erfolgt.} (ou : \all{... ist letzte Woche geleistet worden.}) — \all{erfolgen} (intransitif) : \cle{Perfekt actif} avec \all{sein} (\all{ist erfolgt}) = « a eu lieu » ; \all{leisten} (transitif) : \cle{Passif Perfekt} (\all{ist geleistet worden}) = « a été effectué ».
\end{enumerate}
\textbf{Points de vigilance :} majuscule à tous les noms (\all{Preise, Produkte, Brief, Direktor, Markt, Katalog, Waren, Rechnung, Angebot, Zahlung}) ; \all{als} et non \all{wie} après un \cle{comparatif} ; place de \all{werden} en fin de phrase dans le \cle{passif avec modal}.
\end{corrige}

\begin{corrige}[Exercice 6 — Transformation : voix active $\rightarrow$ voix passive]
\begin{enumerate}
    \item \all{Die Maschinen werden von der Firma Müller geliefert.} — Présent : \all{werden} + \cle{Participe II} ; \all{von} + \cle{Datif} féminin (\all{von der}).
    \item \all{Die Rechnung wurde vom Kunden nicht bezahlt.} — \cle{Präteritum} : \all{wurde} ; la négation \all{nicht} se place avant le participe.
    \item \all{Die Pakete sind vom Spediteur abgeholt worden.} — \cle{Perfekt passif} : \all{sind} (auxiliaire \all{sein} pour le passif parfait) + \cle{Participe II} + \all{worden} (et non \all{geworden}) ; \all{Spediteur} masculin $\rightarrow$ \all{vom}.
    \item \all{Der Vertrag wird nächste Woche von uns unterschrieben werden.} — \cle{Futur I passif} : \all{wird} + \cle{Participe II} + \all{werden} ; \all{man/wir} devient \all{von uns}.
    \item \all{Die Waren müssen sorgfältig geprüft werden.} — \all{man} impersonnel disparaît au passif ; modal conjugué (\all{müssen}) + \cle{Participe II} + \all{werden}.
    \item \all{Die Entscheidung konnte vom Chef nicht getroffen werden.} — \cle{Präteritum} du modal \all{können} : \all{konnte} + \cle{Participe II} + \all{werden}.
    \item \all{Der Bericht wird auf Deutsch geschrieben.} — Sans agent (\all{man} non repris) ; notez la majuscule à \all{Deutsch} (langue \cle{nominalisée}).
    \item \all{Das Fenster ist geöffnet worden.} — Agent indéfini \all{jemand} : on ne le reprend pas ; \cle{Perfekt passif} : \all{ist ... worden}.
\end{enumerate}
\textbf{Rappel de méthode :} l'objet \cle{accusatif} de l'actif devient le \cle{sujet nominatif} du passif ; \all{werden} se conjugue au temps de la phrase active ; le \cle{Participe II} va en fin de phrase (avant \all{werden}/\all{worden} s'il y en a).
\end{corrige}

\begin{corrige}[Exercice 7 — Texte à trous : comparatif et superlatif]
Texte complété :
\begin{quote}
\all{Österreich und Deutschland sind wichtige Handelspartner. Der deutsche Markt ist \textbf{größer} als der österreichische Markt, aber die österreichische Exportquote ist oft \textbf{höher}. Schweizer Uhren sind weltweit \textbf{am bekanntesten} für ihre Präzision.}

\all{Im vergangenen Jahr haben deutsche Unternehmen \textbf{mehr} in Österreich investiert als je zuvor. Die Arbeitslosenzahlen in Deutschland sind \textbf{niedriger} als in vielen anderen EU-Ländern. Österreichische Unternehmen arbeiten \textbf{lieber} mit deutschen Partnern zusammen, weil die Sprache \textbf{einfacher} ist.}

\all{Das \textbf{beste} Ergebnis erzielt man jedoch durch gemeinsame Forschung. Die Zusammenarbeit ist \textbf{wichtiger} geworden in den letzten Jahren.}
\end{quote}
Formes à retenir :
\begin{center}
\begin{tabular}{|l|l|l|}
\hline
\rowcolor{popblueL} \textbf{Positif} & \textbf{Comparatif} & \textbf{Superlatif (prédicatif : \all{am ... sten})} \\
\hline
\all{groß} & \all{größer} (Umlaut) & \all{am größten} \\
\hline
\all{hoch} & \all{höher} (perte du \all{-c-}) & \all{am höchsten} \\
\hline
\all{bekannt} & \all{bekannter} & \all{am bekanntesten} (\all{-esten} après \all{-t}) \\
\hline
\all{viel} & \all{mehr} (irrégulier) & \all{am meisten} \\
\hline
\all{niedrig} & \all{niedriger} & \all{am niedrigsten} \\
\hline
\all{gern} & \all{lieber} (irrégulier) & \all{am liebsten} \\
\hline
\all{einfach} & \all{einfacher} & \all{am einfachsten} (\all{-esten} après \all{-ch}) \\
\hline
\all{gut} & \all{besser} (irrégulier) & \all{am besten} \\
\hline
\all{wichtig} & \all{wichtiger} & \all{am wichtigsten} \\
\hline
\end{tabular}
\end{center}
\textbf{Points de vigilance :} \all{hoch} perd le \all{-c-} au comparatif (\all{höher}) ; \all{bekannt} et \all{einfach} prennent \all{-esten} au superlatif (euphonie après \all{-t}, \all{-ch}) ; \all{das beste Ergebnis} : \cle{superlatif attributif} avec adjectif décliné (\all{-e} après article défini \all{das}).
\end{corrige}

\begin{corrige}[Exercice 8 — Appariement : structure de la lettre commerciale et du compte rendu]
Paires correctes :
\begin{center}
\begin{tabular}{|l|l|}
\hline
\rowcolor{popblueL} \textbf{Élément} & \textbf{Réponse} \\
\hline
1. \all{Der Briefkopf (Absender)} & C (coordonnées de l'expéditeur) \\
\hline
2. \all{Das Datum} & E (lieu et date : \all{Yaoundé, den 15. Oktober 2025}) \\
\hline
3. \all{Die Anrede} & A (\all{Sehr geehrte Damen und Herren,}) \\
\hline
4. \all{Die Betreffzeile} & B (objet, en gras, sans le mot \all{Betreff}) \\
\hline
5. \all{Der Fließtext (Hauptteil)} & D (corps : motifs, arguments, demandes) \\
\hline
6. \all{Die Grußformel} & F (\all{Mit freundlichen Grüßen}) \\
\hline
7. \all{Die Unterschrift} & G (signature manuscrite) \\
\hline
8. \all{Die Anlagen} & H (pièces jointes) \\
\hline
9. \all{Einleitung (Bericht)} & J (sujet, date, lieu, participants, objectif) \\
\hline
10. \all{Hauptteil (Bericht)} & I (résultats, décisions, mesures) \\
\hline
11. \all{Schluss (Bericht)} & K (résumé, suites, signature, date) \\
\hline
\end{tabular}
\end{center}
\textbf{Point de vigilance :} contrairement au français, la \cle{Betreffzeile} allemande s'écrit aujourd'hui \emph{sans} le mot \all{Betreff} devant (norme DIN 5008) ; elle se place \emph{avant} l'\cle{Anrede}. Après \all{Mit freundlichen Grüßen}, il n'y a \textbf{pas de virgule}. Après l'\all{Anrede}, virgule obligatoire, puis le corps commence par une \emph{minuscule}.
\end{corrige}

\begin{corrige}[Exercice 9 — Compréhension et langue : Les exportations suisses]
\textbf{Barème indicatif : 10 points} (compréhension 3 pts ; grammaire 4 pts ; vocabulaire 1 pt ; expression 2 pts).

\begin{enumerate}
    \item \all{Die drei wichtigsten Exportgüter der Schweiz sind chemisch-pharmazeutische Produkte, Uhren sowie Maschinen und Präzisionsinstrumente.} (1,5 pt)
    \item \all{Der wichtigste Handelspartner der Schweiz ist Deutschland.} (1,5 pt)
    \item \all{... investieren Firmen wie Novartis oder Roche \textbf{viel} in Forschung.} (positif, pas de comparaison ici) ; \all{Dies macht die Schweiz \textbf{innovativer} als viele Konkurrenten.} (\cle{comparatif} + \all{als}) ; \all{Er macht die Produkte \textbf{teurer} für ausländische Käufer.} (\cle{comparatif} de \all{teuer} : \all{teurer}, sans \all{-e-} supplémentaire). (3 pts)
    \item \all{Wenn der Franken \textbf{nicht so stark wäre}, \textbf{würden die Exporte} wahrscheinlich weiter \textbf{steigen}.} — \cle{Konjunktiv II} de \all{sein} : \all{wäre} ; principale avec \all{würden} + infinitif et \cle{inversion du sujet} (\all{würden die Exporte}). (1 pt)
    \item \all{„Par excellence“ bedeutet hier „in höchstem Maße“ oder „wie kaum ein anderes Land“: Die Schweiz ist ein typisches, vorbildliches Exportland.} (1 pt)
    \item \all{Der starke Franken ist ein Problem, weil er die Schweizer Produkte für ausländische Käufer teurer macht.} (2 pts)
\end{enumerate}
\textbf{Points de vigilance :} \all{teuer $\rightarrow$ teurer} (et non \all{teuerer}) ; dans la \cle{phrase conditionnelle} avec \all{wenn}, le verbe conjugué de la subordonnée va en \cle{fin de phrase} (\all{wäre}), et la principale commence par le verbe (\all{würden ... steigen}).
\end{corrige}

\begin{corrige}[Exercice 10 — Thème et Version : épreuves types Bac]
\textbf{Barème indicatif : 10 points} (Thème : 5 x 1 pt = 5 pts ; Version : 5 pts, dont fidélité du sens 3 pts, qualité du français 2 pts).

\textbf{Partie A — Thème (corrigé)}
\begin{enumerate}
    \item \all{Wenn die Preise niedriger wären, würden wir mehr Produkte kaufen.}
    \item \all{Der Brief wurde gestern vom Direktor geschrieben.}
    \item \all{Der deutsche Markt ist größer als der österreichische Markt.}
    \item \all{Könnten Sie uns bitte Ihren Katalog schicken?}
    \item \all{Die Waren müssen vor Freitag geliefert werden.}
\end{enumerate}
Critères d'évaluation par phrase : choix du temps/mode (0,25), ordre des mots et place du verbe (0,25), cas et articles (0,25), orthographe et majuscules (0,25).

\textbf{Partie B — Version (proposition de traduction)}

D'abord, complétons les vides du texte allemand : \all{... dass die Investitionen in den nächsten Jahren \textbf{steigen werden}. Wenn die Infrastruktur \textbf{besser wäre}, \textbf{würden mehr Firmen kommen}.}

Traduction :
\begin{quote}
Les relations économiques entre l'Allemagne et le Cameroun ont une longue tradition. L'Allemagne importe du Cameroun surtout du cacao, du bois et de l'aluminium. En contrepartie, l'Allemagne livre des machines, des véhicules et des produits chimiques.

Un nouvel axe prioritaire est la promotion des énergies renouvelables. Des entreprises allemandes construisent des installations solaires dans les régions rurales. Cela crée des emplois et améliore l'approvisionnement en électricité. Le gouvernement camerounais espère que les investissements augmenteront dans les années à venir. Si l'infrastructure était meilleure, davantage d'entreprises viendraient.
\end{quote}
\textbf{Points de vigilance :} \all{im Gegenzug} = « en contrepartie » ; \all{die Stromversorgung} = « l'approvisionnement en électricité » (faux ami : pas « la version du courant ») ; \all{Arbeitsplätze schaffen} = « créer des emplois » ; rendre le \cle{Futur I} (\all{werden steigen}) et le \cle{Konjunktiv II} (\all{wären, würden kommen}) par le futur et le conditionnel français.
\end{corrige}

\begin{corrige}[Exercice 11 — Expression écrite : lettre de réclamation et compte rendu]
\textbf{Barème indicatif : 10 points} (5 pts par sujet : respect de la structure 1 pt ; contenu et complétude 1,5 pt ; correction grammaticale — ordre des mots, cas, temps — 1,5 pt ; vocabulaire économique et formules 1 pt).

\textbf{Sujet 1 — Proposition de rédaction modèle (\all{Beschwerdebrief})}
\begin{quote}
\all{Nkam Business SARL\\
B.P. 123, Douala\\
Kamerun}

\all{TechImport GmbH\\
Musterstraße 10\\
10115 Berlin}

\all{Douala, den 15. Oktober 2025}

\all{\textbf{Beschwerde über die Lieferung vom 10. Oktober 2025}}

\all{Sehr geehrte Damen und Herren,}

\all{am 20. September haben wir 50 Laptops bei Ihnen bestellt. Leider ist die Lieferung erst mit zehn Tagen Verspätung angekommen. Außerdem sind fünf Geräte beschädigt worden.}

\all{Wir bitten Sie, die defekten Geräte so schnell wie möglich zu ersetzen. Für die Verspätung erwarten wir außerdem einen Preisnachlass von 10 \%.}

\all{Wir hoffen auf eine schnelle Lösung und verbleiben}

\all{mit freundlichen Grüßen}

\all{Jean Nkam}\\\all{Geschäftsführer}
\end{quote}
(« Nous vous prions de remplacer les appareils défectueux... Pour le retard, nous attendons en outre une remise de 10 \%. »)

\textbf{Sujet 2 — Proposition de rédaction modèle (\all{Bericht})}
\begin{quote}
\all{\textbf{Bericht über den Besuch der Firma CacaoTransform SA}}

\all{Am 12. Oktober 2025 besuchte die Klasse Terminale A die Fabrik CacaoTransform SA in Douala. Ziel des Besuchs war es, die Verarbeitung des Kakaos kennenzulernen.}

\all{Wir wurden von der Direktorin herzlich empfangen. Zuerst wurden uns die Produktionsanlagen gezeigt. Danach sahen wir, wie der Kakao geröstet und gemahlen wird. Die Qualität der Produkte wurde streng kontrolliert. Zum Schluss konnten wir der Direktorin Fragen stellen.}

\all{Der Besuch war sehr lehrreich. Wir danken der Firma für den freundlichen Empfang.}

\all{Douala, den 13. Oktober 2025}\\\all{Aminatou Bello, Klassensprecherin}
\end{quote}
(« Nous avons été chaleureusement accueillis... D'abord on nous a montré les installations... La qualité était strictement contrôlée. »)

\textbf{Points de vigilance :}
\begin{itemize}
    \item Pas de virgule après \all{Mit freundlichen Grüßen} ; virgule obligatoire après l'\all{Anrede}, puis le corps commence par une \emph{minuscule} (\all{am 20. September...}).
    \item Date : \all{Douala, den 15. Oktober 2025} (accusatif avec \all{den}).
    \item Compte rendu au passé : privilégier le \cle{Präteritum} (\all{besuchte, wurden, sahen, war, konnten, danken}) et le \cle{passif} (\all{wurden ... gezeigt, wurde kontrolliert}).
    \item Majuscules : \all{Laptops, Geräte, Verspätung, Preisnachlass, Besuch, Fabrik, Qualität, Empfang}.
    \item Éviter les calques : « nous espérons une solution rapide » $\rightarrow$ \all{wir hoffen auf eine schnelle Lösung} (avec \all{auf} + accusatif) ; « nous vous prions de... » $\rightarrow$ \all{Wir bitten Sie, ... zu + Infinitiv}.
\end{itemize}
\end{corrige}

\newpage

\coursec{Fiche bilan}

\begin{popfigure}
\begin{tikzpicture}[mindmap, grow cyclic, every node/.style={concept, text=white, font=\small},
root concept/.append style={concept color=popblue, font=\bfseries},
level 1/.append style={level distance=3.6cm, sibling angle=51},
level 2/.append style={level distance=2.2cm, sibling angle=40, font=\scriptsize}]
\node[root concept, align=center]{Vie économique\\ (Wirtschaft)}
child[concept color=poporange]{ node[align=center]{Lexique\\ der Betrieb, die Ware}
  child{ node{die Firma, -n} }
  child{ node{der Handel} }
}
child[concept color=popgreen]{ node[align=center]{Comparatif\\ Superlatif}
  child{ node{größer als} }
  child{ node{am größten} }
}
child[concept color=poppurple]{ node{Konjunktiv II}
  child{ node{würde + Inf.} }
  child{ node{hätte, wäre} }
}
child[concept color=poppink]{ node{Passif}
  child{ node{werden + Part. II} }
  child{ node{von + Dat.} }
}
child[concept color=popgold]{ node[align=center]{Lettre\\ commerciale}
  child{ node{Betreff, Anrede} }
}
child[concept color=popblue!70]{ node[align=center]{Compte rendu\\ (der Bericht)}
  child{ node{faits + passé} }
}
child[concept color=popdark]{ node[align=center]{Traduction\\ thème / version}
  child{ node{ordre des mots} }
};
\end{tikzpicture}
\legende{Carte mentale de la leçon : production et traduction sur la vie économique}
\end{popfigure}

\begin{bilan}
\textbf{1. Lexique clé de la vie économique}\par
\begin{itemize}
\item \all{die Wirtschaft, -en} : l'économie ; \all{der Betrieb, -e} : l'entreprise ; \all{die Firma, Firmen} : la société
\item \all{die Ware, -n} : la marchandise ; \all{der Preis, -e} : le prix ; \all{der Markt, die Märkte} : le marché
\item \all{der Kunde, -n / die Kundin, -nen} : le/la client(e) ; \all{der Arbeitgeber, - / die Arbeitnehmerin, -nen} : employeur / employée
\item \all{der Export, -e / der Import, -e} : l'exportation / l'importation ; \all{die Rechnung, -en} : la facture
\item \all{bestellen} : commander ; \all{liefern} : livrer ; \all{bezahlen} : payer ; \all{produzieren} : produire
\item \all{die Arbeitslosigkeit} : le chômage ; \all{das Gehalt, die Gehälter} : le salaire ; \all{die Steuer, -n} : l'impôt, la taxe
\end{itemize}

\textbf{2. Grammaire à connaître par cœur}\par
\begin{center}
\begin{tabularx}{\linewidth}{|l|X|X|}
\hline
\rowcolor{popblueL} \textbf{Point} & \textbf{Règle} & \textbf{Exemple} \\
\hline
Comparatif & adjectif + \all{-er} + \all{als} & \all{Kakao ist teurer als Kaffee.} \\
\hline
Superlatif & \all{am} + adjectif + \all{-sten} & \all{Der Hafen ist am wichtigsten.} \\
\hline
Konjunktiv II & \all{würde} + infinitif ; \all{hätte}, \all{wäre}, \all{könnte} & \all{Ich würde mehr exportieren.} \\
\hline
Passif présent & \all{werden} (conjugué) + participe II & \all{Die Ware wird geliefert.} \\
\hline
Passif passé & \all{wurde(n)} + participe II ; Perfekt : \all{ist ... worden} & \all{Die Rechnung wurde bezahlt.} \\
\hline
Agent au passif & \all{von} + datif & \all{Die Ware wird von der Firma geliefert.} \\
\hline
\end{tabularx}
\end{center}

\textbf{3. La lettre commerciale (der Geschäftsbrief)}\par
\begin{itemize}
\item Expéditeur et date en haut, destinataire, objet (\all{Betreff}), formule d'appel : \all{Sehr geehrte Damen und Herren,}
\item Corps : demande, offre ou réclamation ; conclusion : \all{Mit freundlichen Grüßen} + signature.
\end{itemize}

\textbf{4. Le compte rendu (der Bericht)}\par
\begin{itemize}
\item Titre, introduction (qui ? quoi ? où ? quand ?), faits au Präteritum ou au Perfekt, conclusion objective.
\item Style neutre, phrases claires, verbe en position 2.
\end{itemize}

\textbf{5. Méthodes express}\par
\begin{itemize}
\item \textbf{Thème (français $\rightarrow$ allemand)} : repérer le temps, placer le verbe conjugué en position 2, vérifier cas, genre et majuscules des noms.
\item \textbf{Version (allemand $\rightarrow$ français)} : identifier sujet/verbe/compléments, traduire le passif par « être + participe » ou « on », restituer un français naturel.
\item \textbf{Relecture} : verbe conjugué en 2\up{e} position, participe/infinitif en fin de phrase, majuscule à tous les noms.
\end{itemize}
\end{bilan}

\begin{aretenir}[Checklist avant le Bac]
\begin{itemize}
\item[$\square$] Je connais le lexique économique essentiel (entreprise, marché, prix, export, facture).
\item[$\square$] Je forme le comparatif avec \all{-er als} et le superlatif avec \all{am ... -sten}.
\item[$\square$] Je connais les comparatifs irréguliers : \all{gut -- besser -- am besten}, \all{viel -- mehr -- am meisten}, \all{gern -- lieber -- am liebsten}.
\item[$\square$] Je construis le Konjunktiv II avec \all{würde} + infinitif et je connais \all{hätte}, \all{wäre}, \all{könnte}, \all{müsste}.
\item[$\square$] Je forme le passif : \all{werden} + participe II, et le Perfekt passif : \all{ist ... worden}.
\item[$\square$] Je traduis l'agent au passif par \all{von} + datif.
\item[$\square$] Je respecte la structure d'une lettre commerciale : \all{Betreff}, \all{Sehr geehrte Damen und Herren}, \all{Mit freundlichen Grüßen}.
\item[$\square$] Je rédige un compte rendu objectif au passé avec les faits essentiels.
\item[$\square$] En thème, je place le verbe conjugué en position 2 et l'infinitif/participe en fin de phrase.
\item[$\square$] Je mets une majuscule à tous les noms allemands et j'utilise les guillemets „ … “.
\item[$\square$] Je me relis : cas (Akk./Dat.), genre des noms, orthographe, ponctuation.
\item[$\square$] Je gère mon temps : lecture du sujet, brouillon, rédaction, relecture.
\end{itemize}
\end{aretenir}

\findecours

\end{document}
